\input{preamble}

% OK, start here.
%
\begin{document}

\title{Limites de schémas}


\maketitle

\phantomsection
\label{section-phantom}

\tableofcontents

\section{Introduction}
\label{section-introduction}

\noindent
Dans ce chapitre, nous rassemblons des résultats relatifs aux limites de
schémas. Nous étudions principalement les limites de systèmes projectifs
indexés par des ensembles ordonnés filtrants
(Catégories, Définition \ref{categories-definition-directed-set})
dont les morphismes de transition sont affines. Nous traitons de
l'approximation noethérienne absolue. Nous caractérisons les schémas
localement de présentation finie sur une base comme ceux dont le foncteur
des points associé commute aux limites. Comme application de
l'approximation noethérienne absolue, nous démontrons que l'image d'un
schéma affine par un morphisme entier est affine. Nous établissons en outre
quelques variantes très générales du lemme de Chow. Une référence de base
est \cite{EGA}.




\section{Limites projectives de schémas à morphismes de transition affines}
\label{section-limits}

\noindent
Dans cette section, nous construisons la limite.

\begin{lemma}
\label{lemma-directed-inverse-system-affine-schemes-has-limit}
Soit $I$ un ensemble ordonné filtrant. Soit $(S_i, f_{ii'})$ un système
projectif de schémas indexé par $I$. Si tous les schémas $S_i$ sont
affines, alors la limite $S = \lim_i S_i$ existe dans la catégorie des
schémas. En fait, $S$ est affine et $S = \Spec(\colim_i R_i)$, où
$R_i = \Gamma(S_i, \mathcal{O})$.
\end{lemma}

\begin{proof}
Il suffit de poser $S = \Spec(\colim_i R_i)$.
Il résulte de Schémas, Lemme \ref{schemes-lemma-morphism-into-affine}
que $S$ est même la limite dans la catégorie des espaces localement
annelés.
\end{proof}

\begin{lemma}
\label{lemma-directed-inverse-system-has-limit}
Soit $I$ un ensemble ordonné filtrant. Soit $(S_i, f_{ii'})$ un système
projectif de schémas indexé par $I$. Si tous les morphismes
$f_{ii'} : S_i \to S_{i'}$ sont affines, alors la limite $S = \lim_i S_i$
existe dans la catégorie des schémas. De plus,
\begin{enumerate}
\item chacun des morphismes $f_i : S \to S_i$ est affine,
\item pour un élément $0 \in I$ et tout sous-schéma ouvert
$U_0 \subset S_0$, on a
$$
f_0^{-1}(U_0) = \lim_{i \geq 0} f_{i0}^{-1}(U_0)
$$
dans la catégorie des schémas.
\end{enumerate}
\end{lemma}

\begin{proof}
Choisissons un élément $0 \in I$. Remarquons que $I$ est non vide puisqu'il
est filtrant. Pour tout $i \geq 0$, considérons le faisceau quasi-cohérent
d'$\mathcal{O}_{S_0}$-algèbres
$\mathcal{A}_i = f_{i0, *}\mathcal{O}_{S_i}$. Rappelons que
$S_i = \underline{\Spec}_{S_0}(\mathcal{A}_i)$ ; voir Morphismes,
Lemme \ref{morphisms-lemma-characterize-affine}. Posons
$\mathcal{A} = \colim_{i \geq 0} \mathcal{A}_i$. C'est un faisceau
quasi-cohérent d'$\mathcal{O}_{S_0}$-algèbres ; voir Schémas,
section \ref{schemes-section-quasi-coherent}. Posons
$S = \underline{\Spec}_{S_0}(\mathcal{A})$. Par Morphismes,
Lemme \ref{morphisms-lemma-affine-equivalence-algebras}, nous obtenons,
pour $i \geq 0$, des morphismes $f_i : S \to S_i$ compatibles avec les
morphismes de transition. Les morphismes $f_i$ sont affines, par exemple
par Morphismes, Lemme \ref{morphisms-lemma-affine-permanence}. D'après le
Lemme \ref{lemma-directed-inverse-system-affine-schemes-has-limit} ci-dessus,
pour tout ouvert affine $U_0 \subset S_0$, l'image réciproque
$U = f_0^{-1}(U_0) \subset S$ est la limite, dans la catégorie des schémas,
du système des ouverts $U_i = f_{i0}^{-1}(U_0)$, $i \geq 0$.

\medskip\noindent
Soit $T$ un schéma. Soit $g_i : T \to S_i$ un système compatible de
morphismes. Pour montrer que $S = \lim_i S_i$, il faut établir qu'il existe
un unique morphisme $g : T \to S$ tel que $g_i = f_i \circ g$ pour tout
$i \in I$. Pour tout $t \in T$, il existe un ouvert affine
$U_0 \subset S_0$ contenant $g_0(t)$. Soit
$V \subset g_0^{-1}(U_0)$ un voisinage ouvert affine de $t$.
D'après les remarques précédentes, nous obtenons un unique morphisme
$g_V : V \to U = f_0^{-1}(U_0)$ tel que $f_i \circ g_V = g_i|_{U_i}$
pour tout $i$. Les ouverts $V \subset T$ ainsi construits forment une base
de la topologie de $T$. Les morphismes $g_V$ se recollent en un morphisme
$g : T \to S$ grâce à la propriété d'unicité. On obtient ainsi le morphisme
cherché $g : T \to S$.

\medskip\noindent
La dernière assertion résulte clairement de la construction précédente de
la limite.
\end{proof}

\begin{lemma}
\label{lemma-scheme-over-limit}
Soit $I$ un ensemble ordonné filtrant.
Soit $(S_i, f_{ii'})$ un système projectif de schémas indexé par $I$.
Supposons tous les morphismes $f_{ii'} : S_i \to S_{i'}$ affines.
Soit $S = \lim_i S_i$. Soit $0 \in I$. Supposons que $T$ soit un schéma
sur $S_0$. Alors
$$
T \times_{S_0} S = \lim_{i \geq 0} T \times_{S_0} S_i
$$
\end{lemma}

\begin{proof}
Le membre de droite est un schéma d'après le
Lemme \ref{lemma-directed-inverse-system-has-limit}. L'égalité est formelle ;
voir Catégories, Lemme \ref{categories-lemma-colimits-commute}.
\end{proof}



\section{Produits infinis}
\label{section-inifinite-products}

\noindent
En général, les produits infinis de schémas n'existent pas. Par exemple,
il est montré dans Exemples, section \ref{examples-section-not-algebraic},
qu'un produit infini de copies de $\mathbf{P}^1$ n'est même pas un espace
algébrique.

\medskip\noindent
En revanche, les produits infinis de schémas affines existent et sont
affines. Par Schémas, Lemme \ref{schemes-lemma-morphism-into-affine}, cela
correspond au fait que la catégorie des anneaux admet des coproduits
infinis : si $I$ est un ensemble et si $R_i$ est un anneau pour tout $i$,
on peut considérer l'anneau
$$
R = \otimes R_i =
\colim_{\{i_1, \ldots, i_n\} \subset I}
R_{i_1} \otimes_\mathbf{Z} \ldots \otimes_\mathbf{Z} R_{i_n}
$$
Étant donné un autre anneau $A$, une application $R \to A$ équivaut à une
famille d'homomorphismes d'anneaux $R_i \to A$, pour tout $i \in I$, comme
il résulte de la propriété correspondante des produits tensoriels finis.

\begin{lemma}
\label{lemma-infinite-product}
\begin{slogan}
Les produits infinis de schémas affines existent et sont affines.
\end{slogan}
Soit $S$ un schéma. Soit $I$ un ensemble et, pour tout $i \in I$, soit
$f_i : T_i \to S$ un morphisme affine. Alors le produit $T = \prod T_i$
existe dans la catégorie des schémas sur $S$. En fait, on a
$$
T = \lim_{\{i_1, \ldots, i_n\} \subset I}
T_{i_1} \times_S \ldots \times_S T_{i_n}
$$
et les morphismes de projection
$T \to T_{i_1} \times_S \ldots \times_S T_{i_n}$ sont affines.
\end{lemma}

\begin{proof}
Omis. Indication : raisonner comme dans la discussion qui précède le lemme
et utiliser le Lemme \ref{lemma-directed-inverse-system-has-limit} pour
l'existence de la limite.
\end{proof}

\begin{lemma}
\label{lemma-infinite-product-surjective}
Soit $S$ un schéma. Soit $I$ un ensemble et, pour tout $i \in I$, soit
$f_i : T_i \to S$ un morphisme affine surjectif. Alors le produit
$T = \prod T_i$ dans la catégorie des schémas sur $S$
(Lemme \ref{lemma-infinite-product}) se projette surjectivement sur $S$.
\end{lemma}

\begin{proof}
Soit $s \in S$. Choisissons $t_i \in T_i$ d'image $s$.
Choisissons une extension de corps suffisamment grande $K/\kappa(s)$ telle
que $\kappa(s_i)$ se plonge dans $K$ pour tout $i$. Nous obtenons alors des
morphismes $\Spec(K) \to T_i$ d'image $s_i$ qui coïncident comme morphismes
vers $S$. Il en résulte un morphisme $\Spec(K) \to T$, ce qui prouve qu'il
existe un point de $T$ d'image $s$.
\end{proof}

\begin{lemma}
\label{lemma-infinite-product-integral}
Soit $S$ un schéma. Soit $I$ un ensemble et, pour tout $i \in I$, soit
$f_i : T_i \to S$ un morphisme entier. Alors le produit $T = \prod T_i$
dans la catégorie des schémas sur $S$ (Lemme \ref{lemma-infinite-product})
est entier sur $S$.
\end{lemma}

\begin{proof}
Omis. Indication : sur des morceaux affines, cela se ramène au fait
algébrique suivant : si $A \to B_i$ est entier pour tout $i$, alors
$A \to \otimes_A B_i$ est entier.
\end{proof}





\section{Descente des propriétés}
\label{section-descent}

\noindent
D'abord quelques lemmes élémentaires décrivant la topologie d'une limite.

\begin{lemma}
\label{lemma-inverse-limit-sets}
Soit $S = \lim S_i$ la limite d'un système projectif filtrant de schémas
dont les morphismes de transition sont affines
(Lemme \ref{lemma-directed-inverse-system-has-limit}). Alors
$S_{set} = \lim_i S_{i, set}$, où $S_{set}$ désigne l'ensemble sous-jacent
au schéma $S$.
\end{lemma}

\begin{proof}
Fixons $i \in I$ et prenons un ouvert affine $U_i \subset S_i$.
Notons $U_{i'} = f_{i'i}^{-1}(U_i)$ et $U = f_i^{-1}(U_i)$. Ici,
$f_{i'i} : S_{i'} \to S_i$ est le morphisme de transition et
$f_i : S \to S_i$ la projection. D'après le
Lemme \ref{lemma-directed-inverse-system-has-limit}, on a
$U = \lim_{i' \geq i} U_i$. Supposons que l'on puisse montrer
$U_{set} = \lim_{i' \geq i} U_{i', set}$. Le lemme en résulte alors par un
argument simple utilisant un recouvrement affine de $S_i$. Nous pouvons
donc supposer que tous les $S_i$ et $S$ sont affines. Cela nous ramène à la
question algébrique examinée au paragraphe suivant.

\medskip\noindent
Soit donné un système d'anneaux $(A_i, \varphi_{ii'})$ indexé par $I$.
Posons $A = \colim_i A_i$, avec les homomorphismes canoniques
$\varphi_i : A_i \to A$. Alors
$$
\Spec(A) = \lim_i \Spec(A_i)
$$
En effet, supposons donnés des idéaux premiers
$\mathfrak p_i \subset A_i$ tels que
$\mathfrak p_i = \varphi_{ii'}^{-1}(\mathfrak p_{i'})$ pour tout
$i' \geq i$. Il suffit alors de poser
$$
\mathfrak p =
\{x \in A
\mid
\exists i, x_i \in \mathfrak p_i \text{ avec }\varphi_i(x_i) = x\}
$$
Il est clair que c'est un idéal et qu'il vérifie
$\varphi_i^{-1}(\mathfrak p) = \mathfrak p_i$. On en déduit aisément que
c'est également un idéal premier.
\end{proof}

\begin{lemma}
\label{lemma-inverse-limit-top}
\begin{reference}
\cite[IV, Proposition 8.2.9]{EGA}
\end{reference}
Soit $S = \lim S_i$ la limite d'un système projectif filtrant de schémas
dont les morphismes de transition sont affines
(Lemme \ref{lemma-directed-inverse-system-has-limit}). Alors
$S_{top} = \lim_i S_{i, top}$, où $S_{top}$ désigne l'espace topologique
sous-jacent au schéma $S$.
\end{lemma}

\begin{proof}
Nous utiliserons le critère de Topologie,
Lemme \ref{topology-lemma-characterize-limit}. Nous avons vu que
$S_{set} = \lim_i S_{i, set}$ dans le
Lemme \ref{lemma-inverse-limit-sets}. Les applications $f_i : S \to S_i$ sont des
morphismes de schémas, donc sont continues. Ainsi, $f_i^{-1}(U_i)$ est
ouvert pour tout ouvert $U_i \subset S_i$. Enfin, soient $s \in S$ et
$s \in V \subset S$ un voisinage ouvert. Choisissons $0 \in I$ et un
voisinage ouvert affine $U_0 \subset S_0$ de l'image de $s$. Alors
$f_0^{-1}(U_0) = \lim_{i \geq 0} f_{i0}^{-1}(U_0)$ ; voir le
Lemme \ref{lemma-directed-inverse-system-has-limit}. Les schémas
$f_0^{-1}(U_0)$ et $f_{i0}^{-1}(U_0)$ sont affines et
$$
\mathcal{O}_S(f_0^{-1}(U_0)) =
\colim_{i \geq 0} \mathcal{O}_{S_i}(f_{i0}^{-1}(U_0))
$$
soit d'après la démonstration du
Lemme \ref{lemma-directed-inverse-system-has-limit}, soit d'après le
Lemme \ref{lemma-directed-inverse-system-affine-schemes-has-limit}.
Choisissons $a \in \mathcal{O}_S(f_0^{-1}(U_0))$ tel que
$s \in D(a) \subset V$. C'est possible puisque les ouverts principaux
forment une base de la topologie du schéma affine $f_0^{-1}(U_0)$.
On peut alors choisir $i \geq 0$ et
$a_i \in \mathcal{O}_{S_i}(f_{i0}^{-1}(U_0))$ d'image $a$. Il en résulte
que $D(a_i) \subset f_{i0}^{-1}(U_0) \subset S_i$ est un ouvert dont
l'image réciproque dans $S$ est $D(a)$. Cela achève la démonstration.
\end{proof}

\begin{lemma}
\label{lemma-limit-nonempty}
Soit $S = \lim S_i$ la limite d'un système projectif filtrant de schémas
dont les morphismes de transition sont affines
(Lemme \ref{lemma-directed-inverse-system-has-limit}). Si tous les schémas
$S_i$ sont non vides et quasi-compacts, alors la limite
$S = \lim_i S_i$ est non vide.
\end{lemma}

\begin{proof}
Choisissons $0 \in I$. Remarquons que $I$ est non vide puisqu'il est
filtrant. Choisissons un recouvrement ouvert affine
$S_0 = \bigcup_{j = 1, \ldots, m} U_j$. Comme $I$ est filtrant, il existe
$j \in \{1, \ldots, m\}$ tel que
$f_{i0}^{-1}(U_j) \not = \emptyset$ pour tout $i \geq 0$. Par suite,
$\lim_{i \geq 0} f_{i0}^{-1}(U_j)$ n'est pas vide, car une limite
inductive filtrante d'anneaux non nuls est non nulle (puisque
$1 \not = 0$). Comme $\lim_{i \geq 0} f_{i0}^{-1}(U_j)$ est un sous-schéma
ouvert de la limite, on conclut.
\end{proof}

\begin{lemma}
\label{lemma-inverse-limit-irreducibles}
Soit $S = \lim S_i$ la limite d'un système projectif filtrant de schémas
dont les morphismes de transition sont affines
(Lemme \ref{lemma-directed-inverse-system-has-limit}). Soit $s \in S$
d'images $s_i \in S_i$. Alors
\begin{enumerate}
\item $s = \lim s_i$ comme schémas, c'est-à-dire
$\kappa(s) = \colim \kappa(s_i)$,
\item $\overline{\{s\}} = \lim \overline{\{s_i\}}$ comme ensembles, et
\item $\overline{\{s\}} = \lim \overline{\{s_i\}}$ comme schémas, où
$\overline{\{s\}}$ et $\overline{\{s_i\}}$ sont munis de leur structure
induite de schéma réduit.
\end{enumerate}
\end{lemma}

\begin{proof}
Choisissons $0 \in I$ et un recouvrement ouvert affine
$S_0 = \bigcup_{j \in J} U_{0, j}$. Pour $i \geq 0$, posons
$U_{i, j} = f_{i, 0}^{-1}(U_{0, j})$ et
$U_j = f_0^{-1}(U_{0, j})$. Ici, $f_{i'i} : S_{i'} \to S_i$ est le
morphisme de transition et $f_i : S \to S_i$ la projection. Pour
$j \in J$, les conditions suivantes sont équivalentes :
(a) $s \in U_j$, (b) $s_0 \in U_{0, j}$,
(c) $s_i \in U_{i, j}$ pour tout $i \geq 0$. Soit $J' \subset J$
l'ensemble des indices pour lesquels (a), (b) et (c) sont vérifiées.
Alors $\overline{\{s\}} = \bigcup_{j \in J'} (\overline{\{s\}} \cap U_j)$,
et de même pour $\overline{\{s_i\}}$ lorsque $i \geq 0$. Remarquons que
$\overline{\{s\}} \cap U_j$ est l'adhérence de $\{s\}$ dans l'espace
topologique $U_j$. Il en va de même de
$\overline{\{s_i\}} \cap U_{i, j}$ pour $i \geq 0$. Il suffit donc de
démontrer le lemme lorsque $S$ et $S_i$ sont affines pour tout $i$. Cela nous
ramène à la question algébrique du paragraphe suivant.

\medskip\noindent
Soit donné un système d'anneaux $(A_i, \varphi_{ii'})$ indexé par $I$.
Posons $A = \colim_i A_i$, avec les homomorphismes canoniques
$\varphi_i : A_i \to A$. Soit $\mathfrak p \subset A$ un idéal premier et
posons $\mathfrak p_i = \varphi_i^{-1}(\mathfrak p)$. Alors
$$
V(\mathfrak p) = \lim_i V(\mathfrak p_i)
$$
Cela résulte du Lemme \ref{lemma-inverse-limit-sets}, puisque
$A/\mathfrak p = \colim A_i/\mathfrak p_i$. Cette égalité d'anneaux montre
aussi la dernière assertion concernant les structures induites de schéma
réduit. L'égalité $\kappa(\mathfrak p) = \colim \kappa(\mathfrak p_i)$
résulte également de cette assertion.
\end{proof}

\noindent
Dans la suite de cette section, nous nous plaçons dans la situation suivante.

\begin{situation}
\label{situation-descent}
Soit $S = \lim_{i \in I} S_i$ la limite d'un système projectif filtrant de schémas
dont les morphismes de transition $f_{i'i} : S_{i'} \to S_i$ sont affines
(Lemme \ref{lemma-directed-inverse-system-has-limit}). Nous supposons que
$S_i$ est quasi-compact et quasi-séparé pour tout $i \in I$. Nous notons
$f_i : S \to S_i$ la projection. Nous choisissons aussi un élément
$0 \in I$.
\end{situation}

\noindent
Dans cette situation, le morphisme $S \to S_0$ est affine. Il s'ensuit que
$S$ est quasi-compact et quasi-séparé\footnote{Cela résulte de Morphismes,
Lemme \ref{morphisms-lemma-affine-separated}, Topologie,
Définition \ref{topology-definition-quasi-compact}, et Schémas,
Lemme \ref{schemes-lemma-separated-permanence}.}. Le type de résultat que
nous recherchons est le suivant : si l'on dispose d'un objet sur $S$, alors,
pour un certain $i$, il existe un objet analogue sur $S_i$.

\begin{lemma}
\label{lemma-topology-limit}
Dans la Situation \ref{situation-descent}.
\begin{enumerate}
\item On a $S_{set} = \lim_i S_{i, set}$, où $S_{set}$ désigne l'ensemble
sous-jacent au schéma $S$.
\item On a $S_{top} = \lim_i S_{i, top}$, où $S_{top}$ désigne l'espace
topologique sous-jacent au schéma $S$.
\item Si $s, s' \in S$ et si $s'$ n'est pas une spécialisation de $s$,
alors, pour un certain $i \in I$, l'image $s'_i \in S_i$ de $s'$ n'est pas
une spécialisation de l'image $s_i \in S_i$ de $s$.
\item Ajouter ici d'autres faits élémentaires sur la topologie de $S$.
(Exigence : tout ajout doit être facile dans le cas affine.)
\end{enumerate}
\end{lemma}

\begin{proof}
L'assertion (1) est un cas particulier du Lemme \ref{lemma-inverse-limit-sets}.

\medskip\noindent
L'assertion (2) est un cas particulier du Lemme \ref{lemma-inverse-limit-top}.

\medskip\noindent
L'assertion (3) est un cas particulier du
Lemme \ref{lemma-inverse-limit-irreducibles}.
\end{proof}

\begin{lemma}
\label{lemma-descend-section}
Dans la Situation \ref{situation-descent}. Supposons que $\mathcal{F}_0$
soit un faisceau quasi-cohérent sur $S_0$. Posons
$\mathcal{F}_i = f_{i0}^*\mathcal{F}_0$ pour $i \geq 0$ et
$\mathcal{F} = f_0^*\mathcal{F}_0$. Alors
$$
\Gamma(S, \mathcal{F}) = \colim_{i \geq 0} \Gamma(S_i, \mathcal{F}_i)
$$
\end{lemma}

\begin{proof}
Écrivons $\mathcal{A}_j = f_{i0, *} \mathcal{O}_{S_i}$. C'est un faisceau
quasi-cohérent d'$\mathcal{O}_{S_0}$-algèbres (voir Morphismes,
Lemme \ref{morphisms-lemma-affine-equivalence-algebras}) et $S_i$ est le
spectre relatif de $\mathcal{A}_i$ sur $S_0$. Dans la démonstration du
Lemme \ref{lemma-directed-inverse-system-has-limit}, nous avons construit
$S$ comme le spectre relatif de
$\mathcal{A} = \colim_{i \geq 0} \mathcal{A}_i$ sur $S_0$. Posons
$$
\mathcal{M}_i = \mathcal{F}_0 \otimes_{\mathcal{O}_{S_0}} \mathcal{A}_i
$$
et
$$
\mathcal{M} = \mathcal{F}_0 \otimes_{\mathcal{O}_{S_0}} \mathcal{A}.
$$
On a alors $f_{i0, *} \mathcal{F}_i = \mathcal{M}_i$ et
$f_{0, *}\mathcal{F} = \mathcal{M}$. Comme $\mathcal{A}$ est la limite
inductive des faisceaux $\mathcal{A}_i$ et que le produit tensoriel commute
aux limites inductives filtrantes, on en déduit
$\mathcal{M} = \colim_{i \geq 0} \mathcal{M}_i$. Puisque $S_0$ est
quasi-compact et quasi-séparé, on obtient
\begin{eqnarray*}
\Gamma(S, \mathcal{F})
& = &
\Gamma(S_0, \mathcal{M}) \\
& = &
\Gamma(S_0, \colim_{i \geq 0} \mathcal{M}_i) \\
& = &
\colim_{i \geq 0} \Gamma(S_0, \mathcal{M}_i) \\
& = &
\colim_{i \geq 0} \Gamma(S_i, \mathcal{F}_i)
\end{eqnarray*}
voir Faisceaux, Lemme \ref{sheaves-lemma-directed-colimits-sections}, et
Topologie, Lemme \ref{topology-lemma-topology-quasi-separated-scheme}, pour
l'égalité intermédiaire.
\end{proof}

\begin{lemma}
\label{lemma-limit-closed-nonempty}
Dans la Situation \ref{situation-descent}. Supposons donné, pour chaque $i$,
un fermé non vide $Z_i \subset S_i$ tel que
$f_{i'i}(Z_{i'}) \subset Z_i$ pour tout $i' \geq i$. Alors il existe un
point $s \in S$ tel que $f_i(s) \in Z_i$ pour tout $i$.
\end{lemma}

\begin{proof}
Notons encore $Z_i \subset S_i$ le sous-schéma fermé réduit associé à
$Z_i$ ; voir Schémas,
Définition \ref{schemes-definition-reduced-induced-scheme}. Une immersion
fermée est affine et une composée de morphismes affines est affine (voir
Morphismes, Lemmes \ref{morphisms-lemma-closed-immersion-affine} et
\ref{morphisms-lemma-composition-affine}) ; par conséquent,
$Z_{i'} \to S_i$ est affine lorsque $i' \geq i$. On en déduit que le morphisme
$f_{i'i} : Z_{i'} \to Z_i$ est affine par
Morphismes, Lemme \ref{morphisms-lemma-affine-permanence}. Chacun des schémas $Z_i$ est
quasi-compact, car c'est un sous-schéma fermé d'un schéma quasi-compact.
Nous pouvons donc appliquer le Lemme \ref{lemma-limit-nonempty} pour voir
que $Z = \lim_i Z_i$ est non vide. Comme il existe un morphisme canonique
$Z \to S$, on conclut.
\end{proof}

\begin{lemma}
\label{lemma-limit-fibre-product-empty}
Dans la Situation \ref{situation-descent}. Supposons donnés un $i$ et un
morphisme $T \to S_i$ tels que
\begin{enumerate}
\item $T \times_{S_i} S = \emptyset$, et
\item $T$ est quasi-compact.
\end{enumerate}
Alors $T \times_{S_i} S_{i'} = \emptyset$ pour tout $i'$ assez grand.
\end{lemma}

\begin{proof}
D'après le Lemme \ref{lemma-scheme-over-limit}, on a
$T \times_{S_i} S = \lim_{i' \geq i} T \times_{S_i} S_{i'}$. Le résultat
découle donc du Lemme \ref{lemma-limit-nonempty}.
\end{proof}

\begin{lemma}
\label{lemma-limit-contained-in-constructible}
Dans la Situation \ref{situation-descent}. Supposons donnés un $i$ et une
partie localement constructible $E \subset S_i$ telle que
$f_i(S) \subset E$. Alors $f_{i'i}(S_{i'}) \subset E$ pour tout $i'$ assez
grand.
\end{lemma}

\begin{proof}
En écrivant $S_i$ comme réunion finie de sous-schémas ouverts affines, on
se ramène au cas où $S_i$ est affine et $E$ constructible ; voir le
Lemme \ref{lemma-directed-inverse-system-has-limit} et Propriétés,
Lemme \ref{properties-lemma-locally-constructible}. Dans ce cas, le
complémentaire $S_i \setminus E$ est lui aussi constructible. Il existe donc
un schéma affine $T$ et un morphisme $T \to S_i$ dont l'image est
$S_i \setminus E$ ; voir Algèbre,
Lemme \ref{algebra-lemma-constructible-is-image}. D'après le
Lemme \ref{lemma-limit-fibre-product-empty}, le schéma
$T \times_{S_i} S_{i'}$ est vide pour tout $i'$ assez grand, et par
conséquent $f_{i'i}(S_{i'}) \subset E$ pour tout $i'$ assez grand.
\end{proof}

\begin{lemma}
\label{lemma-descend-opens}
Dans la Situation \ref{situation-descent}, on a les assertions suivantes :
\begin{enumerate}
\item Pour tout ouvert quasi-compact $V \subset S = \lim_i S_i$, il existe
$i \in I$ et un ouvert quasi-compact $V_i \subset S_i$ tels que
$f_i^{-1}(V_i) = V$.
\item Étant donnés des ouverts quasi-compacts $V_i \subset S_i$ et
$V_{i'} \subset S_{i'}$ tels que
$f_i^{-1}(V_i) = f_{i'}^{-1}(V_{i'})$, il existe un indice
$i'' \geq i, i'$ tel que
$f_{i''i}^{-1}(V_i) = f_{i''i'}^{-1}(V_{i'})$.
\item Si $V_{1, i}, \ldots, V_{n, i} \subset S_i$ sont des ouverts
quasi-compacts et si
$S = f_i^{-1}(V_{1, i}) \cup \ldots \cup f_i^{-1}(V_{n, i})$, alors
$S_{i'} = f_{i'i}^{-1}(V_{1, i}) \cup \ldots \cup f_{i'i}^{-1}(V_{n, i})$
pour un certain $i' \geq i$.
\end{enumerate}
\end{lemma}

\begin{proof}
Choisissons $i_0 \in I$. Remarquons que $I$ est non vide puisqu'il est
filtrant. Pour simplifier les notations, écrivons $S_0 = S_{i_0}$ et
$i_0 = 0$. Choisissons un recouvrement ouvert affine
$S_0 = U_{1, 0} \cup \ldots \cup U_{m, 0}$. Notons
$U_{j, i} \subset S_i$ l'image réciproque de $U_{j, 0}$ par le morphisme de
transition pour $i \geq 0$. Notons $U_j$ l'image réciproque de $U_{j, 0}$ dans $S$.
Remarquons que $U_j = \lim_i U_{j, i}$ est une limite de schémas affines.

\medskip\noindent
Nous démontrons d'abord l'assertion d'unicité. Soient des ouverts
quasi-compacts $V_i \subset S_i$ et $V_{i'} \subset S_{i'}$ tels que
$f_i^{-1}(V_i) = f_{i'}^{-1}(V_{i'})$. Il suffit de montrer que
$f_{i''i}^{-1}(V_i \cap U_{j, i''})$ et
$f_{i''i'}^{-1}(V_{i'} \cap U_{j, i''})$ deviennent égaux pour $i''$ assez
grand. On se ramène ainsi au cas d'une limite de schémas affines. Dans ce
cas, écrivons $S = \Spec(R)$ et $S_i = \Spec(R_i)$ pour tout $i \in I$.
On peut écrire $V_i = S_i \setminus V(h_1, \ldots, h_m)$ et
$V_{i'} = S_{i'} \setminus V(g_1, \ldots, g_n)$. L'hypothèse signifie que
les idéaux $\sum g_jR$ et $\sum h_jR$ ont le même radical dans $R$.
Cela signifie que $g_j^N = \sum a_{jj'}h_{j'}$ et
$h_j^N = \sum b_{jj'} g_{j'}$ pour un certain $N \gg 0$ et certains
$a_{jj'}$, $b_{jj'}$ dans $R$. Puisque $R = \colim_i R_i$, nous pouvons
choisir un indice $i'' \geq i$ tel que les égalités
$g_j^N = \sum a_{jj'}h_{j'}$ et
$h_j^N = \sum b_{jj'} g_{j'}$ soient vérifiées dans $R_{i''}$ pour certains
$a_{jj'}$ et $b_{jj'}$ dans $R_{i''}$. Il en résulte que les idéaux
$\sum g_jR_{i''}$ et $\sum h_jR_{i''}$ ont le même radical dans $R_{i''}$,
comme voulu.

\medskip\noindent
Montrons l'existence. Si $S_0$ est affine, alors
$S_i = \Spec(R_i)$ pour tout $i \geq 0$ et $S = \Spec(R)$, où
$R = \colim R_i$. On a alors
$V = S \setminus V(g_1, \ldots, g_n)$ pour certains
$g_1, \ldots, g_n \in R$. Choisissons $i$ assez grand pour que chacun des
$g_j$ provienne d'un élément $g_{j, i} \in R_i$, et prenons
$V_i = S_i \setminus V(g_{1, i}, \ldots, g_{n, i})$. Pour $S_0$ quelconque,
les ouverts $V \cap U_j$ sont quasi-compacts puisque $S$ est quasi-séparé.
Le cas affine montre donc que, pour chaque $j = 1, \ldots, m$, il existe
$i_j \in I$ et un ouvert quasi-compact $V_{i_j} \subset U_{j, i_j}$ dont
l'image réciproque dans $U_j$ est $V \cap U_j$. Posons
$i = \max(i_1, \ldots, i_m)$ et
$V_i = \bigcup f_{ii_j}^{-1}(V_{i_j})$.

\medskip\noindent
L'assertion sur les recouvrements découle de l'assertion d'unicité appliquée
aux ouverts $V_{1, i} \cup \ldots \cup V_{n, i}$ et $S_i$ de $S_i$.
\end{proof}

\begin{lemma}
\label{lemma-limit-quasi-affine}
Dans la Situation \ref{situation-descent}, si $S$ est quasi-affine, alors,
pour un certain $i_0 \in I$, les schémas $S_i$ sont quasi-affines pour tout
$i \geq i_0$.
\end{lemma}

\begin{proof}
Choisissons $i_0 \in I$. Remarquons que $I$ est non vide puisqu'il est
filtrant. Pour simplifier les notations, écrivons $S_0 = S_{i_0}$ et
$i_0 = 0$. Soit $s \in S$. Nous pouvons choisir un ouvert affine
$U_0 \subset S_0$ contenant $f_0(s)$. Puisque $S$ est quasi-affine, nous
pouvons choisir $a \in \Gamma(S, \mathcal{O}_S)$ tel que
$s \in D(a) \subset f_0^{-1}(U_0)$ et tel que $D(a)$ soit affine. D'après le
Lemme \ref{lemma-descend-section}, il existe $i \geq 0$ tel que $a$ provienne
d'un élément $a_i \in \Gamma(S_i, \mathcal{O}_{S_i})$. Pour tout indice
$j \geq i$, notons $a_j$ l'image de $a_i$ dans les sections globales du
faisceau structural de $S_j$. Considérons les ouverts $D(a_j) \subset S_j$
et $U_j = f_{j0}^{-1}(U_0)$. Remarquons que $U_j$ est affine et que
$D(a_j)$ est un ouvert quasi-compact de $S_j$ ; voir par exemple Propriétés,
Lemme \ref{properties-lemma-affine-cap-s-open}. Nous pouvons donc appliquer
le Lemme \ref{lemma-descend-opens} aux ouverts $U_j$ et
$U_j \cup D(a_j)$ pour conclure que $D(a_j) \subset U_j$ pour un certain
$j \geq i$. Pour un tel indice $j$, $D(a_j) \subset S_j$ est un ouvert
affine (car $D(a_j)$ est un ouvert affine standard de l'ouvert affine $U_j$)
contenant l'image $f_j(s)$.

\medskip\noindent
Nous en déduisons que, pour tout $s \in S$, il existe un indice $i \in I$
et une section globale $a \in \Gamma(S_i, \mathcal{O}_{S_i})$ tels que
$D(a) \subset S_i$ soit un ouvert affine contenant $f_i(s)$. Comme $S$ est
quasi-compact, nous pouvons choisir un seul indice $i \in I$ et des sections
globales $a_1, \ldots, a_m \in \Gamma(S_i, \mathcal{O}_{S_i})$ tels que
chaque $D(a_j) \subset S_i$ soit un ouvert affine et que l'image de
$f_i : S \to S_i$ soit contenue dans la réunion
$W_i = \bigcup_{j = 1, \ldots, m} D(a_j)$. Pour $i' \geq i$, posons
$W_{i'} = f_{i'i}^{-1}(W_i)$. Puisque $f_i^{-1}(W_i)$ est égal à $S$, une nouvelle
application du Lemme \ref{lemma-descend-opens} montre que
pour un $i' \geq i$ convenable, on a $S_{i'} = W_{i'}$. Nous pouvons donc remplacer
$i$ par $i'$ et supposer
$S_i = \bigcup_{j = 1, \ldots, m} D(a_j)$. Il en résulte que
$\mathcal{O}_{S_i}$ est un faisceau inversible ample sur $S_i$ (voir
Propriétés, Définition \ref{properties-definition-ample}) et donc que $S_i$
est quasi-affine ; voir Propriétés,
Lemme \ref{properties-lemma-quasi-affine-O-ample}. On conclut.
\end{proof}

\begin{lemma}
\label{lemma-limit-affine}
Dans la Situation \ref{situation-descent}, si $S$ est affine, alors, pour
un certain $i_0 \in I$, les schémas $S_i$ sont affines pour tout
$i \geq i_0$.
\end{lemma}

\begin{proof}
D'après le Lemme \ref{lemma-limit-quasi-affine}, nous pouvons supposer que
$S_0$ est quasi-affine pour un certain $0 \in I$. Posons
$R_0 = \Gamma(S_0, \mathcal{O}_{S_0})$. Alors $S_0$ est un ouvert
quasi-compact de $T_0 = \Spec(R_0)$. Notons $j_0 : S_0 \to T_0$
l'immersion ouverte quasi-compacte correspondante. Pour $i \geq 0$, posons
$\mathcal{A}_i = f_{i0, *}\mathcal{O}_{S_i}$. Comme $f_{i0}$ est affine,
on a $S_i = \underline{\Spec}_{S_0}(\mathcal{A}_i)$. Posons
$T_i = \underline{\Spec}_{T_0}(j_{0, *}\mathcal{A}_i)$. Alors
$T_i \to T_0$ est affine, donc $T_i$ est affine. Ainsi, $T_i$ est le spectre
de
$$
R_i = \Gamma(T_0, j_{0, *}\mathcal{A}_i) = \Gamma(S_0, \mathcal{A}_i) =
\Gamma(S_i, \mathcal{O}_{S_i}).
$$
Écrivons $S = \Spec(R)$. On a $R = \colim_i R_i$
d'après le Lemme \ref{lemma-descend-section}. On a donc aussi $S = \lim_i T_i$. Comme la formation du
spectre relatif commute au changement de base, l'image réciproque de
l'ouvert $S_0 \subset T_0$ dans $T_i$ est $S_i$. Soit
$Z_0 = T_0 \setminus S_0$, et soit $Z_i \subset T_i$ l'image réciproque de
$Z_0$. Puisque $S_i = T_i \setminus Z_i$, il suffit de montrer que $Z_i$ est
vide pour un certain $i$. Supposons, en vue d'une contradiction, que $Z_i$
soit non vide pour tout $i$. D'après le
Lemme \ref{lemma-limit-closed-nonempty}, il existe un point $s$ de
$S = \lim T_i$ dont l'image appartient à $Z_i$ pour tout $i$. Mais
$S = \lim_i S_i$, ce qui contredit le Lemme \ref{lemma-topology-limit}.
\end{proof}

\begin{lemma}
\label{lemma-limit-separated}
Dans la Situation \ref{situation-descent}, si $S$ est séparé, alors, pour un
certain $i_0 \in I$, les schémas $S_i$ sont séparés pour tout $i \geq i_0$.
\end{lemma}

\begin{proof}
Choisissons un recouvrement ouvert affine fini
$S_0 = U_{0, 1} \cup \ldots \cup U_{0, m}$. Notons
$U_{i, j} \subset S_i$ et $U_j \subset S$ les images réciproques de
$U_{0, j}$. Remarquons que $U_{i, j}$ et $U_j$ sont affines. Comme $S$ est
séparé, les intersections $U_{j_1} \cap U_{j_2}$ sont affines. Puisque
$U_{j_1} \cap U_{j_2} = \lim_{i \geq 0} U_{i, j_1} \cap U_{i, j_2}$,
on voit que $U_{i, j_1} \cap U_{i, j_2}$ est affine pour $i$ assez grand,
d'après le Lemme \ref{lemma-limit-affine}. Pour montrer
que $S_i$ est séparé lorsque $i$ est assez grand, il suffit maintenant de
montrer que
$$
\mathcal{O}_{S_i}(U_{i, j_1})
\otimes_{\mathcal{O}_S(S)}
\mathcal{O}_{S_i}(U_{i, j_2})
\longrightarrow
\mathcal{O}_{S_i}(U_{i, j_1} \cap U_{i, j_2})
$$
est surjectif pour $i$ assez grand (Schémas,
Lemme \ref{schemes-lemma-characterize-separated}).

\medskip\noindent
Pour alléger les indices, supposons donnés des ouverts affines
$U, V \subset S_0$ tels que $U \cap V$ soit également affine. Notons
$U_i, V_i \subset S_i$, resp.\ $U, V \subset S$, leurs images
réciproques. Il faut montrer que
$\mathcal{O}(U_i) \otimes \mathcal{O}(V_i) \to
\mathcal{O}(U_i \cap V_i)$
est surjectif pour $i$ assez grand, sachant que
$\mathcal{O}(U) \otimes \mathcal{O}(V) \to \mathcal{O}(U \cap V)$
est surjectif. Remarquons que
$\mathcal{O}(U_0) \otimes \mathcal{O}(V_0) \to
\mathcal{O}(U_0 \cap V_0)$
est de type fini, car le morphisme diagonal $S_i \to S_i \times S_i$ est une
immersion (Schémas, Lemme \ref{schemes-lemma-diagonal-immersion}), donc est
localement de type fini (Morphismes,
Lemmes \ref{morphisms-lemma-locally-finite-type-characterize} et
\ref{morphisms-lemma-immersion-locally-finite-type}). Nous pouvons donc
choisir des éléments
$f_{0, 1}, \ldots, f_{0, n} \in \mathcal{O}(U_0 \cap V_0)$
qui engendrent $\mathcal{O}(U_0 \cap V_0)$ sur
$\mathcal{O}(U_0) \otimes \mathcal{O}(V_0)$. Observons que, pour $i \geq 0$,
le diagramme de schémas
$$
\xymatrix{
U_i \cap V_i \ar[r] \ar[d] & U_i \ar[d] \\
U_0 \cap V_0 \ar[r] & U_0
}
$$
est cartésien. Nous voyons donc que les images
$f_{i, 1}, \ldots, f_{i, n} \in \mathcal{O}(U_i \cap V_i)$
engendrent $\mathcal{O}(U_i \cap V_i)$ sur
$\mathcal{O}(U_i) \otimes \mathcal{O}(V_0)$ et, a fortiori, sur
$\mathcal{O}(U_i) \otimes \mathcal{O}(V_i)$. Par hypothèse, les images
$f_1, \ldots, f_n \in \mathcal{O}(U \otimes V)$ appartiennent à l'image de
l'application
$\mathcal{O}(U) \otimes \mathcal{O}(V) \to \mathcal{O}(U \cap V)$.
Puisque
$\mathcal{O}(U) \otimes \mathcal{O}(V) =
\colim \mathcal{O}(U_i) \otimes \mathcal{O}(V_i)$
nous voyons qu'elles appartiennent à l'image de l'application à un certain
niveau fini, ce qui démontre le lemme.
\end{proof}

\begin{lemma}
\label{lemma-limit-ample}
Dans la Situation \ref{situation-descent}, soit $\mathcal{L}_0$ un faisceau
inversible de modules sur $S_0$. Si son image réciproque $\mathcal{L}$ sur
$S$ est ample, alors, pour un certain $i \in I$, son image réciproque
$\mathcal{L}_i$ sur $S_i$ est ample.
\end{lemma}

\begin{proof}
L'hypothèse signifie qu'il existe un nombre fini de sections
$s_1, \ldots, s_m \in \Gamma(S, \mathcal{L})$ telles que $S_{s_j}$ soit
affine et que $S = \bigcup S_{s_j}$ ; voir Propriétés,
Définition \ref{properties-definition-ample}. D'après le
Lemme \ref{lemma-descend-section}, nous pouvons trouver $i \in I$ et des
sections $s_{i, j} \in \Gamma(S_i, \mathcal{L}_i)$ d'images $s_j$.
Après avoir augmenté $i$, le Lemme \ref{lemma-limit-affine} permet de
supposer que $(S_i)_{s_{i, j}}$ est affine pour $j = 1, \ldots, m$.
Après avoir augmenté $i$ une dernière fois, le
Lemme \ref{lemma-descend-opens} permet de supposer que
$S_i = \bigcup (S_i)_{s_{i, j}}$. Alors $\mathcal{L}_i$ est ample par
définition.
\end{proof}

\begin{lemma}
\label{lemma-finite-type-eventually-closed}
Soit $S$ un schéma. Soit $X = \lim X_i$ une limite projective de schémas sur
$S$ dont les morphismes de transition sont affines. Soit $Y \to X$ un
morphisme de schémas sur $S$.
\begin{enumerate}
\item Si $Y \to X$ est une immersion fermée, si $X_i$ est quasi-compact et
si $Y$ est localement de type fini sur $S$, alors $Y \to X_i$ est une
immersion fermée pour $i$ assez grand.
\item Si $Y \to X$ est une immersion, si $X_i$ est quasi-séparé, si
$Y \to S$ est localement de type fini et si $Y$ est quasi-compact, alors
$Y \to X_i$ est une immersion pour $i$ assez grand.
\item Si $Y \to X$ est un isomorphisme, si $X_i$ est quasi-compact, si
$X_i \to S$ est localement de type fini, si les morphismes de transition
$X_{i'} \to X_i$ sont des immersions fermées et si $Y \to S$ est localement
de présentation finie, alors $Y \to X_i$ est un isomorphisme pour $i$ assez
grand.
\end{enumerate}
\end{lemma}

\begin{proof}
Démonstration de (1). Choisissons $0 \in I$ et un recouvrement ouvert affine
fini $X_0 = U_{0, 1} \cup \ldots \cup U_{0, m}$ tel que $U_{0, j}$ s'envoie
dans un ouvert affine $W_j \subset S$. Soient $V_j \subset Y$,
resp.\ $U_{i, j} \subset X_i$, $i \geq 0$, respectivement
$U_j \subset X$, les images réciproques de $U_{0, j}$. Il suffit de montrer
que $V_j \to U_{i, j}$ est une immersion fermée pour $i$ assez grand, et
nous savons que $V_j \to U_j$ est une immersion fermée. On se ramène ainsi
au fait algébrique suivant : si $A = \colim A_i$ est une limite inductive
filtrante de $R$-algèbres, si $A \to B$ est une surjection de $R$-algèbres
et si $B$ est une $R$-algèbre de type fini, alors $A_i \to B$ est surjectif
pour $i$ assez grand.

\medskip\noindent
Démonstration de (2). Choisissons $0 \in I$. Choisissons un ouvert
quasi-compact $X'_0 \subset X_0$ tel que $Y \to X_0$ se factorise par $X'_0$.
Après avoir remplacé $X_i$ par l'image réciproque de
$X'_0$ pour $i \geq 0$, nous pouvons supposer tous les $X_i'$ quasi-compacts et
quasi-séparés. Soit $U \subset X$ un ouvert quasi-compact tel que
$Y \to X$ se factorise par une immersion fermée $Y \to U$ (un tel $U$
existe puisque $Y$ est quasi-compact). D'après le
Lemme \ref{lemma-descend-opens}, nous pouvons supposer
$U = \lim U_i$, avec $U_i \subset X_i$ ouvert quasi-compact. L'assertion (1)
montre que $Y \to U_i$ est une immersion fermée pour un certain $i$. Ainsi,
(2) est vérifiée.

\medskip\noindent
Démonstration de (3). En travaillant localement sur des ouverts affines de
$X_0$, pour un certain $0 \in I$, comme dans la démonstration de (1), on se
ramène au fait algébrique suivant : si $A = \lim A_i$ est une limite
inductive filtrante de $R$-algèbres dont les homomorphismes de transition
sont surjectifs, et si $A$ est de présentation finie sur $A_0$, alors
$A = A_i$ pour un certain $i$. En effet, écrivons
$A = A_0/(f_1, \ldots, f_n)$. Choisissons $i$ tel que
$f_1, \ldots, f_n$ aient une image nulle par l'homomorphisme surjectif
$A_0 \to A_i$.
\end{proof}

\begin{lemma}
\label{lemma-eventually-separated}
Soit $S$ un schéma. Soit $X = \lim X_i$ une limite projective de schémas sur
$S$ dont les morphismes de transition sont affines. Supposons que
\begin{enumerate}
\item $S$ soit quasi-séparé,
\item $X_i$ soit quasi-compact et quasi-séparé,
\item $X \to S$ soit séparé.
\end{enumerate}
Alors $X_i \to S$ est séparé pour tout $i$ assez grand.
\end{lemma}

\begin{proof}
Soit $0 \in I$. Remarquons que $I$ est non vide puisqu'il est filtrant.
Comme $X_0$ est quasi-compact, nous pouvons trouver un nombre fini d'ouverts
affines $U_1, \ldots, U_n \subset S$ tels que l'image de $X_0 \to S$ soit
contenue dans $U_1 \cup \ldots \cup U_n$. Notons $h_i : X_i \to S$ le
morphisme structural. Il suffit de vérifier que, pour un certain $i \geq 0$,
les morphismes $h_i^{-1}(U_j) \to U_j$ sont séparés pour
$j = 1, \ldots,  n$. Puisque $S$ est quasi-séparé, les morphismes
$U_j \to S$ sont quasi-compacts. Par conséquent, $h_i^{-1}(U_j)$ est
quasi-compact et quasi-séparé. On se ramène ainsi au cas où $S$ est affine.
Dans ce cas, il faut montrer que $X_i$ est séparé, sachant que $X$ l'est.
Le lemme résulte donc du Lemme \ref{lemma-limit-separated}.
\end{proof}

\begin{lemma}
\label{lemma-eventually-affine}
Soit $S$ un schéma. Soit $X = \lim X_i$ une limite projective de schémas sur
$S$ dont les morphismes de transition sont affines. Supposons que
\begin{enumerate}
\item $S$ soit quasi-compact et quasi-séparé,
\item $X_i$ soit quasi-compact et quasi-séparé,
\item $X \to S$ soit affine.
\end{enumerate}
Alors $X_i \to S$ est affine pour $i$ assez grand.
\end{lemma}

\begin{proof}
Choisissons un recouvrement ouvert affine fini
$S = \bigcup_{j = 1, \ldots, n} V_j$. Notons $f : X \to S$ et
$f_i : X_i \to S$ les morphismes structuraux. Pour chaque $j$, le schéma
$f^{-1}(V_j) = \lim_i f_i^{-1}(V_j)$ est affine (car un morphisme fini est
affine par définition). D'après le Lemme \ref{lemma-limit-affine}, il existe
donc $i \in I$ tel que chaque $f_i^{-1}(V_j)$ soit affine. Autrement dit,
$f_i : X_i \to S$ est affine pour $i$ assez grand ; voir Morphismes,
Lemme \ref{morphisms-lemma-characterize-affine}.
\end{proof}

\begin{lemma}
\label{lemma-eventually-finite}
Soit $S$ un schéma. Soit $X = \lim X_i$ une limite projective de schémas sur
$S$ dont les morphismes de transition sont affines. Supposons que
\begin{enumerate}
\item $S$ soit quasi-compact et quasi-séparé,
\item $X_i$ soit quasi-compact et quasi-séparé,
\item les morphismes de transition $X_{i'} \to X_i$ soient finis,
\item $X_i \to S$ soit localement de type fini,
\item $X \to S$ soit entier.
\end{enumerate}
Alors $X_i \to S$ est fini pour $i$ assez grand.
\end{lemma}

\begin{proof}
D'après le Lemme \ref{lemma-eventually-affine}, nous pouvons supposer que
$X_i \to S$ est affine pour tout $i$. Choisissons un recouvrement ouvert
affine fini $S = \bigcup_{j = 1, \ldots, n} V_j$. Notons $f : X \to S$ et
$f_i : X_i \to S$ les morphismes structuraux. Il suffit de montrer qu'il
existe $i$ tel que $f_i^{-1}(V_j)$ soit fini sur $V_j$ pour
$j = 1, \ldots, m$ (Morphismes,
Lemme \ref{morphisms-lemma-finite-local}). En effet, pour $i' \geq i$, la
composée $X_{i'} \to X_i \to S$ sera finie en tant que composée de morphismes
finis (Morphismes, Lemme \ref{morphisms-lemma-composition-finite}). On se
ramène ainsi au cas affine. Soit $R$ un anneau et soit $A = \colim A_i$, avec
$R \to A$ entier et $A_i \to A_{i'}$ fini pour tout $i \leq i'$. De plus,
$R \to A_i$ est de type fini pour tout $i$. Il faut montrer que $A_i$ est
fini sur $R$ pour un certain $i$. Pour cela, choisissons $i \in I$ et des
générateurs $x_1, \ldots, x_m \in A_i$ de $A_i$ comme $R$-algèbre. Puisque
$A$ est entier sur $R$, nous pouvons trouver des polynômes unitaires
$P_j \in R[T]$ tels que $P_j(x_j) = 0$ dans $A$. Il existe donc
$i' \geq i$ tel que $P_j(x_j) = 0$ dans $A_{i'}$ pour
$j = 1, \ldots, m$. Alors l'image $A'_i$ de $A_i$ dans $A_{i'}$ est finie
sur $R$ d'après Algèbre,
Lemme \ref{algebra-lemma-characterize-finite-in-terms-of-integral}. Comme
$A'_i \subset A_{i'}$ est également fini, on en déduit qu'$A_{i'}$ est fini
sur $R$ par Algèbre, Lemme \ref{algebra-lemma-finite-transitive}.
\end{proof}

\begin{lemma}
\label{lemma-eventually-closed-immersion}
Soit $S$ un schéma. Soit $X = \lim X_i$ une limite projective de schémas sur
$S$ dont les morphismes de transition sont affines. Supposons que
\begin{enumerate}
\item $S$ soit quasi-compact et quasi-séparé,
\item $X_i$ soit quasi-compact et quasi-séparé,
\item les morphismes de transition $X_{i'} \to X_i$ soient des immersions
fermées,
\item $X_i \to S$ soit localement de type fini,
\item $X \to S$ soit une immersion fermée.
\end{enumerate}
Alors $X_i \to S$ est une immersion fermée pour $i$ assez grand.
\end{lemma}

\begin{proof}
D'après le Lemme \ref{lemma-eventually-affine}, nous pouvons supposer que
$X_i \to S$ est affine pour tout $i$. Choisissons un recouvrement ouvert
affine fini $S = \bigcup_{j = 1, \ldots, n} V_j$. Notons $f : X \to S$ et
$f_i : X_i \to S$ les morphismes structuraux. Il suffit de montrer qu'il
existe $i$ tel que $f_i^{-1}(V_j)$ soit un sous-schéma fermé de $V_j$ pour
$j = 1, \ldots, m$ (Morphismes,
Lemme \ref{morphisms-lemma-closed-immersion}). On se ramène ainsi au cas
affine. Soit $R$ un anneau et soit $A = \colim A_i$, avec $R \to A$
surjectif et $A_i \to A_{i'}$ surjectif pour tout $i \leq i'$. De plus,
$R \to A_i$ est de type fini pour tout $i$. Il faut montrer que
$R \to A_i$ est surjectif pour un certain $i$. Pour cela, choisissons
$i \in I$ et des générateurs $x_1, \ldots, x_m \in A_i$ de $A_i$ comme
$R$-algèbre. Puisque $R \to A$ est surjectif, nous pouvons trouver
$r_j \in R$ tels que $r_j$ ait pour image $x_j$ dans $A$. Il existe donc $i' \geq i$ tel
que $r_j$ ait pour image celle de $x_j$ dans $A_{i'}$ pour
$j = 1, \ldots, m$. Comme $A_i \to A_{i'}$ est surjectif, il en résulte que
$R \to A_{i'}$ est surjectif.
\end{proof}

\begin{lemma}
\label{lemma-eventually-immersion}
Soit $S$ un schéma. Soit $X = \lim X_i$ une limite projective de schémas sur
$S$ dont les morphismes de transition sont affines. Supposons que
\begin{enumerate}
\item $S$ soit quasi-séparé,
\item $X_i$ soit quasi-compact et quasi-séparé,
\item les morphismes de transition $X_{i'} \to X_i$ soient des immersions
fermées,
\item $X_i \to S$ soit localement de type fini, et
\item $X \to S$ soit une immersion.
\end{enumerate}
Alors $X_i \to S$ est une immersion pour $i$ assez grand.
\end{lemma}

\begin{proof}
Choisissons un sous-schéma ouvert $U \subset S$ tel que $X \to S$ soit la
composée d'une immersion fermée $X \to U$ et du morphisme d'inclusion
$U \to S$. Puisque $X$ est quasi-compact, nous pouvons rétrécir $U$ et
supposer $U$ quasi-compact. Notons $V_i \subset X_i$ l'image réciproque de $U$.
Comme l'image réciproque de $V_i$ est $X$, le
Lemme \ref{lemma-descend-opens} montre que $V_i = X_i$ pour tout $i$ assez
grand. Nous pouvons donc supposer $X = \lim X_i$ dans la catégorie des
schémas sur $U$. Alors $X_i \to U$ est une immersion fermée pour $i$ assez grand,
d'après le Lemme \ref{lemma-eventually-closed-immersion}. Cela
démontre le lemme.
\end{proof}










\section{Approximation noethérienne absolue}
\label{section-approximation}

\noindent
Une bonne référence pour cette section est l'appendice C de l'article de
Thomason et Trobaugh \cite{TT}. Pour nos conventions sur les systèmes
filtrants, voir Catégories, section \ref{categories-section-posets-limits}.
Nous utiliserons sans autre mention le résultat d'existence et les propriétés
de la limite établis à la section \ref{section-limits}.

\begin{lemma}
\label{lemma-quasi-affine-finite-type-over-Z}
Soit $W$ un schéma quasi-affine de type fini sur $\mathbf{Z}$. Supposons que
$W \to \Spec(R)$ soit une immersion ouverte dans un schéma affine. Il existe
une $\mathbf{Z}$-algèbre de type fini $A \subset R$ qui induit une immersion
ouverte $W \to \Spec(A)$. De plus, $R$ est la limite inductive filtrante de
telles sous-algèbres.
\end{lemma}

\begin{proof}
Choisissons un recouvrement ouvert affine
$W = \bigcup_{i = 1, \ldots, n} W_i$ tel que chaque $W_i$ soit un ouvert
affine standard de $\Spec(R)$. Autrement dit, si l'on écrit
$W_i = \Spec(R_i)$, alors $R_i = R_{f_i}$ pour un certain $f_i \in R$.
Choisissons un nombre fini d'éléments $x_{ij} \in R_i$ qui engendrent $R_i$
sur $\mathbf{Z}$. Prenons $N \gg 0$ tel que chaque $f_i^Nx_{ij}$ provienne
d'un élément de $R$, disons $y_{ij} \in R$. Soit $A$ la
$\mathbf{Z}$-algèbre engendrée par les $f_i$ et les $y_{ij}$, ainsi que,
éventuellement, par un nombre fini d'éléments supplémentaires de $R$.
Alors $A$ convient. Les détails sont omis.
\end{proof}

\begin{lemma}
\label{lemma-diagram}
Supposons donné un diagramme cartésien d'anneaux
$$
\xymatrix{
B \ar[r]_s & R \\
B'\ar[u] \ar[r] & R' \ar[u]_t
}
$$
Soit $W' \subset \Spec(R')$ un ouvert de la forme
$W' = D(f_1) \cup \ldots \cup D(f_n)$ tel que $t(f_i) = s(g_i)$ pour un
certain $g_i \in B$ et que $B_{g_i} \cong R_{s(g_i)}$. Alors $B' \to R'$
induit une immersion ouverte de $W'$ dans $\Spec(B')$.
\end{lemma}

\begin{proof}
Posons $h_i = (g_i, f_i) \in B'$. Compléments d'algèbre,
Lemme \ref{more-algebra-lemma-diagram-localize}, montre que
$(B')_{h_i} \cong (R')_{f_i}$, comme voulu.
\end{proof}

\noindent
Le lemme suivant donne un énoncé précis de l'approximation noethérienne.

\begin{lemma}
\label{lemma-approximate}
Soit $S$ un schéma quasi-compact et quasi-séparé. Soit $V \subset S$ un
ouvert quasi-compact. Soit $I$ un ensemble ordonné filtrant et soit
$(V_i, f_{ii'})$ un système projectif de schémas indexé par $I$, à
morphismes de transition affines, chaque $V_i$ étant de type fini sur
$\mathbf{Z}$, et tel que $V = \lim V_i$. Il existe alors
\begin{enumerate}
\item un ensemble ordonné filtrant $J$,
\item un système projectif de schémas $(S_j, g_{jj'})$ indexé par $J$,
\item une application croissante $\alpha : J \to I$,
\item des sous-schémas ouverts $V'_j \subset S_j$, et
\item des isomorphismes $V'_j \to V_{\alpha(j)}$
\end{enumerate}
tels que
\begin{enumerate}
\item les morphismes de transition $g_{jj'} : S_j \to S_{j'}$ sont affines,
\item chaque $S_j$ est de type fini sur $\mathbf{Z}$,
\item $g_{jj'}^{-1}(V'_{j'}) = V'_j$,
\item $S = \lim S_j$ et $V = \lim V'_j$, et
\item les diagrammes
$$
\vcenter{
\xymatrix{
V \ar[d] \ar[rd] \\
V'_j \ar[r] & V_{\alpha(j)}
}
}
\quad\text{et}\quad
\vcenter{
\xymatrix{
V'_j \ar[r] \ar[d] & V_{\alpha(j)} \ar[d] \\
V'_{j'} \ar[r] & V_{\alpha(j')}
}
}
$$
sont commutatifs.
\end{enumerate}
\end{lemma}

\begin{proof}
Posons $Z = S \setminus V$. Choisissons des ouverts affines
$U_1, \ldots, U_m \subset S$ tels que
$Z \subset \bigcup_{l = 1, \ldots, m} U_l$. Considérons les ouverts
$$
V \subset V \cup U_1 \subset V \cup U_1 \cup U_2 \subset
\ldots \subset V \cup \bigcup\nolimits_{l = 1, \ldots, m} U_l = S
$$
Si nous pouvons démontrer le lemme successivement dans chacun des cas
$$
V \cup U_1 \cup \ldots \cup U_l
\subset
V \cup U_1 \cup \ldots \cup U_{l + 1}
$$
alors le lemme en résultera pour $V \subset S$. Dans chaque cas, nous
ajoutons un seul ouvert affine. Nous pouvons donc supposer que
\begin{enumerate}
\item $S = U \cup V$,
\item $U$ est un ouvert affine de $S$,
\item $V$ est un ouvert quasi-compact de $S$, et
\item $V = \lim_i V_i$, où $(V_i, f_{ii'})$ est un système projectif indexé
par un ensemble ordonné filtrant $I$, chaque $f_{ii'}$ étant affine et
chaque $V_i$ de type fini sur $\mathbf{Z}$.
\end{enumerate}
Notons $f_i : V \to V_i$ les projections. Posons $W = U \cap V$. Comme $S$
est quasi-séparé, c'est un ouvert quasi-compact de $V$. D'après le
Lemme \ref{lemma-descend-opens} (et après avoir restreint $I$), nous pouvons
supposer qu'il existe des ouverts $W_i \subset V_i$ tels que
$f_{ii'}^{-1}(W_{i'}) = W_i$ et $f_i^{-1}(W_i) = W$. Puisque $W$ est un
ouvert quasi-compact de $U$, il est quasi-affine. Nous pouvons donc supposer,
après avoir restreint $I$ une nouvelle fois, que $W_i$ est quasi-affine pour
tout $i$ ; voir le Lemme \ref{lemma-limit-quasi-affine}.

\medskip\noindent
Écrivons $U = \Spec(B)$. Posons $R = \Gamma(W, \mathcal{O}_W)$ et
$R_i = \Gamma(W_i, \mathcal{O}_{W_i})$. D'après le
Lemme \ref{lemma-descend-section}, on a $R = \colim_i R_i$. Nous disposons
maintenant des homomorphismes d'anneaux
$$
\xymatrix{
B \ar[r]_s & R \\
& R_i \ar[u]_{t_i}
}
$$
Posons $B_i = \{(b, r) \in B \times R_i \mid s(b) = t_i(r)\}$, de sorte que
nous ayons un diagramme cartésien
$$
\xymatrix{
B \ar[r]_s & R \\
B_i \ar[u] \ar[r] & R_i \ar[u]_{t_i}
}
$$
pour chaque $i$. Les homomorphismes de transition $R_i \to R_{i'}$
induisent des homomorphismes $B_i \to B_{i'}$. Il est clair que
$B = \colim_i B_i$. Au paragraphe suivant, nous montrons que, pour tout $i$
assez grand, la composée $W_i \to \Spec(R_i) \to \Spec(B_i)$ est une
immersion ouverte.

\medskip\noindent
Comme $W$ est un ouvert quasi-compact de $U = \Spec(B)$, nous pouvons trouver
un nombre fini d'éléments $g_l \in B$, $l = 1, \ldots, m$, tels que
$D(g_l) \subset W$ et $W = \bigcup_{l = 1, \ldots, m} D(g_l)$. Remarquons
que cela implique $D(g_l) = W_{s(g_l)}$ comme ouverts de $U$, où
$W_{s(g_l)}$ désigne le plus grand ouvert de $W$ sur lequel $s(g_l)$ est
inversible. Par conséquent,
$$
B_{g_l} =
\Gamma(D(g_l), \mathcal{O}_U) =
\Gamma(W_{s(g_l)}, \mathcal{O}_W) = R_{s(g_l)},
$$
où la dernière égalité découle de Propriétés,
Lemme \ref{properties-lemma-invert-f-sections}. Comme $W_{s(g_l)}$ est
affine, cela implique aussi que $D(s(g_l)) = W_{s(g_l)}$ comme ouverts de
$\Spec(R)$. Puisque $R = \colim_i R_i$, nous pouvons, après avoir restreint
$I$, supposer qu'il existe des $g_{l, i} \in R_i$ pour tout $i \in I$ tels
que $s(g_l) = t_i(g_{l, i})$. Bien entendu, nous choisissons les $g_{l, i}$
de sorte que $g_{l, i}$ s'envoie sur $g_{l, i'}$ par les homomorphismes de
transition $R_i \to R_{i'}$. D'après le Lemme \ref{lemma-descend-opens}, nous
pouvons alors, après avoir restreint $I$ une nouvelle fois, supposer que les
ouverts correspondants $D(g_{l, i}) \subset \Spec(R_i)$ sont contenus dans
$W_i$, pour $l = 1, \ldots, m$, et recouvrent $W_i$. Nous en déduisons que
le morphisme $W_i \to \Spec(R_i) \to \Spec(B_i)$ est une immersion ouverte ;
voir le Lemme \ref{lemma-diagram}.

\medskip\noindent
D'après le Lemme \ref{lemma-quasi-affine-finite-type-over-Z}, nous pouvons
écrire $B_i$ comme une limite inductive filtrante de sous-algèbres
$A_{i, p} \subset B_i$, $p \in P_i$, chacune de type fini sur $\mathbf{Z}$,
et telles que $W_i$ s'identifie à un sous-schéma ouvert de
$\Spec(A_{i, p})$. Soit $S_{i, p}$ le schéma obtenu en recollant $V_i$ et
$\Spec(A_{i, p})$ le long de l'ouvert $W_i$ ; voir Schémas,
section \ref{schemes-section-glueing-schemes}. Voici le diagramme commutatif
de schémas qui en résulte :
$$
\xymatrix{
& & V \ar[lld] \ar[d] & W \ar[l] \ar[lld] \ar[d] \\
V_i \ar[d] & W_i \ar[l] \ar[d] & S \ar[lld] & U \ar[lld] \ar[l] \\
S_{i, p} & \Spec(A_{i, p}) \ar[l]
}
$$
Le morphisme $S \to S_{i, p}$ existe parce que le carré supérieur droit est
un coproduit amalgamé dans la catégorie des schémas. Remarquons que
$S_{i, p}$ est de type fini sur $\mathbf{Z}$ puisqu'il admet un recouvrement
ouvert affine fini dont les membres sont les spectres de
$\mathbf{Z}$-algèbres de type fini. Nous définissons un préordre sur
$J = \coprod_{i \in I} P_i$ par la règle suivante :
$(i', p') \geq (i, p)$ si et seulement si $i' \geq i$ et si
l'homomorphisme $B_i \to B_{i'}$ envoie $A_{i, p}$ dans $A_{i', p'}$.
C'est exactement la condition requise pour définir un morphisme
$S_{i', p'} \to S_{i, p}$ : on forme en effet un diagramme commutatif comme
ci-dessus à l'aide des morphismes de transition $V_{i'} \to V_i$ et
$W_{i'} \to W_i$, ainsi que du morphisme
$\Spec(A_{i', p'}) \to \Spec(A_{i, p})$ induit par l'homomorphisme d'anneaux
$A_{i, p} \to A_{i', p'}$. Les commutativités nécessaires ont été intégrées
aux constructions. Nous affirmons que $S$ est la limite projective des
schémas $S_{i, p}$. Puisque, par construction, les schémas $V_i$ ont pour
limite $V$, cela se ramène au fait que $B$ est la limite des anneaux
$A_{i, p}$, ce qui est vrai par construction. L'application
$\alpha : J \to I$ est donnée par la règle $j = (i, p) \mapsto i$. Le
sous-schéma ouvert $V'_j$ est simplement l'image de $V_i \to S_{i, p}$
ci-dessus. La commutativité des diagrammes de (5) résulte clairement de la
construction. Cela achève la démonstration du lemme.
\end{proof}

\begin{proposition}
\label{proposition-approximate}
Soit $S$ un schéma quasi-compact et quasi-séparé. Il existe un ensemble
ordonné filtrant $I$ et un système projectif de schémas $(S_i, f_{ii'})$
indexé par $I$ tels que
\begin{enumerate}
\item les morphismes de transition $f_{ii'}$ sont affines,
\item chaque $S_i$ est de type fini sur $\mathbf{Z}$, et
\item $S = \lim_i S_i$.
\end{enumerate}
\end{proposition}

\begin{proof}
C'est un cas particulier du Lemme \ref{lemma-approximate}, avec
$V = \emptyset$.
\end{proof}






\section{Limites et morphismes de présentation finie}
\label{section-finite-presentation}

\noindent
Le résultat suivant généralise Algèbre,
Lemme \ref{algebra-lemma-characterize-finite-presentation}.

\begin{proposition}
\label{proposition-characterize-locally-finite-presentation}
\begin{reference}
\cite[IV, Proposition 8.14.2]{EGA}
\end{reference}
Soit $f : X \to S$ un morphisme de schémas. Les conditions suivantes sont
équivalentes :
\begin{enumerate}
\item Le morphisme $f$ est localement de présentation finie.
\item Pour tout ensemble ordonné filtrant $I$ et tout système projectif
$(T_i, f_{ii'})$ de $S$-schémas indexé par $I$, chaque $T_i$ étant affine,
on a
$$
\Mor_S(\lim_i T_i, X) =
\colim_i \Mor_S(T_i, X)
$$
\item Pour tout ensemble ordonné filtrant $I$ et tout système projectif
$(T_i, f_{ii'})$ de $S$-schémas indexé par $I$, chaque $f_{ii'}$ étant affine
et chaque $T_i$ étant quasi-compact et quasi-séparé comme schéma, on a
$$
\Mor_S(\lim_i T_i, X) =
\colim_i \Mor_S(T_i, X)
$$
\end{enumerate}
\end{proposition}

\begin{proof}
Il est clair que (3) implique (2).

\medskip\noindent
Montrons que (2) implique (1). Supposons (2). Choisissons des ouverts affines
$U \subset X$ et $V \subset S$ tels que $f(U) \subset V$. Il faut montrer
que $\mathcal{O}_S(V) \to \mathcal{O}_X(U)$ est de présentation finie.
Soit $(A_i, \varphi_{ii'})$ un système inductif filtrant
d'$\mathcal{O}_S(V)$-algèbres. Posons $A = \colim_i A_i$. D'après Algèbre,
Lemme \ref{algebra-lemma-characterize-finite-presentation}, il faut montrer
que
$$
\Hom_{\mathcal{O}_S(V)}(\mathcal{O}_X(U), A) =
\colim_i \Hom_{\mathcal{O}_S(V)}(\mathcal{O}_X(U), A_i)
$$
Considérons les schémas $T_i = \Spec(A_i)$. Ils forment un système projectif
de $V$-schémas indexé par $I$, dont les morphismes de transition
$f_{ii'} : T_i \to T_{i'}$ sont induits par les homomorphismes
d'$\mathcal{O}_S(V)$-algèbres $\varphi_{i'i}$. Posons
$T := \Spec(A) = \lim_i T_i$. En termes d'ensembles de morphismes de
schémas, la formule ci-dessus devient
$$
\Mor_V(\lim_i T_i, U) =
\colim_i \Mor_V(T_i, U).
$$
Observons d'abord que
$\Mor_V(T_i, U) = \Mor_S(T_i, U)$
et
$\Mor_V(T, U) = \Mor_S(T, U)$.
Il faut donc montrer que
$$
\Mor_S(\lim_i T_i, U) =
\colim_i \Mor_S(T_i, U)
$$
et l'on sait que
$$
\Mor_S(\lim_i T_i, X) =
\colim_i \Mor_S(T_i, X).
$$
Il suffit donc de prouver ceci : étant donné un morphisme
$g_i : T_i \to X$ sur $S$ tel que la composée $T \to T_i \to X$ ait son
image dans $U$, il existe $i' \geq i$ tel que la composée
$g_{i'} : T_{i'} \to T_i \to X$ ait son image dans $U$. Notons
$Z_{i'} = g_{i'}^{-1}(X \setminus U)$. Supposons tous les $Z_{i'}$ non vides,
en vue d'une contradiction. D'après le
Lemme \ref{lemma-limit-closed-nonempty}, il existe un point $t$ de $T$ qui
s'envoie dans $Z_{i'}$ pour tout $i' \geq i$. Un tel point ne s'envoie pas
dans $U$, contradiction.

\medskip\noindent
Enfin, montrons que (1) implique (3). Supposons (1). Soit donné un système
projectif filtrant $(T_i, f_{ii'})$ de $S$-schémas. Supposons les morphismes
$f_{ii'}$ affines et chaque $T_i$ quasi-compact et quasi-séparé comme schéma.
Soit $T = \lim_i T_i$. Notons $f_i : T \to T_i$ les morphismes de
projection. Il faut montrer :
\begin{enumerate}
\item[(a)] étant donnés des morphismes $g_i, g'_i : T_i \to X$ sur $S$ tels
que $g_i \circ f_i = g'_i \circ f_i$, il existe $i' \geq i$ tel que
$g_i \circ f_{i'i} = g'_i \circ f_{i'i}$ ;
\item[(b)] étant donné tout morphisme $g : T \to X$ sur $S$, il existe
$i \in I$ et un morphisme $g_i : T_i \to X$ tels que
$g = f_i \circ g_i$.
\end{enumerate}

\noindent
Montrons d'abord l'unicité (a). Soient $g_i, g'_i : T_i \to X$ des
morphismes tels que $g_i \circ f_i = g'_i \circ f_i$. Pour tout
$i' \geq i$, posons $g_{i'} = g_i \circ f_{i'i}$ et
$g'_{i'} = g'_i \circ f_{i'i}$. Posons également
$g = g_i \circ f_i = g'_i \circ f_i$. Considérons le morphisme
$(g_i, g'_i) : T_i \to X \times_S X$. Posons
$$
W =
\bigcup\nolimits_{U \subset X\text{ ouvert affine},
V \subset S\text{ ouvert affine}, f(U) \subset V}
U \times_V U.
$$
C'est un ouvert de $X \times_S X$ tel que le morphisme $\Delta_{X/S}$ se
factorise par une immersion fermée dans $W$ ; voir la démonstration de
Schémas, Lemme \ref{schemes-lemma-diagonal-immersion}. Remarquons que la
composée $(g_i, g'_i) \circ f_i : T \to X \times_S X$ est un morphisme dans
$W$, puisqu'elle se factorise par la diagonale par hypothèse. Posons
$Z_{i'} = (g_{i'}, g'_{i'})^{-1}(X \times_S X \setminus W)$. Si tous les
$Z_{i'}$ sont non vides, le Lemme \ref{lemma-limit-closed-nonempty} fournit
un point $t \in T$ qui s'envoie dans $Z_{i'}$ pour tout $i' \geq i$. Cela
contredit le fait que $T$ s'envoie dans $W$. Nous pouvons donc augmenter $i$
et supposer que $(g_i, g'_i) : T_i \to X \times_S X$ est un morphisme dans
$W$. Par construction de $W$, et puisque $T_i$ est quasi-compact, nous
pouvons trouver un recouvrement ouvert affine fini
$T_i = T_{1, i} \cup \ldots \cup T_{n, i}$ tel que
$(g_i, g'_i)|_{T_{j, i}}$ soit un morphisme dans $U \times_V U$ pour une
certaine paire $(U, V)$ intervenant dans la définition de $W$ ci-dessus.
Comme il suffit de prouver que $g_{i'}$ et $g'_{i'}$ coïncident sur chacun
des $f_{i'i}^{-1}(T_{j, i})$, on se ramène au cas affine. Celui-ci résulte
d'Algèbre, Lemme \ref{algebra-lemma-characterize-finite-presentation}, et du
fait que l'homomorphisme d'anneaux
$\mathcal{O}_S(V) \to \mathcal{O}_X(U)$ est de présentation finie (voir
Morphismes,
Lemme \ref{morphisms-lemma-locally-finite-presentation-characterize}).

\medskip\noindent
Enfin, montrons l'existence (b). Soit $g : T \to X$ un morphisme de schémas
sur $S$. Nous pouvons trouver un recouvrement ouvert affine fini
$T = W_1 \cup \ldots \cup W_n$ tel que, pour tout
$j \in \{1, \ldots, n\}$, il existe des ouverts affines $U_j \subset X$ et
$V_j \subset S$ avec $f(U_j) \subset V_j$ et $g(W_j) \subset U_j$.
D'après les Lemmes \ref{lemma-descend-opens} et \ref{lemma-limit-affine}, et
après avoir éventuellement restreint $I$, nous pouvons supposer qu'il existe
des recouvrements ouverts affines
$T_i = W_{1, i} \cup \ldots \cup W_{n, i}$ compatibles aux morphismes de
transition et tels que $W_j = \lim_i W_{j, i}$. Appliquons Algèbre,
Lemme \ref{algebra-lemma-characterize-finite-presentation}, aux anneaux
correspondant aux schémas affines $U_j$, $V_j$, $W_{j, i}$ et $W_j$, en
utilisant que $\mathcal{O}_S(V_j) \to \mathcal{O}_X(U_j)$ est de présentation
finie (voir Morphismes,
Lemme \ref{morphisms-lemma-locally-finite-presentation-characterize}). Pour
chaque $j$, nous pouvons ainsi trouver un indice $i_j \in I$ et un morphisme
$g_{j, i_j} : W_{j, i_j} \to X$ tels que
$g_{j, i_j} \circ f_i|_{W_j} : W_j \to W_{j, i} \to X$ soit égal à
$g|_{W_j}$. Par l'assertion (a) démontrée ci-dessus, et en utilisant la
quasi-compacité de $W_{j_1, i} \cap W_{j_2, i}$, qui résulte de ce que $T_i$
est quasi-séparé, nous pouvons trouver un indice $i' \in I$ supérieur à tous
les $i_j$ tel que
$$
g_{j_1, i_{j_1}} \circ f_{i'i_{j_1}}|_{W_{j_1, i'} \cap W_{j_2, i'}} =
g_{j_2, i_{j_2}} \circ f_{i'i_{j_2}}|_{W_{j_1, i'} \cap W_{j_2, i'}}
$$
pour tous $j_1, j_2 \in \{1, \ldots, n\}$. Les morphismes
$g_{j, i_j} \circ f_{i'i_j}|_{W_{j, i'}}$ se recollent donc pour donner le
morphisme cherché $T_{i'} \to X$.
\end{proof}

\begin{remark}
\label{remark-limit-preserving}
Soit $S$ un schéma. Disons qu'un foncteur
$F : (\Sch/S)^{opp} \to \textit{Ensembles}$ est
{\it compatible aux limites} si, pour tout système projectif filtrant
$\{T_i\}_{i \in I}$ de schémas affines de limite $T$, on a
$F(T) = \colim_i F(T_i)$. Soit $X$ un schéma sur $S$ et soit
$h_X : (\Sch/S)^{opp} \to \textit{Ensembles}$ son foncteur des points ; voir
Schémas, section \ref{schemes-section-representable}. Dans cette
terminologie, la
Proposition \ref{proposition-characterize-locally-finite-presentation}
affirme qu'un schéma $X$ est localement de présentation finie sur $S$ si et
seulement si $h_X$ est compatible aux limites.
\end{remark}

\begin{lemma}
\label{lemma-surjection-is-enough}
Soit $f : X \to S$ un morphisme de schémas. Si, pour toute limite projective
$T = \lim_{i \in I} T_i$ de schémas affines sur $S$, l'application
$$
\colim \Mor_S(T_i, X) \longrightarrow \Mor_S(T, X)
$$
est surjective, alors $f$ est localement de présentation finie. Autrement
dit, dans les assertions (2) et (3) de la
Proposition \ref{proposition-characterize-locally-finite-presentation}, il
suffit de vérifier la surjectivité de l'application.
\end{lemma}

\begin{proof}
La démonstration est exactement la même que celle de l'implication
« (2) implique (1) » dans la
Proposition \ref{proposition-characterize-locally-finite-presentation}.
Choisissons des ouverts affines $U \subset X$ et $V \subset S$ tels que
$f(U) \subset V$. Il faut montrer que
$\mathcal{O}_S(V) \to \mathcal{O}_X(U)$ est de présentation finie. Soit
$(A_i, \varphi_{ii'})$ un système inductif filtrant d'$\mathcal{O}_S(V)$-algèbres.
Posons $A = \colim_i A_i$. D'après Algèbre,
Lemme \ref{algebra-lemma-characterize-finite-presentation}, il suffit de
montrer que
$$
\colim_i \Hom_{\mathcal{O}_S(V)}(\mathcal{O}_X(U), A_i) \to
\Hom_{\mathcal{O}_S(V)}(\mathcal{O}_X(U), A)
$$
est surjective. Considérons les schémas $T_i = \Spec(A_i)$. Ils forment un
système projectif de $V$-schémas indexé par $I$, dont les morphismes de
transition $f_{ii'} : T_i \to T_{i'}$ sont induits par les homomorphismes
d'$\mathcal{O}_S(V)$-algèbres $\varphi_{i'i}$. Posons
$T := \Spec(A) = \lim_i T_i$. En termes d'ensembles de morphismes de
schémas, la formule ci-dessus devient
$$
\colim_i \Mor_V(T_i, U) \to \Mor_V(\lim_i T_i, U)
$$
Observons d'abord que
$\Mor_V(T_i, U) = \Mor_S(T_i, U)$
et
$\Mor_V(T, U) = \Mor_S(T, U)$.
Il faut donc montrer que
$$
\colim_i \Mor_S(T_i, U) \to
\Mor_S(\lim_i T_i, U)
$$
est surjective, et l'on sait que
$$
\colim_i \Mor_S(T_i, X) \to
\Mor_S(\lim_i T_i, X)
$$
est surjective. Il suffit donc de prouver ceci : étant donné un morphisme
$g_i : T_i \to X$ sur $S$ tel que la composée $T \to T_i \to X$ ait son
image dans $U$, il existe $i' \geq i$ tel que la composée
$g_{i'} : T_{i'} \to T_i \to X$ ait son image dans $U$. Notons
$Z_{i'} = g_{i'}^{-1}(X \setminus U)$. Supposons tous les $Z_{i'}$ non vides,
en vue d'une contradiction. D'après le
Lemme \ref{lemma-limit-closed-nonempty}, il existe un point $t$ de $T$ qui
s'envoie dans $Z_{i'}$ pour tout $i' \geq i$. Un tel point ne s'envoie pas
dans $U$, contradiction.
\end{proof}

\noindent
Voici un exemple d'application de la
Proposition \ref{proposition-characterize-locally-finite-presentation}.

\begin{lemma}
\label{lemma-morphism-glueing-near-closed-point}
Soit $S$ un schéma. Soient $X$ et $Y$ des schémas sur $S$. Supposons que
$Y$ soit localement de présentation finie sur $S$. Soit $x \in X$ un point
fermé tel que $U = X \setminus \{x\} \to X$ soit quasi-compact. En posant
$V = \Spec(\mathcal{O}_{X, x}) \setminus \{x\}$, on a une bijection
$$
\left\{
\begin{matrix}
\text{morphismes }X \to Y\text{ sur }S
\end{matrix}
\right\}
\longrightarrow
\left\{
\begin{matrix}
(a, b)\text{ où }
a : U \to Y\text{ et }b : \Spec(\mathcal{O}_{X, x}) \to Y\\
\text{ sont des morphismes sur }S
\text{ qui coïncident sur }V
\end{matrix}
\right\}
$$
\end{lemma}

\begin{proof}
Soit $W \subset X$ un voisinage ouvert de $x$. Par recollement des schémas,
voir Schémas, section \ref{schemes-section-glueing-schemes}, le résultat est
valable si l'on considère des paires de morphismes $a : U \to Y$ et
$c : W \to Y$ qui coïncident sur $U \cap W$. On a
$\mathcal{O}_{X, x} = \colim \mathcal{O}_W(W)$, où $W$ parcourt les
voisinages ouverts affines de $x$ dans $X$. Ainsi,
$\Spec(\mathcal{O}_{X, x}) = \lim W$, où $W$ parcourt les voisinages ouverts
affines de $s$. D'après la
Proposition \ref{proposition-characterize-locally-finite-presentation}, tout
morphisme $b : \Spec(\mathcal{O}_{X, x}) \to Y$ sur $S$ provient donc d'un
morphisme $c : W \to Y$ pour un certain $W$ comme ci-dessus (et $c$ est
unique après un éventuel rétrécissement supplémentaire de $W$). Pour tout
ouvert affine $x \in W$, l'ouvert $U \cap W$ est quasi-compact puisque
$U \to X$ est quasi-compact. Par conséquent,
$V = \lim W \cap U = \lim W \setminus \{x\}$ est une limite de schémas
quasi-compacts et quasi-séparés (voir le
Lemme \ref{lemma-directed-inverse-system-has-limit}). Ainsi, si $a$ et $b$
coïncident sur $V$, alors, après avoir rétréci $W$, $a$ et $c$ coïncident sur
$U \cap W$ (par la même proposition). Le lemme en résulte.
\end{proof}







\section{Approximation relative}
\label{section-relative-approximation}

\noindent
Nous examinons des variantes sur une base de la
Proposition \ref{proposition-approximate}.

\begin{lemma}
\label{lemma-approximate-morphism}
Soit $f : X \to S$ un morphisme de schémas quasi-compacts et quasi-séparés.
Il existe alors un ensemble ordonné filtrant $I$ et un système projectif
$(f_i : X_i \to S_i)$ de morphismes de schémas indexé par $I$, tels que les
morphismes de transition $X_i \to X_{i'}$ et $S_i \to S_{i'}$ soient
affines, que $X_i$ et $S_i$ soient de type fini sur $\mathbf{Z}$, et que
$(X \to S) = \lim (X_i \to S_i)$.
\end{lemma}

\begin{proof}
Écrivons $X = \lim_{a \in A} X_a$ et $S = \lim_{b \in B} S_b$ comme dans la
Proposition \ref{proposition-approximate}, c'est-à-dire avec $X_a$ et $S_b$
de type fini sur $\mathbf{Z}$ et des morphismes de transition affines.

\medskip\noindent
Fixons $b \in B$. En appliquant la
Proposition \ref{proposition-characterize-locally-finite-presentation} à
$S_b$ et à $X = \lim X_a$ sur $\mathbf{Z}$, nous trouvons $a \in A$ et un
morphisme $f_{a, b} : X_a \to S_b$ rendant commutatif le diagramme
$$
\xymatrix{
X \ar[d] \ar[r] & S \ar[d] \\
X_a \ar[r] & S_b
}
$$
Soit $I$ l'ensemble des triplets $(a, b, f_{a, b})$ ainsi obtenus.

\medskip\noindent
Soient $(a, b, f_{a, b})$ et $(a', b', f_{a', b'})$ dans $I$. Soit
$b'' \leq \min(b, b')$. Une nouvelle application de la
Proposition \ref{proposition-characterize-locally-finite-presentation}
fournit $a'' \geq \max(a, a')$ tel que les composées
$X_{a''} \to X_a \to S_b \to S_{b''}$ et
$X_{a''} \to X_{a'} \to S_{b'} \to S_{b''}$ soient égales. Munissons $I$
du préordre
$$
(a, b, f_{a, b}) \geq (a', b', f_{a', b'})
\Leftrightarrow
a \geq a',\ b \geq b',\text{ et }
g_{b, b'} \circ f_{a, b} = f_{a', b'} \circ h_{a, a'}
$$
où $h_{a, a'} : X_a \to X_{a'}$ et $g_{b, b'} : S_b \to S_{b'}$ sont les
morphismes de transition. Les remarques précédentes montrent que $I$ est
filtrant et que les applications $I \to A$,
$(a, b, f_{a, b}) \mapsto a$, et $I \to B$, $(a, b, f_{a, b})$, sont
cofinales. Si, pour $i = (a, b, f_{a, b})$, nous posons $X_i = X_a$,
$S_i = S_b$ et $f_i = f_{a, b}$, nous obtenons un système projectif de
morphismes indexé par $I$, et nous avons
$$
\lim_{i \in I} X_i = \lim_{a \in A} X_a = X
\quad\text{et}\quad
\lim_{i \in I} S_i = \lim_{b \in B} S_b = S
$$
d'après Catégories, Lemme \ref{categories-lemma-initial} (rappelons que les
limites indexées par $I$ sont en réalité des limites sur la catégorie opposée
associée à $I$, de sorte que « cofinal » devient « initial »). Cela achève la
démonstration.
\end{proof}

\begin{lemma}
\label{lemma-relative-approximation}
Soit $f : X \to S$ un morphisme de schémas. Supposons que
\begin{enumerate}
\item $X$ soit quasi-compact et quasi-séparé, et
\item $S$ soit quasi-séparé.
\end{enumerate}
Alors $X = \lim X_i$ est la limite d'un système projectif filtrant de schémas $X_i$
de présentation finie sur $S$, dont les morphismes de transition sur $S$
sont affines.
\end{lemma}

\begin{proof}
Comme $f(X)$ est quasi-compact, nous pouvons remplacer $S$ par un ouvert
quasi-compact contenant $f(X)$. Nous pouvons donc supposer $S$
quasi-compact. D'après le Lemme \ref{lemma-approximate-morphism}, on peut
écrire $(X \to S) = \lim (X_i \to S_i)$ pour un certain système projectif
filtrant de morphismes entre schémas de type fini sur $\mathbf{Z}$, à
morphismes de transition affines. Comme les limites commutent aux limites
(Catégories, Lemme \ref{categories-lemma-colimits-commute}), on a
$X = \lim X_i \times_{S_i} S$. Soient $i \geq i'$ dans $I$. Le morphisme
$X_i \times_{S_i} S \to X_{i'} \times_{S_{i'}} S$ est affine, car c'est la
composée
$$
X_i \times_{S_i} S \to X_i \times_{S_{i'}} S \to X_{i'} \times_{S_{i'}} S
$$
où le premier morphisme est une immersion fermée (d'après Schémas,
Lemme \ref{schemes-lemma-fibre-product-after-map}) et le second est le
changement de base d'un morphisme affine (Morphismes,
Lemme \ref{morphisms-lemma-base-change-affine}) ; or la composée de morphismes
affines est affine (Morphismes,
Lemme \ref{morphisms-lemma-composition-affine}). Les morphismes $f_i$ sont de
présentation finie (Morphismes, Lemmes
\ref{morphisms-lemma-noetherian-finite-type-finite-presentation} et
\ref{morphisms-lemma-finite-presentation-permanence}) ; par conséquent,
leurs changements de base $X_i \times_{f_i, S_i} S \to S$ sont de
présentation finie (Morphismes,
Lemme \ref{morphisms-lemma-base-change-finite-presentation}).
\end{proof}

\begin{lemma}
\label{lemma-integral-limit-finite-and-finite-presentation}
Soit $X \to S$ un morphisme entier, où $S$ est quasi-compact et
quasi-séparé. Alors $X = \lim X_i$, avec $X_i \to S$ fini et de présentation
finie.
\end{lemma}

\begin{proof}
Considérons le faisceau $\mathcal{A} = f_*\mathcal{O}_X$. C'est un faisceau
quasi-cohérent d'$\mathcal{O}_S$-algèbres ; voir Schémas,
Lemme \ref{schemes-lemma-push-forward-quasi-coherent}. D'après
Propriétés,
Lemme \ref{properties-lemma-integral-algebra-directed-colimit-finite}, nous
pouvons écrire $\mathcal{A} = \colim_i \mathcal{A}_i$ comme limite inductive
filtrante d'$\mathcal{O}_S$-algèbres finies et de présentation finie. Alors
$$
X_i = \underline{\Spec}_S(\mathcal{A}_i)
\longrightarrow
S
$$
est un morphisme de schémas fini et de présentation finie. Par construction,
$X = \lim_i X_i$, ce qui démontre le lemme.
\end{proof}








\section{Descente des propriétés des morphismes}
\label{section-descent-of-properties}

\noindent
Cette section est l'analogue de la section \ref{section-descent} pour les
propriétés des morphismes sur $S$. Nous nous placerons dans la situation
suivante.

\begin{situation}
\label{situation-descent-property}
Soit $S = \lim S_i$ la limite d'un système projectif filtrant de schémas dont les
morphismes de transition sont affines
(Lemme \ref{lemma-directed-inverse-system-has-limit}). Soit $0 \in I$ et soit
$f_0 : X_0 \to Y_0$ un morphisme de schémas sur $S_0$. Supposons $S_0$,
$X_0$ et $Y_0$ quasi-compacts et quasi-séparés. Soit
$f_i : X_i \to Y_i$ le changement de base de $f_0$ à $S_i$, et soit
$f : X \to Y$ le changement de base de $f_0$ à $S$.
\end{situation}

\begin{lemma}
\label{lemma-descend-affine-finite-presentation}
Notations et hypothèses comme dans la
Situation \ref{situation-descent-property}. Si $f$ est affine, il existe un
indice $i \geq 0$ tel que $f_i$ soit affine.
\end{lemma}

\begin{proof}
Soit $Y_0 = \bigcup_{j = 1, \ldots, m} V_{j, 0}$ un recouvrement ouvert
affine fini. Posons $U_{j, 0} = f_0^{-1}(V_{j, 0})$. Pour $i \geq 0$, notons
$V_{j, i}$ l'image réciproque de $V_{j, 0}$ dans $Y_i$ et
$U_{j, i} = f_i^{-1}(V_{j, i})$. De même, on a $U_j = f^{-1}(V_j)$. Alors
$U_j = \lim_{i \geq 0} U_{j, i}$ (voir le
Lemme \ref{lemma-directed-inverse-system-has-limit}). Comme $U_j$ est affine
par hypothèse, le Lemme \ref{lemma-limit-affine} montre que chaque $U_{j, i}$
est affine pour $i$ assez grand. Puisqu'il n'y a qu'un nombre fini de $j$,
nous pouvons choisir un $i$ qui convienne à tous les $j$. Ainsi, $f_i$ est affine
pour $i$ assez grand ; voir Morphismes,
Lemme \ref{morphisms-lemma-characterize-affine}.
\end{proof}

\begin{lemma}
\label{lemma-descend-finite-finite-presentation}
Notations et hypothèses comme dans la
Situation \ref{situation-descent-property}. Si
\begin{enumerate}
\item $f$ est un morphisme fini, et
\item $f_0$ est localement de type fini,
\end{enumerate}
alors il existe $i \geq 0$ tel que $f_i$ soit fini.
\end{lemma}

\begin{proof}
Un morphisme fini est affine ; voir Morphismes,
Définition \ref{morphisms-definition-integral}. D'après le
Lemme \ref{lemma-descend-affine-finite-presentation} ci-dessus, nous pouvons,
après avoir augmenté $0$, supposer que $f_0$ est affine. En écrivant $Y_0$
comme réunion finie d'ouverts affines, on se ramène à démontrer le résultat
lorsque $X_0$ et $Y_0$ sont affines et s'envoient dans un même ouvert affine
$W \subset S_0$. L'énoncé algébrique correspondant résulte d'Algèbre,
Lemme \ref{algebra-lemma-colimit-finite}.
\end{proof}

\begin{lemma}
\label{lemma-descend-unramified}
Notations et hypothèses comme dans la
Situation \ref{situation-descent-property}. Si
\begin{enumerate}
\item $f$ est non ramifié, et
\item $f_0$ est localement de type fini,
\end{enumerate}
alors il existe $i \geq 0$ tel que $f_i$ soit non ramifié.
\end{lemma}

\begin{proof}
Choisissons un recouvrement ouvert affine fini
$Y_0 = \bigcup_{j = 1, \ldots, m} Y_{j, 0}$ tel que chaque $Y_{j, 0}$
s'envoie dans un ouvert affine $S_{j, 0} \subset S_0$. Pour chaque $j$, soit
$f_0^{-1}Y_{j, 0} = \bigcup_{k = 1, \ldots, n_j} X_{k, 0}$ un recouvrement
ouvert affine fini. Comme la propriété d'être non ramifié est locale, il
suffit de démontrer le lemme pour les morphismes entre schémas affines
$X_{k, i} \to Y_{j, i} \to S_{j, i}$ obtenus par changement de base
de $X_{k, 0} \to Y_{j, 0} \to S_{j, 0}$ à $S_i$. On se ramène ainsi au cas où
$X_0, Y_0, S_0$ sont affines.

\medskip\noindent
Dans le cas affine, on se ramène au résultat algébrique suivant. Supposons
$R = \colim_{i \in I} R_i$. Pour un certain $0 \in I$, supposons donné un
homomorphisme de type fini de $R_0$-algèbres $A_i \to B_i$. Si
$R \otimes_{R_0} A_0 \to R \otimes_{R_0} B_0$ est non ramifié, alors, pour
un certain $i \geq 0$, l'homomorphisme
$R_i \otimes_{R_0} A_0 \to R_i \otimes_{R_0} B_0$ est non ramifié. Cela
résulte d'Algèbre, Lemme \ref{algebra-lemma-colimit-unramified}.
\end{proof}

\begin{lemma}
\label{lemma-descend-closed-immersion-finite-presentation}
Notations et hypothèses comme dans la
Situation \ref{situation-descent-property}. Si
\begin{enumerate}
\item $f$ est une immersion fermée, et
\item $f_0$ est localement de type fini,
\end{enumerate}
alors il existe $i \geq 0$ tel que $f_i$ soit une immersion fermée.
\end{lemma}

\begin{proof}
Une immersion fermée est affine ; voir Morphismes,
Lemme \ref{morphisms-lemma-closed-immersion-affine}. D'après le
Lemme \ref{lemma-descend-affine-finite-presentation} ci-dessus, nous pouvons,
après avoir augmenté $0$, supposer que $f_0$ est affine. En écrivant $Y_0$
comme réunion finie d'ouverts affines, on se ramène à démontrer le résultat
lorsque $X_0$ et $Y_0$ sont affines et s'envoient dans un même ouvert affine
$W \subset S_0$. L'énoncé algébrique correspondant est une conséquence
d'Algèbre, Lemme \ref{algebra-lemma-colimit-surjective}.
\end{proof}

\begin{lemma}
\label{lemma-descend-separated-finite-presentation}
Notations et hypothèses comme dans la
Situation \ref{situation-descent-property}. Si $f$ est séparé, alors $f_i$
est séparé pour un certain $i \geq 0$.
\end{lemma}

\begin{proof}
Appliquons le Lemme \ref{lemma-descend-closed-immersion-finite-presentation}
au morphisme diagonal
$\Delta_{X_0/S_0} : X_0 \to X_0 \times_{S_0} X_0$. (C'est permis puisque
les morphismes diagonaux sont localement de type fini et que le produit
fibré $X_0 \times_{S_0} X_0$ est quasi-compact et quasi-séparé ; voir
Schémas, Lemme \ref{schemes-lemma-diagonal-immersion}, Morphismes,
Lemme \ref{morphisms-lemma-immersion-locally-finite-type}, et Schémas,
Remarque \ref{schemes-remark-quasi-compact-and-quasi-separated}.
\end{proof}

\begin{lemma}
\label{lemma-descend-flat-finite-presentation}
Notations et hypothèses comme dans la
Situation \ref{situation-descent-property}. Si
\begin{enumerate}
\item $f$ est plat,
\item $f_0$ est localement de présentation finie,
\end{enumerate}
alors $f_i$ est plat pour un certain $i \geq 0$.
\end{lemma}

\begin{proof}
Choisissons un recouvrement ouvert affine fini
$Y_0 = \bigcup_{j = 1, \ldots, m} Y_{j, 0}$ tel que chaque $Y_{j, 0}$
s'envoie dans un ouvert affine $S_{j, 0} \subset S_0$. Pour chaque $j$, soit
$f_0^{-1}Y_{j, 0} = \bigcup_{k = 1, \ldots, n_j} X_{k, 0}$ un recouvrement
ouvert affine fini. Comme la propriété d'être plat est locale, il suffit de
démontrer le lemme pour les morphismes entre schémas affines
$X_{k, i} \to Y_{j, i} \to S_{j, i}$ obtenus par changement de base
de $X_{k, 0} \to Y_{j, 0} \to S_{j, 0}$ à $S_i$. On se ramène ainsi au cas où
$X_0, Y_0, S_0$ sont affines.

\medskip\noindent
Dans le cas affine, on se ramène au résultat algébrique suivant. Supposons
$R = \colim_{i \in I} R_i$. Pour un certain $0 \in I$, supposons donné un
homomorphisme de présentation finie de $R_0$-algèbres $A_i \to B_i$. Si
$R \otimes_{R_0} A_0 \to R \otimes_{R_0} B_0$ est plat, alors, pour un
certain $i \geq 0$, l'homomorphisme
$R_i \otimes_{R_0} A_0 \to R_i \otimes_{R_0} B_0$ est plat. Cela résulte
d'Algèbre, Lemme \ref{algebra-lemma-flat-finite-presentation-limit-flat},
assertion (3).
\end{proof}

\begin{lemma}
\label{lemma-descend-finite-locally-free}
Notations et hypothèses comme dans la
Situation \ref{situation-descent-property}. Si
\begin{enumerate}
\item $f$ est fini et localement libre (de degré $d$),
\item $f_0$ est localement de présentation finie,
\end{enumerate}
alors $f_i$ est fini et localement libre (de degré $d$) pour un certain
$i \geq 0$.
\end{lemma}

\begin{proof}
D'après les Lemmes \ref{lemma-descend-flat-finite-presentation} et
\ref{lemma-descend-finite-finite-presentation}, il existe $i$ tel que $f_i$
soit plat et fini. D'autre part, $f_i$ est localement de présentation finie.
Ainsi, $f_i$ est fini et localement libre d'après Morphismes,
Lemme \ref{morphisms-lemma-finite-flat}. Si, de plus, $f$ est fini et
localement libre de degré $d$, alors l'image de $Y \to Y_i$ est contenue dans
le lieu ouvert et fermé $W_d \subset Y_i$ au-dessus duquel $f_i$ est de
degré $d$. D'après le Lemme \ref{lemma-limit-contained-in-constructible},
pour un certain $i' \geq i$, l'image de $Y_{i'} \to Y_i$ est contenue dans $W_d$.
Alors $f_{i'}$ est fini et localement libre de degré $d$.
\end{proof}

\begin{lemma}
\label{lemma-descend-smooth}
Notations et hypothèses comme dans la
Situation \ref{situation-descent-property}. Si
\begin{enumerate}
\item $f$ est lisse,
\item $f_0$ est localement de présentation finie,
\end{enumerate}
alors $f_i$ est lisse pour un certain $i \geq 0$.
\end{lemma}

\begin{proof}
La lissité est locale sur la source et sur le but (Morphismes,
Lemme \ref{morphisms-lemma-smooth-characterize}) ; nous pouvons donc supposer
$S_0, X_0, Y_0$ affines (détails omis). Le fait algébrique correspondant
est Algèbre, Lemme \ref{algebra-lemma-colimit-smooth}.
\end{proof}

\begin{lemma}
\label{lemma-descend-etale}
Notations et hypothèses comme dans la
Situation \ref{situation-descent-property}. Si
\begin{enumerate}
\item $f$ est étale,
\item $f_0$ est localement de présentation finie,
\end{enumerate}
alors $f_i$ est étale pour un certain $i \geq 0$.
\end{lemma}

\begin{proof}
La propriété d'être étale est locale sur la source et sur le but (Morphismes,
Lemme \ref{morphisms-lemma-etale-characterize}) ; nous pouvons donc supposer
$S_0, X_0, Y_0$ affines (détails omis). Le fait algébrique correspondant
est Algèbre, Lemme \ref{algebra-lemma-colimit-etale}.
\end{proof}

\begin{lemma}
\label{lemma-descend-isomorphism}
Notations et hypothèses comme dans la
Situation \ref{situation-descent-property}. Si
\begin{enumerate}
\item $f$ est un isomorphisme, et
\item $f_0$ est localement de présentation finie,
\end{enumerate}
alors $f_i$ est un isomorphisme pour un certain $i \geq 0$.
\end{lemma}

\begin{proof}
D'après les Lemmes \ref{lemma-descend-etale} et
\ref{lemma-descend-closed-immersion-finite-presentation}, nous pouvons
trouver $i$ tel que $f_i$ soit plat et soit une immersion fermée. Alors $f_i$
identifie $X_i$ à un sous-schéma ouvert et fermé de $Y_i$ ; voir Morphismes,
Lemme \ref{morphisms-lemma-flat-closed-immersions-finite-presentation}. Par
hypothèse, l'image de $Y \to Y_i$ s'envoie dans $f_i(X_i)$. Le
Lemme \ref{lemma-limit-contained-in-constructible} montre donc que $Y_{i'}$
s'envoie dans $f_i(X_i)$ pour un certain $i' \geq i$. Il en résulte que
$X_{i'} \to Y_{i'}$ est surjectif, et on conclut.
\end{proof}

\begin{lemma}
\label{lemma-descend-open-immersion}
Notations et hypothèses comme dans la
Situation \ref{situation-descent-property}. Si
\begin{enumerate}
\item $f$ est une immersion ouverte, et
\item $f_0$ est localement de présentation finie,
\end{enumerate}
alors $f_i$ est une immersion ouverte pour un certain $i \geq 0$.
\end{lemma}

\begin{proof}
D'après le Lemme \ref{lemma-descend-etale}, nous pouvons trouver $i$ tel que
$f_i$ soit étale. Alors $V_i = f_i(X_i)$ est un sous-schéma ouvert
quasi-compact de $Y_i$ (Morphismes,
Lemme \ref{morphisms-lemma-etale-open}). Soient $V$ et $V_{i'}$, pour
$i' \geq i$, les images réciproques de $V_i$ dans $Y$ et $Y_{i'}$. Alors
$f : X \to V$ est un isomorphisme (c'est en effet une immersion ouverte
surjective). D'après le Lemme \ref{lemma-descend-isomorphism}, le morphisme
$X_{i'} \to V_{i'}$ est donc un isomorphisme pour un certain $i' \geq i$,
comme voulu.
\end{proof}

\begin{lemma}
\label{lemma-descend-immersion}
Notations et hypothèses comme dans la
Situation \ref{situation-descent-property}. Si
\begin{enumerate}
\item $f$ est une immersion, et
\item $f_0$ est localement de type fini,
\end{enumerate}
alors $f_i$ est une immersion pour un certain $i \geq 0$.
\end{lemma}

\begin{proof}
Il existe un ouvert $V \subset Y$ tel que le morphisme $f$ se factorise en
$X \to V \to Y$ et que $X \to V$ soit une immersion fermée ; voir la
discussion de Schémas, section \ref{schemes-section-immersions}. Puisque $X$
est quasi-compact, nous pouvons et allons supposer que $V$ est un ouvert
quasi-compact de $Y$. Après avoir augmenté $0$, le
Lemme \ref{lemma-descend-opens} fournit un ouvert quasi-compact
$V_0 \subset Y_0$ tel que $V$ soit l'image réciproque de $V_0$. L'image réciproque de
$V_0$ dans $X_0$ est alors un ouvert quasi-compact dont l'image réciproque
dans $X$ est $X$. En appliquant le même lemme à $X = \lim X_i$, et après
avoir augmenté $0$, nous pouvons donc supposer donnée la factorisation
$X_0 \to V_0 \to Y_0$. Pour $i \geq 0$ assez grand, le morphisme
$X_i \to V_i$, où $V_i = Y_i \times_{Y_0} V_0$, est alors une immersion
fermée d'après le
Lemme \ref{lemma-descend-closed-immersion-finite-presentation}. La
démonstration est achevée.
\end{proof}

\begin{lemma}
\label{lemma-descend-monomorphism}
Notations et hypothèses comme dans la
Situation \ref{situation-descent-property}. Si
\begin{enumerate}
\item $f$ est un monomorphisme, et
\item $f_0$ est localement de type fini,
\end{enumerate}
alors $f_i$ est un monomorphisme pour un certain $i \geq 0$.
\end{lemma}

\begin{proof}
Rappelons qu'un morphisme de schémas $V \to W$ est un monomorphisme si et
seulement si la diagonale $V \to V \times_W V$ est un isomorphisme (Schémas,
Lemme \ref{schemes-lemma-monomorphism}). Le morphisme
$X_0 \to X_0 \times_{Y_0} X_0$ est localement de présentation finie d'après
Morphismes, Lemme \ref{morphisms-lemma-diagonal-morphism-finite-type}.
Puisque $X_0 \times_{Y_0} X_0$ est quasi-compact et quasi-séparé (Schémas,
Remarque \ref{schemes-remark-quasi-compact-and-quasi-separated}), le
Lemme \ref{lemma-descend-isomorphism} montre que
$\Delta_i : X_i \to X_i \times_{Y_i} X_i$ est un isomorphisme pour un
certain $i \geq 0$. Pour cet $i$, le morphisme $f_i$ est un monomorphisme.
\end{proof}

\begin{lemma}
\label{lemma-descend-surjective}
Notations et hypothèses comme dans la
Situation \ref{situation-descent-property}. Si
\begin{enumerate}
\item $f$ est surjectif, et
\item $f_0$ est localement de présentation finie,
\end{enumerate}
alors il existe $i \geq 0$ tel que $f_i$ soit surjectif.
\end{lemma}

\begin{proof}
Le morphisme $f_0$ est de présentation finie. Ainsi,
$E = f_0(X_0)$ est une partie constructible de $Y_0$ ; voir Morphismes,
Lemme \ref{morphisms-lemma-chevalley}. Comme $f_i$ est le changement de base
de $f_0$ par $Y_i \to Y_0$, l'image de $f_i$ est l'image réciproque de $E$
dans $Y_i$. De plus, nous savons que $Y \to Y_0$ s'envoie dans $E$. On
conclut donc par le Lemme \ref{lemma-limit-contained-in-constructible}.
\end{proof}

\begin{lemma}
\label{lemma-descend-syntomic}
Notations et hypothèses comme dans la
Situation \ref{situation-descent-property}. Si
\begin{enumerate}
\item $f$ est syntomique, et
\item $f_0$ est localement de présentation finie,
\end{enumerate}
alors il existe $i \geq 0$ tel que $f_i$ soit syntomique.
\end{lemma}

\begin{proof}
Choisissons un recouvrement ouvert affine fini
$Y_0 = \bigcup_{j = 1, \ldots, m} Y_{j, 0}$ tel que chaque $Y_{j, 0}$
s'envoie dans un ouvert affine $S_{j, 0} \subset S_0$. Pour chaque $j$, soit
$f_0^{-1}Y_{j, 0} = \bigcup_{k = 1, \ldots, n_j} X_{k, 0}$ un recouvrement
ouvert affine fini. Comme la propriété d'être syntomique est locale, il
suffit de démontrer le lemme pour les morphismes entre schémas affines
$X_{k, i} \to Y_{j, i} \to S_{j, i}$ obtenus par changement de base
de $X_{k, 0} \to Y_{j, 0} \to S_{j, 0}$ à $S_i$. On se ramène ainsi au cas où
$X_0, Y_0, S_0$ sont affines.

\medskip\noindent
Dans le cas affine, on se ramène au résultat algébrique suivant. Supposons
$R = \colim_{i \in I} R_i$. Pour un certain $0 \in I$, supposons donné un
homomorphisme de présentation finie de $R_0$-algèbres $A_i \to B_i$. Si
$R \otimes_{R_0} A_0 \to R \otimes_{R_0} B_0$ est syntomique, alors, pour
un certain $i \geq 0$, l'homomorphisme
$R_i \otimes_{R_0} A_0 \to R_i \otimes_{R_0} B_0$ est syntomique. Cela
résulte d'Algèbre, Lemme \ref{algebra-lemma-colimit-lci}.
\end{proof}





\section{Immersion fermée d'un morphisme de type fini dans un morphisme de présentation finie}
\label{section-finite-type-closed-in-finite-presentation}

\noindent
Un résultat de ce type est \cite[Satz 2.10]{Kiehl}. Une autre référence est
\cite{Conrad-Nagata}.

\begin{lemma}
\label{lemma-locally-finite-type-in-finite-presentation}
Soit $f : X \to S$ un morphisme de schémas. Faisons les hypothèses suivantes :
\begin{enumerate}
\item le morphisme $f$ est localement de type fini ;
\item le schéma $X$ est quasi-compact et quasi-séparé.
\end{enumerate}
Alors il existe un morphisme de présentation finie $f' : X' \to S$ et une
immersion $X \to X'$ de schémas sur $S$.
\end{lemma}

\begin{proof}
D'après la Proposition \ref{proposition-approximate}, on peut écrire
$X = \lim_i X_i$, chaque $X_i$ étant de type fini sur $\mathbf{Z}$ et les
morphismes de transition $f_{ii'} : X_i \to X_{i'}$ étant affines.
Considérons le diagramme commutatif
$$
\xymatrix{
X \ar[r] \ar[rd] & X_{i, S} \ar[r] \ar[d] & X_i \ar[d] \\
& S \ar[r] & \Spec(\mathbf{Z})
}
$$
Remarquons que $X_i$ est de présentation finie sur $\Spec(\mathbf{Z})$ ;
voir Morphismes,
Lemme \ref{morphisms-lemma-noetherian-finite-type-finite-presentation}.
Ainsi, le changement de base $X_{i, S} \to S$ est de présentation finie par
Morphismes, Lemme \ref{morphisms-lemma-base-change-finite-presentation}. Il
suffit donc de montrer que la flèche $X \to X_{i, S}$ est une immersion pour
$i$ assez grand.

\medskip\noindent
Pour cela, choisissons un recouvrement ouvert affine fini
$X = V_1 \cup \ldots \cup V_n$ tel que $f$ envoie chaque $V_j$ dans un ouvert
affine $U_j \subset S$. Soit $h_{j, a} \in \mathcal{O}_X(V_j)$ une famille
finie d'éléments qui engendrent $\mathcal{O}_X(V_j)$ comme
$\mathcal{O}_S(U_j)$-algèbre ; voir Morphismes,
Lemme \ref{morphisms-lemma-locally-finite-type-characterize}. D'après les
Lemmes \ref{lemma-descend-opens} et \ref{lemma-limit-affine}, et après avoir
éventuellement restreint $I$, nous pouvons supposer qu'il existe des
recouvrements ouverts affines $X_i = V_{1, i} \cup \ldots \cup V_{n, i}$,
compatibles aux morphismes de transition et tels que
$V_j = \lim_i V_{j, i}$. D'après le Lemme \ref{lemma-descend-section}, nous
pouvons choisir $i$ assez grand pour que chaque $h_{j, a}$ provienne d'un
élément $h_{j, a, i} \in \mathcal{O}_{X_i}(V_{j, i})$. Ainsi, la flèche dans
$$
V_j \longrightarrow U_j \times_{\Spec(\mathbf{Z})} V_{j, i} =
(V_{j, i})_{U_j} \subset (V_{j, i})_S \subset X_{i, S}
$$
est une immersion fermée. Puisque $\bigcup (V_{j, i})_{U_j}$ est un ouvert
de $X_{i, S}$ et que l'image réciproque de $(V_{j, i})_{U_j}$ dans $X$ est
$V_j$, il en résulte que $X \to X_{i, S}$ est une immersion.
\end{proof}

\begin{remark}
\label{remark-cannot-do-better}
On ne peut améliorer cet énoncé sans hypothèses supplémentaires sur $S$ et
sur le morphisme $f : X \to S$. Par exemple, il ne sera généralement pas
possible de trouver comme dans le lemme une immersion {\it fermée}
$X \to X'$. En effet, cela impliquerait que $f$ soit quasi-compact, ce qui
peut ne pas être le cas. On peut prendre pour exemple $S$ égal à l'espace
affine de dimension infinie dont le point $0$ est dédoublé, et $X$ égal à
l'un des deux espaces affines de dimension infinie.
\end{remark}

\begin{lemma}
\label{lemma-finite-type-closed-in-finite-presentation}
Soit $f : X \to S$ un morphisme de schémas. Faisons les hypothèses suivantes :
\begin{enumerate}
\item le morphisme $f$ est localement de type fini ;
\item le schéma $X$ est quasi-compact et quasi-séparé ;
\item le schéma $S$ est quasi-séparé.
\end{enumerate}
Alors il existe un morphisme de présentation finie $f' : X' \to S$ et une
immersion fermée $X \to X'$ de schémas sur $S$.
\end{lemma}

\begin{proof}
D'après le Lemme \ref{lemma-locally-finite-type-in-finite-presentation}
ci-dessus, il existe un morphisme de présentation finie $Y \to S$ et une
immersion $i : X \to Y$ de schémas sur $S$. Pour tout point $x \in X$, il
existe un ouvert affine $V_x \subset Y$ tel que
$i^{-1}(V_x) \to V_x$ soit une immersion fermée. Puisque $X$ est
quasi-compact, nous pouvons trouver un nombre fini d'ouverts affines
$V_1, \ldots, V_n \subset Y$ tels que
$i(X) \subset V_1 \cup \ldots \cup V_n$ et que
$i^{-1}(V_j) \to V_j$ soit une immersion fermée. Autrement dit,
$i : X \to X' = V_1 \cup \ldots \cup V_n$ est une immersion fermée de
schémas sur $S$. Comme $S$ est quasi-séparé et $Y$ quasi-séparé sur $S$, on
en déduit que $Y$ est quasi-séparé ; voir Schémas,
Lemme \ref{schemes-lemma-separated-permanence}. Ainsi, l'immersion ouverte
$X' = V_1 \cup \ldots \cup V_n \to Y$ est quasi-compacte. Il en résulte que
$X' \to Y$ est de présentation finie ; voir Morphismes,
Lemme \ref{morphisms-lemma-quasi-compact-open-immersion-finite-presentation}.
On conclut, car $X' \to Y \to S$ est alors une composée de morphismes de
présentation finie, donc est de présentation finie (voir Morphismes,
Lemme \ref{morphisms-lemma-composition-finite-presentation}).
\end{proof}

\begin{lemma}
\label{lemma-closed-is-limit-closed-and-finite-presentation}
Soit $X \to Y$ une immersion fermée de schémas. Supposons $Y$ quasi-compact
et quasi-séparé. Alors $X$ peut s'écrire comme une limite projective
$X = \lim X_i$ de schémas sur $Y$, où $X_i \to Y$ est une immersion fermée
de présentation finie.
\end{lemma}

\begin{proof}
Soit $\mathcal{I} \subset \mathcal{O}_Y$ le faisceau quasi-cohérent d'idéaux
qui définit $X$ comme sous-schéma fermé de $Y$. D'après Propriétés,
Lemme \ref{properties-lemma-quasi-coherent-colimit-finite-type}, nous pouvons
écrire $\mathcal{I}$ comme limite inductive filtrante
$\mathcal{I} = \colim_{i \in I} \mathcal{I}_i$ de ses faisceaux
quasi-cohérents d'idéaux de type fini. Soit $X_i \subset Y$ le sous-schéma
fermé défini par $\mathcal{I}_i$. Ils forment un système projectif de schémas
indexé par $I$. Les morphismes de transition $X_i \to X_{i'}$ sont affines
parce que ce sont des immersions fermées. Chaque $X_i$ est quasi-compact et
quasi-séparé, puisqu'il est un sous-schéma fermé de $Y$ et que $Y$ est
quasi-compact et quasi-séparé par hypothèse. On a $X = \lim_i X_i$, comme il
résulte directement du fait que
$\mathcal{I} = \colim_{i \in I} \mathcal{I}_a$. Chacun des morphismes
$X_i \to Y$ est de présentation finie ; voir Morphismes,
Lemme \ref{morphisms-lemma-closed-immersion-finite-presentation}.
\end{proof}

\begin{lemma}
\label{lemma-finite-type-is-limit-finite-presentation}
Soit $f : X \to S$ un morphisme de schémas. Faisons les hypothèses suivantes :
\begin{enumerate}
\item le morphisme $f$ est localement de type fini ;
\item le schéma $X$ est quasi-compact et quasi-séparé ;
\item le schéma $S$ est quasi-séparé.
\end{enumerate}
Alors $X = \lim X_i$, où les $X_i \to S$ sont de présentation finie, les
$X_i$ sont quasi-compacts et quasi-séparés, et les morphismes de transition
$X_{i'} \to X_i$ sont des immersions fermées (ce qui implique que
$X \to X_i$ est une immersion fermée pour tout $i$).
\end{lemma}

\begin{proof}
D'après le Lemme \ref{lemma-finite-type-closed-in-finite-presentation}, il
existe une immersion fermée $X \to Y$, où $Y \to S$ est de présentation
finie. Le schéma $Y$ est alors quasi-séparé par Schémas,
Lemme \ref{schemes-lemma-separated-permanence}. Comme $X$ est quasi-compact,
nous pouvons supposer $Y$ quasi-compact en remplaçant $Y$ par un ouvert
quasi-compact contenant $X$. On a $X = \lim X_i$, avec
$X_i \to Y$ immersion fermée de présentation finie, d'après le
Lemme \ref{lemma-closed-is-limit-closed-and-finite-presentation}.
Les morphismes $X_i \to S$ sont de présentation finie par Morphismes,
Lemme \ref{morphisms-lemma-composition-finite-presentation}.
\end{proof}

\begin{proposition}
\label{proposition-separated-closed-in-finite-presentation}
Soit $f : X \to S$ un morphisme de schémas. Faisons les hypothèses suivantes :
\begin{enumerate}
\item $f$ est de type fini et séparé, et
\item $S$ est quasi-compact et quasi-séparé.
\end{enumerate}
Alors il existe un morphisme séparé de présentation finie $f' : X' \to S$
et une immersion fermée $X \to X'$ de schémas sur $S$.
\end{proposition}

\begin{proof}
Appliquons le Lemme \ref{lemma-finite-type-is-limit-finite-presentation} et
remarquons que $X_i \to S$ est séparé pour $i$ assez grand d'après le
Lemme \ref{lemma-eventually-separated}, puisque nous avons supposé
$X \to S$ séparé.
\end{proof}

\begin{lemma}
\label{lemma-finite-closed-in-finite-finite-presentation}
Soit $f : X \to S$ un morphisme de schémas. Faisons les hypothèses suivantes :
\begin{enumerate}
\item $f$ est fini, et
\item $S$ est quasi-compact et quasi-séparé.
\end{enumerate}
Alors il existe un morphisme $f' : X' \to S$ fini et de présentation finie,
ainsi qu'une immersion fermée $X \to X'$ de schémas sur $S$.
\end{lemma}

\begin{proof}
Nous pouvons écrire $X = \lim X_i$ comme dans le
Lemme \ref{lemma-finite-type-is-limit-finite-presentation}. En appliquant le
Lemme \ref{lemma-eventually-finite}, nous voyons que $X_i \to S$ est fini
pour $i$ assez grand.
\end{proof}

\begin{lemma}
\label{lemma-finite-in-finite-and-finite-presentation}
Soit $f : X \to S$ un morphisme de schémas. Faisons les hypothèses suivantes :
\begin{enumerate}
\item $f$ est fini, et
\item $S$ est quasi-compact et quasi-séparé.
\end{enumerate}
Alors $X$ est une limite projective $X = \lim X_i$, dont les morphismes de
transition sont des immersions fermées et dont les objets $X_i$ sont finis
et de présentation finie sur $S$.
\end{lemma}

\begin{proof}
Nous pouvons écrire $X = \lim X_i$ comme dans le
Lemme \ref{lemma-finite-type-is-limit-finite-presentation}. En appliquant le
Lemme \ref{lemma-eventually-finite}, nous voyons que $X_i \to S$ est fini
pour $i$ assez grand.
\end{proof}








\section{Descente des objets relatifs}
\label{section-descending-relative}

\noindent
Le lemme suivant est représentatif des résultats de cette section.
Nous en donnons intégralement la démonstration « standard ». Il sera
peut-être plus rapide de se convaincre du résultat que de lire cette démonstration.

\begin{lemma}
\label{lemma-descend-finite-presentation}
Soit $I$ un ensemble ordonné filtrant.
Soit $(S_i, f_{ii'})$ un système projectif de schémas indexé par $I$.
Supposons que
\begin{enumerate}
\item les morphismes $f_{ii'} : S_i \to S_{i'}$ soient affines,
\item les schémas $S_i$ soient quasi-compacts et quasi-séparés.
\end{enumerate}
Soit $S = \lim_i S_i$. On a alors les propriétés suivantes :
\begin{enumerate}
\item Pour tout morphisme de présentation finie $X \to S$,
il existe un indice $i \in I$ et un morphisme de présentation finie
$X_i \to S_i$ tels que $X \cong X_{i, S}$ comme schémas sur $S$.
\item Étant donnés un indice $i \in I$, des schémas
$X_i$, $Y_i$ de présentation finie sur $S_i$ et un morphisme
$\varphi : X_{i, S} \to Y_{i, S}$ sur $S$, il existe un indice
$i' \geq i$ et un morphisme
$\varphi_{i'} : X_{i, S_{i'}} \to Y_{i, S_{i'}}$
dont le changement de base à $S$ est $\varphi$.
\item Étant donnés un indice $i \in I$, des schémas $X_i$, $Y_i$
de présentation finie sur $S_i$ et deux morphismes
$\varphi_i, \psi_i : X_i \to Y_i$
dont les changements de base $\varphi_{i, S} = \psi_{i, S}$ sont égaux,
il existe un indice $i' \geq i$ tel que
$\varphi_{i, S_{i'}} = \psi_{i, S_{i'}}$.
\end{enumerate}
Autrement dit, la catégorie des schémas de présentation finie sur
$S$ est la colimite, indexée par $I$, des catégories de schémas de
présentation finie sur $S_i$.
\end{lemma}

\begin{proof}
Si tous les schémas $S_i$ sont affines et si l'on ne considère
que les schémas affines de présentation finie sur $S_i$, resp.\ $S$,
ce lemme équivaut à
Algèbre, Lemme \ref{algebra-lemma-colimit-category-fp-algebras}.
Nous affirmons que le cas affine entraîne le cas général.

\medskip\noindent
Démontrons (3). Donnons-nous un indice $i \in I$, des schémas
$X_i$, $Y_i$ de présentation finie sur $S_i$ et deux morphismes
$\varphi_i, \psi_i : X_i \to Y_i$. Supposons leurs changements de base
égaux : $\varphi_{i, S} = \psi_{i, S}$. Nous noterons
$X_{i'} = X_{i, S_{i'}}$ et $Y_{i'} = Y_{i, S_{i'}}$ pour
$i' \geq i$. Posons aussi $X = X_{i, S}$ et $Y = Y_{i, S}$.
D'après le Lemme \ref{lemma-scheme-over-limit}, on a
$X = \lim_{i' \geq i} X_{i'}$, et de même pour $Y$.
Notons en outre $\varphi_{i'}$ et $\psi_{i'}$
(resp.\ $\varphi$ et $\psi$)
les changements de base de $\varphi_i$ et $\psi_i$ à $S_{i'}$
(resp.\ $S$). Notre hypothèse signifie donc que $\varphi = \psi$.
Puisque $Y_i$ et $X_i$ sont de présentation finie
sur $S_i$ et que $S_i$ est quasi-compact et quasi-séparé,
$X_i$ et $Y_i$ sont eux aussi quasi-compacts et quasi-séparés
(voir Morphismes,
Lemme \ref{morphisms-lemma-finite-presentation-quasi-compact-quasi-separated}).
On peut donc choisir un recouvrement ouvert affine fini
$Y_i = \bigcup V_{j, i}$ tel que chaque $V_{j, i}$ ait son image dans
un ouvert affine de $S_i$. Comme ci-dessus, notons $V_{j, i'}$ l'image
réciproque de $V_{j, i}$ dans $Y_{i'}$ et $V_j$ son image réciproque dans $Y$.
Les immersions $V_{j, i'} \to Y_{i'}$ sont quasi-compactes, et les images réciproques
$U_{j, i'} = \varphi_i^{-1}(V_{j, i'})$ et
$U_{j, i'}' = \psi_i^{-1}(V_{j, i'})$
sont des ouverts quasi-compacts de $X_{i'}$. Par hypothèse, les images réciproques
de $V_j$ par $\varphi$ et $\psi$ dans $X$ sont égales.
D'après le Lemme \ref{lemma-descend-opens},
il existe donc un indice $i' \geq i$ tel que
$U_{j, i'} = U_{j, i'}'$ dans $X_{i'}$.
Choisissons un recouvrement ouvert affine fini
$U_{j, i'} = U_{j, i'}' = \bigcup W_{j, k, i'}$
qui induit des recouvrements $U_{j, i''} = U_{j, i''}' = \bigcup W_{j, k, i''}$
pour tout $i'' \geq i'$.
D'après le cas affine, il existe
un indice $i''$ tel que
$\varphi_{i''}|_{W_{j, k, i''}} = \psi_{i''}|_{W_{j, k, i''}}$
pour tous $j, k$. Alors $i''$ vérifie
$\varphi_{i''} = \psi_{i''}$, ce qui démontre (3).

\medskip\noindent
Démontrons (2). Donnons-nous un indice $i \in I$, des schémas
$X_i$, $Y_i$ de présentation finie sur $S_i$ et un morphisme
$\varphi : X_{i, S} \to Y_{i, S}$. Nous noterons
$X_{i'} = X_{i, S_{i'}}$ et $Y_{i'} = Y_{i, S_{i'}}$ pour
$i' \geq i$. Posons aussi $X = X_{i, S}$ et $Y = Y_{i, S}$.
D'après le Lemme \ref{lemma-scheme-over-limit}, on a
$X = \lim_{i' \geq i} X_{i'}$, et de même pour $Y$.
Puisque $Y_i$ et $X_i$ sont de présentation finie
sur $S_i$ et que $S_i$ est quasi-compact et quasi-séparé,
$X_i$ et $Y_i$ sont eux aussi quasi-compacts et quasi-séparés
(voir Morphismes,
Lemme \ref{morphisms-lemma-finite-presentation-quasi-compact-quasi-separated}).
On peut donc choisir un recouvrement ouvert affine fini
$Y_i = \bigcup V_{j, i}$ tel que chaque $V_{j, i}$ ait son image dans
un ouvert affine de $S_i$. Comme ci-dessus, notons $V_{j, i'}$ l'image
réciproque de $V_{j, i}$ dans $Y_{i'}$ et $V_j$ son image réciproque dans $Y$.
Les immersions $V_j \to Y$ sont quasi-compactes, et les images réciproques
$U_j = \varphi^{-1}(V_j)$ sont des ouverts quasi-compacts de $X$.
D'après le Lemme \ref{lemma-descend-opens}, il existe donc un indice
$i' \geq i$ et des ouverts quasi-compacts $U_{j, i'}$ de $X_{i'}$
dont l'image réciproque dans $X$ est $U_j$. Choisissons un recouvrement ouvert affine fini
$U_{j, i'} = \bigcup W_{j, k, i'}$ qui induit des recouvrements ouverts affines
$U_{j, i''} = \bigcup W_{j, k, i''}$
pour tout $i'' \geq i'$ et un recouvrement ouvert affine
$U_j = \bigcup W_{j, k}$. D'après le cas affine, il existe
un indice $i''$ et des morphismes
$\varphi_{j, k, i''} : W_{j, k, i''} \to V_{j, i''}$
tels que
$\varphi|_{W_{j, k}} = \varphi_{j, k, i'', S}$ pour tous $j, k$.
D'après (3), démontré ci-dessus, il existe un indice $i''' \geq i''$
tel que
$$
\varphi_{j_1, k_1, i'', S_{i'''}}|_{W_{j_1, k_1, i'''} \cap W_{j_2, k_2, i'''}}
=
\varphi_{j_2, k_2, i'', S_{i'''}}|_{W_{j_1, k_1, i'''} \cap W_{j_2, k_2, i'''}}
$$
pour tous $j_1, j_2, k_1, k_2$. Alors $i'''$ est un indice pour lequel
il existe un morphisme $\varphi_{i'''} : X_{i'''} \to Y_{i'''}$
dont le changement de base à $S$ est $\varphi$. Ainsi (2) est établi.

\medskip\noindent
Démontrons (1). Donnons-nous un schéma $X$ de présentation finie
sur $S$. Puisque $X$ est de présentation finie
sur $S$ et que $S$ est quasi-compact et quasi-séparé,
$X$ est lui aussi quasi-compact et quasi-séparé
(voir Morphismes,
Lemme \ref{morphisms-lemma-finite-presentation-quasi-compact-quasi-separated}).
Choisissons un recouvrement ouvert affine fini $X = \bigcup U_j$
tel que chaque $U_j$ ait son image dans un ouvert affine $V_j \subset S$.
Notons $U_{j_1j_2} = U_{j_1} \cap U_{j_2}$ et
$U_{j_1j_2j_3} = U_{j_1} \cap U_{j_2} \cap U_{j_3}$.
D'après les Lemmes \ref{lemma-descend-opens} et \ref{lemma-limit-affine},
on peut trouver un indice $i_1$ et des ouverts affines $V_{j, i_1} \subset S_{i_1}$
tels que chaque $V_j$ soit leur image réciproque dans $S$.
Soit $V_{j, i}$ l'image réciproque de $V_{j, i_1}$ dans $S_i$ pour
$i \geq i_1$. D'après le cas affine, on peut trouver un indice $i_2 \geq i_1$ et
des schémas affines $U_{j, i_2} \to V_{j, i_2}$ tels
que $U_j = S \times_{S_{i_2}} U_{j, i_2}$ soit le changement de base.
Notons $U_{j, i} = S_i \times_{S_{i_2}} U_{j, i_2}$ pour $i \geq i_2$.
D'après le Lemme \ref{lemma-descend-opens}, il existe un indice
$i_3 \geq i_2$ et des sous-schémas ouverts
$W_{j_1, j_2, i_3} \subset U_{j_1, i_3}$
dont le changement de base à $S$ est égal à $U_{j_1j_2}$.
Notons $W_{j_1, j_2, i} = S_i \times_{S_{i_3}} W_{j_1, j_2, i_3}$
pour $i \geq i_3$. D'après (2), démontré ci-dessus, il existe un indice
$i_4 \geq i_3$ et des morphismes
$\varphi_{j_1, j_2, i_4} : W_{j_1, j_2, i_4} \to W_{j_2, j_1, i_4}$
dont le changement de base à $S$ est le morphisme identité
$U_{j_1j_2} = U_{j_2j_1}$ pour tous $j_1, j_2$.
Pour tout $i \geq i_4$, notons
$\varphi_{j_1, j_2, i} = \text{id}_S \times \varphi_{j_1, j_2, i_4}$
le changement de base. Nous affirmons que, pour un certain $i_5 \geq i_4$, le système
$((U_{j, i_5})_j, (W_{j_1, j_2, i_5})_{j_1, j_2},
(\varphi_{j_1, j_2, i_5})_{j_1, j_2})$ est une donnée de recollement
au sens de Schémas, section \ref{schemes-section-glueing-schemes}.
Pour le voir, il faut vérifier que, pour $i$ assez grand,
on a
$$
\varphi_{j_1, j_2, i}^{-1}(W_{j_2, j_1, i} \cap W_{j_2, j_3, i})
=
W_{j_1, j_2, i} \cap W_{j_1, j_3, i}
$$
et que, pour $i$ assez grand, la condition de cocycle est satisfaite.
La première condition résulte du Lemme \ref{lemma-descend-opens}
et de l'égalité $U_{j_2j_1j_3} = U_{j_1j_2j_3}$.
La seconde résulte de (3), démontré ci-dessus, et du fait
que la condition de cocycle est satisfaite par les applications
$\text{id} : U_{j_1j_2} \to U_{j_2j_1}$.
On peut donc appliquer Schémas, Lemme \ref{schemes-lemma-glue-schemes}
pour recoller le système
$((U_{j, i_5})_j, (W_{j_1, j_2, i_5})_{j_1, j_2},
(\varphi_{j_1, j_2, i_5})_{j_1, j_2})$ et obtenir un schéma
$X_{i_5} \to S_{i_5}$. Par construction, le changement de base de
$X_{i_5}$ à $S$ est obtenu en recollant les ouverts affines
$U_j$ le long des ouverts $U_{j_1} \leftarrow U_{j_1j_2} \rightarrow U_{j_2}$.
Ainsi $S \times_{S_{i_5}} X_{i_5} \cong X$, comme voulu.
\end{proof}

\begin{lemma}
\label{lemma-descend-modules-finite-presentation}
Soit $I$ un ensemble ordonné filtrant.
Soit $(S_i, f_{ii'})$ un système projectif de schémas indexé par $I$.
Supposons que
\begin{enumerate}
\item tous les morphismes $f_{ii'} : S_i \to S_{i'}$ soient affines,
\item tous les schémas $S_i$ soient quasi-compacts et quasi-séparés.
\end{enumerate}
Soit $S = \lim_i S_i$. On a alors les propriétés suivantes :
\begin{enumerate}
\item Pour tout faisceau de $\mathcal{O}_S$-modules
$\mathcal{F}$ de présentation finie, il existe un indice
$i \in I$ et un faisceau de $\mathcal{O}_{S_i}$-modules de présentation finie
$\mathcal{F}_i$ tels que
$\mathcal{F} \cong f_i^*\mathcal{F}_i$.
\item Donnons-nous un indice $i \in I$, des faisceaux
de $\mathcal{O}_{S_i}$-modules $\mathcal{F}_i$, $\mathcal{G}_i$
de présentation finie et un morphisme
$\varphi : f_i^*\mathcal{F}_i \to f_i^*\mathcal{G}_i$ sur $S$.
Il existe alors un indice $i' \geq i$ et un morphisme
$\varphi_{i'} : f_{i'i}^*\mathcal{F}_i \to f_{i'i}^*\mathcal{G}_i$
dont le changement de base à $S$ est $\varphi$.
\item Donnons-nous un indice $i \in I$, des faisceaux de $\mathcal{O}_{S_i}$-modules
$\mathcal{F}_i$, $\mathcal{G}_i$ de présentation finie
et deux morphismes $\varphi_i, \psi_i : \mathcal{F}_i \to \mathcal{G}_i$.
Supposons leurs changements de base égaux : $f_i^*\varphi_i = f_i^*\psi_i$.
Il existe alors un indice $i' \geq i$ tel que
$f_{i'i}^*\varphi_i = f_{i'i}^*\psi_i$.
\end{enumerate}
Autrement dit, la catégorie des modules
de présentation finie sur $S$ est la colimite, indexée par $I$,
des catégories de modules de présentation finie sur $S_i$.
\end{lemma}

\begin{proof}
Nous esquissons deux démonstrations, en omettant les détails.

\medskip\noindent
Première démonstration. Si $S$ et les $S_i$ sont affines, ce lemme
équivaut à
Algèbre, Lemme \ref{algebra-lemma-colimit-category-fp-modules}.
Dans le cas général, on se ramène au cas affine par recollement de Zariski.

\medskip\noindent
Seconde démonstration. On utilise les faits suivants :
\begin{enumerate}
\item il existe une équivalence de catégories entre les
$\mathcal{O}_S$-modules quasi-cohérents et les fibrés vectoriels sur $S$ ; voir
Constructions, section \ref{constructions-section-vector-bundle},
\item un fibré vectoriel $\mathbf{V}(\mathcal{F}) \to S$ est
de présentation finie sur $S$ si et seulement si $\mathcal{F}$
est un $\mathcal{O}_S$-module de présentation finie.
\end{enumerate}
Ceci étant, on peut utiliser le Lemme \ref{lemma-descend-finite-presentation}
pour montrer que la catégorie des fibrés vectoriels de présentation finie
sur $S$ est la colimite, indexée par $I$,
des catégories de fibrés vectoriels sur $S_i$.
\end{proof}

\begin{lemma}
\label{lemma-descend-invertible-modules}
Soit $S = \lim S_i$ la limite d'un système projectif filtrant de schémas
quasi-compacts et quasi-séparés $S_i$, à morphismes de transition affines. Alors
\begin{enumerate}
\item tout $\mathcal{O}_S$-module localement libre de rang fini est l'image inverse
d'un $\mathcal{O}_{S_i}$-module localement libre de rang fini pour un certain $i$,
\item tout $\mathcal{O}_S$-module inversible est l'image inverse d'un
$\mathcal{O}_{S_i}$-module inversible pour un certain $i$,
\item tout idéal quasi-cohérent de type fini $\mathcal{I} \subset \mathcal{O}_S$
est de la forme $\mathcal{I}_i \cdot \mathcal{O}_S$ pour un certain $i$ et un
idéal quasi-cohérent de type fini $\mathcal{I}_i \subset \mathcal{O}_{S_i}$.
\end{enumerate}
\end{lemma}

\begin{proof}
Soit $\mathcal{E}$ un $\mathcal{O}_S$-module localement libre de rang fini. Puisque
les modules localement libres de rang fini sont de présentation finie, on peut trouver un $i$
et un $\mathcal{O}_{S_i}$-module $\mathcal{E}_i$ de présentation finie
tels que $f_i^*\mathcal{E}_i \cong \mathcal{E}$ ; voir
Lemme \ref{lemma-descend-modules-finite-presentation}.
En augmentant $i$, on peut supposer que $\mathcal{E}_i$ est un
$\mathcal{O}_{S_i}$-module plat ; voir
Algèbre, Lemme \ref{algebra-lemma-flat-finite-presentation-limit-flat}.
(L'emploi de ce lemme n'est pas nécessaire, mais commode.)
Alors $\mathcal{E}_i$ est localement libre de rang fini d'après
Algèbre, Lemme \ref{algebra-lemma-finite-projective}.

\medskip\noindent
Si $\mathcal{L}$ est un $\mathcal{O}_S$-module inversible,
ce qui précède fournit un $i$ et des
$\mathcal{O}_{S_i}$-modules localement libres de rang fini $\mathcal{L}_i$
et $\mathcal{N}_i$ dont les images inverses sont $\mathcal{L}$ et
$\mathcal{L}^{\otimes -1}$. Après avoir éventuellement augmenté
$i$, l'application
$\mathcal{L} \otimes_{\mathcal{O}_X} \mathcal{L}^{\otimes -1}
\to \mathcal{O}_X$ descend en une application
$\mathcal{L}_i \otimes_{\mathcal{O}_{S_i}} \mathcal{N}_i \to
\mathcal{O}_{S_i}$. En augmentant encore $i$, on
peut supposer que c'est un isomorphisme. Il s'ensuit que
$\mathcal{L}_i$ est un module inversible
(Modules, Lemme \ref{modules-lemma-invertible}), ce qui
achève la démonstration de (2).

\medskip\noindent
Étant donné $\mathcal{I}$ comme dans (3), on voit que
$\mathcal{O}_S \to \mathcal{O}_S/\mathcal{I}$
est une application de $\mathcal{O}_S$-modules de présentation finie.
D'après le Lemme \ref{lemma-descend-modules-finite-presentation},
c'est donc l'image inverse d'une
application $\mathcal{O}_{S_i} \to \mathcal{F}_i$ de
$\mathcal{O}_{S_i}$-modules de présentation finie. En augmentant $i$,
on peut supposer cette application surjective (détails omis ; indication : appliquer
Algèbre, Lemme \ref{algebra-lemma-module-map-property-in-colimit}
sur un recouvrement ouvert affine). Le noyau de $\mathcal{O}_{S_i} \to \mathcal{F}_i$
est alors un idéal quasi-cohérent de type fini de $\mathcal{O}_{S_i}$
dont l'image inverse donne $\mathcal{I}$.
\end{proof}

\begin{lemma}
\label{lemma-descend-module-flat-finite-presentation}
Reprenons les notations et les hypothèses du
Lemme \ref{lemma-descend-finite-presentation}.
Soit $i \in I$.
Supposons que $\varphi_i : X_i \to Y_i$ soit un morphisme de schémas
de présentation finie sur $S_i$ et que $\mathcal{F}_i$ soit un
$\mathcal{O}_{X_i}$-module quasi-cohérent de présentation finie.
Si l'image inverse de $\mathcal{F}_i$ sur $X_i \times_{S_i} S$ est plate
sur $Y_i \times_{S_i} S$, il existe un indice $i' \geq i$
tel que l'image inverse de $\mathcal{F}_i$ sur $X_i \times_{S_i} S_{i'}$
soit plate sur $Y_i \times_{S_i} S_{i'}$.
\end{lemma}

\begin{proof}
(Ce lemme est l'analogue du
Lemme \ref{lemma-descend-flat-finite-presentation}
pour les modules.)
Pour $i' \geq i$, notons $X_{i'} = S_{i'} \times_{S_i} X_i$,
$\mathcal{F}_{i'} = (X_{i'} \to X_i)^*\mathcal{F}_i$, et de même
pour $Y_{i'}$. Notons $\varphi_{i'}$ le changement de base
de $\varphi_i$ à $S_{i'}$. Posons aussi $X = S \times_{S_i} X_i$,
$Y =S \times_{S_i} X_i$, $\mathcal{F} = (X \to X_i)^*\mathcal{F}_i$
et notons $\varphi$ le changement de base de $\varphi_i$ à $S$.
Soit $Y_i = \bigcup_{j = 1, \ldots, m} V_{j, i}$ un recouvrement ouvert affine
fini tel que chaque $V_{j, i}$ ait son image dans un ouvert affine de $S_i$.
Pour chaque $j = 1, \ldots m$, soit
$\varphi_i^{-1}(V_{j, i}) = \bigcup_{k = 1, \ldots, m(j)} U_{k, j, i}$
un recouvrement ouvert affine fini. Pour $i' \geq i$, notons
$V_{j, i'}$ l'image réciproque de $V_{j, i}$ dans $Y_{i'}$ et
$U_{k, j, i'}$ celle de $U_{k, j, i}$ dans $X_{i'}$.
On a de même $U_{k, j} \subset X$ et $V_j \subset Y$.
Alors $U_{k, j} = \lim_{i' \geq i} U_{k, j, i'}$
et $V_j = \lim_{i' \geq i} V_j$
(voir Lemme \ref{lemma-directed-inverse-system-has-limit}).
Puisque $X_{i'} = \bigcup_{k, j} U_{k, j, i'}$ est un recouvrement ouvert fini,
il suffit de démontrer le lemme pour chacun des morphismes
$U_{k, j, i} \to V_{j, i}$ et le faisceau $\mathcal{F}_i|_{U_{k, j, i}}$.
On se ramène donc au cas où $X_i$ et
$Y_i$ sont affines et ont leur image dans un ouvert affine de $S_i$, c'est-à-dire
que l'on peut aussi supposer $S$ affine.

\medskip\noindent
Dans le cas affine, on se ramène au résultat d'algèbre suivant.
Supposons que $R = \colim_{i \in I} R_i$. Pour un certain $i \in I$,
donnons-nous une application $A_i \to B_i$ de $R_i$-algèbres de présentation finie.
Soit $N_i$ un $B_i$-module de présentation finie.
Si $R \otimes_{R_i} N_i$ est plat sur $R \otimes_{R_i} A_i$,
alors, pour un certain $i' \geq i$, le module
$R_{i'} \otimes_{R_i} N_i$ est plat sur $R_{i'} \otimes_{R_i} A$.
C'est exactement le résultat démontré dans
Algèbre,
Lemme \ref{algebra-lemma-flat-finite-presentation-limit-flat} (3).
\end{proof}

\begin{lemma}
\label{lemma-descend-finite-presentation-variant}
Pour un schéma $T$, notons $\mathcal{C}_T$ la sous-catégorie pleine des
schémas $W$ sur $T$ tels que $W$ soit quasi-compact et quasi-séparé
et que le morphisme structural $W \to T$ soit
localement de présentation finie.
Soit $S = \lim S_i$ une limite projective de schémas à morphismes
de transition affines. Il existe alors une équivalence
de catégories
$$
\colim \mathcal{C}_{S_i} \longrightarrow \mathcal{C}_S
$$
donnée par les foncteurs de changement de base.
\end{lemma}

\noindent
Attention : n'utilisez pas ce lemme sans comprendre la différence
entre celui-ci et le Lemme \ref{lemma-descend-finite-presentation}.

\begin{proof}
Pleine fidélité. Donnons-nous $i \in I$ et des objets
$X_i$, $Y_i$ de $\mathcal{C}_{S_i}$. Notons
$X = X_i \times_{S_i} S$ et $Y = Y_i \times_{S_i} S$.
Donnons-nous un morphisme $f : X \to Y$ sur $S$.
On peut choisir un recouvrement ouvert affine fini
$Y_i = V_{i, 1} \cup \ldots \cup V_{i, m}$
tel que $V_{i, j} \to Y_i \to S_i$ ait son image dans un
ouvert affine $W_{i, j}$ de $S_i$. Notons $Y = V_1 \cup \ldots \cup V_m$ le
recouvrement ouvert affine induit sur $Y$.
Puisque $f : X \to Y$ est quasi-compact
(Schémas, Lemme \ref{schemes-lemma-quasi-compact-permanence}),
on peut, en augmentant $i$, supposer qu'il existe un
recouvrement ouvert fini $X_i = U_{i, 1} \cup \ldots \cup U_{i, m}$
par des ouverts quasi-compacts tels que l'image réciproque de
$U_{i, j}$ dans $Y$ soit $f^{-1}(V_j)$ ; voir
Lemme \ref{lemma-descend-opens}.
D'après le Lemme \ref{lemma-descend-finite-presentation}
appliqué à $f|_{f^{-1}(V_j)}$ sur $W_j$, on peut supposer, en
augmentant $i$, qu'il existe un morphisme
$f_{i, j} : V_{i, j} \to U_{i, j}$ sur $S$ dont le changement de base
à $S$ est $f|_{f^{-1}(V_j)}$.
En augmentant encore $i$, on peut supposer que $f_{i, j}$ et $f_{i, j'}$
coïncident sur l'ouvert quasi-compact $U_{i, j} \cap U_{i, j'}$.
On peut alors recoller ces morphismes pour obtenir
le morphisme cherché $f_i : X_i \to Y_i$.
Ce morphisme est unique (quitte à augmenter $i$),
car c'est le cas des morphismes $f_{i, j}$.

\medskip\noindent
Pour montrer que le foncteur est essentiellement surjectif, on raisonne
exactement de la même manière. Supposons en effet que
$X$ soit un objet de $\mathcal{C}_S$. Choisissons $i \in I$.
On peut choisir un recouvrement ouvert affine fini
$X = U_1 \cup \ldots \cup U_m$ tel que $U_j \to X \to S \to S_i$
se factorise par un ouvert affine $W_{i, j} \subset S_i$.
Posons $W_j = W_{i, j} \times_{S_i} S$. C'est un ouvert affine de $S$.
D'après le Lemme \ref{lemma-descend-finite-presentation},
on peut, en augmentant $i$, supposer qu'il existe des morphismes
$U_{i, j} \to W_{i, j}$ de présentation finie
dont le changement de base à $W_j$ est $U_j$.
En augmentant $i$, on peut supposer qu'il existe
des ouverts quasi-compacts $U_{i, j, j'} \subset U_{i, j}$
dont les changements de base à $S$ sont égaux à $U_j \cap U_{j'}$.
Nous affirmons qu'en augmentant $i$, on peut supposer que l'image du morphisme
$U_{i, j, j'} \to U_{i, j} \to W_{i, j}$
est contenue dans $W_{i, j} \cap W_{i, j'}$.
En effet, le complémentaire de $W_{i, j} \cap W_{i, j'}$
est fermé dans le schéma affine $W_{i, j}$, donc affine.
Puisque $U_j \cap U_{j'} = \lim U_{i, j, j'}$ a bien son image dans
$W_{i, j} \cap W_{i, j'}$,
on peut appliquer le Lemme \ref{lemma-limit-fibre-product-empty}
pour obtenir l'assertion. Ainsi, on peut considérer
$$
U_{i, j, j'} \quad\text{et}\quad U_{i, j', j}
$$
comme deux schémas sur $W_{i, j'}$ dont les changements de base à $W_{j'}$
redonnent $U_j \cap U_{j'}$. En augmentant $i$ et en utilisant le
Lemme \ref{lemma-descend-finite-presentation},
on peut donc supposer qu'il existe des isomorphismes
$U_{i, j, j'} \to U_{i, j', j}$ sur $W_{i, j'}$, donc
sur $S_i$. En augmentant encore $i$ (détails omis),
on peut supposer que ces isomorphismes
satisfont la condition de cocycle mentionnée dans
Schémas, section \ref{schemes-section-glueing-schemes}.
En appliquant Schémas, Lemme \ref{schemes-lemma-glue},
on obtient un objet $X_i$ de $\mathcal{C}_{S_i}$ dont le
changement de base à $S$ est isomorphe à $X$ ; nous omettons
certaines vérifications.
\end{proof}












\section{Caractérisation des schémas affines}
\label{section-affine}

\noindent
Si $f : X \to S$ est un morphisme entier surjectif de schémas
tel que $X$ soit affine, alors $S$ est lui aussi affine.
Voir \cite[A.2]{Conrad-Nagata}. Notre démonstration repose
sur le cas noethérien, énoncé et démontré dans Cohomologie des schémas,
Lemme \ref{coherent-lemma-image-affine-finite-morphism-affine-Noetherian}.
Voir aussi \cite[II 6.7.1]{EGA}.

\begin{lemma}
\label{lemma-affine}
\begin{slogan}
Un schéma qui admet un morphisme fini surjectif provenant d'un schéma affine est affine.
\end{slogan}
Soit $f : X \to S$ un morphisme de schémas.
Supposons que $f$ soit surjectif et fini, et que $X$ soit affine.
Alors $S$ est affine.
\end{lemma}

\begin{proof}
Puisque $f$ est surjectif et $X$ quasi-compact, $S$ est
quasi-compact. Puisque $X$ est séparé et $f$ surjectif et
universellement fermé (Morphismes, Lemme
\ref{morphisms-lemma-integral-universally-closed}), $S$
est séparé (Morphismes, Lemme
\ref{morphisms-lemma-image-universally-closed-separated}).

\medskip\noindent
D'après le Lemme \ref{lemma-finite-in-finite-and-finite-presentation},
on peut écrire $X = \lim_a X_a$, où $X_a \to S$ est fini et de présentation
finie. D'après le Lemme \ref{lemma-limit-affine}, $X_a$
est affine pour un certain $a \in A$. En remplaçant $X$ par $X_a$, on peut supposer
que $X \to S$ est surjectif, fini et de présentation finie, et
que $X$ est affine.

\medskip\noindent
D'après la Proposition \ref{proposition-approximate}, on peut écrire
$S = \lim_{i \in I} S_i$ comme
une limite projective de schémas de type fini sur $\mathbf{Z}$.
D'après le Lemme \ref{lemma-descend-finite-presentation}, on peut,
en restreignant $I$, supposer qu'il existe des schémas $X_i \to S_i$
de présentation finie tels que $X_{i'} = X_i \times_S S_{i'}$
pour $i' \geq i$ et que $X = \lim_i X_i$. D'après le
Lemme \ref{lemma-descend-finite-finite-presentation}, on peut
également supposer que $X_i \to S_i$ est fini pour tout $i \in I$.
En appliquant de nouveau le Lemme \ref{lemma-limit-affine}, on peut supposer $X_i$
affine pour tout $i \in I$. Le résultat découle donc du
cas noethérien ; voir Cohomologie des schémas,
Lemme \ref{coherent-lemma-image-affine-finite-morphism-affine-Noetherian}.
\end{proof}

\begin{proposition}
\label{proposition-affine}
\begin{slogan}
Un schéma qui admet un morphisme entier surjectif provenant d'un schéma affine est affine.
\end{slogan}
Soit $f : X \to S$ un morphisme de schémas. Supposons $X$ affine
et $f$ surjectif et universellement fermé\footnote{Un morphisme entier
est universellement fermé ; voir
Morphismes, Lemme \ref{morphisms-lemma-integral-universally-closed}.}.
Alors $S$ est affine.
\end{proposition}

\begin{proof}
D'après Morphismes, Lemme \ref{morphisms-lemma-image-universally-closed-separated},
le schéma $S$ est séparé. D'après
Morphismes, Lemme \ref{morphisms-lemma-affine-permanence},
$f$ est alors affine. Dès lors, d'après
Morphismes, Lemme \ref{morphisms-lemma-integral-universally-closed},
$f$ est entier.

\medskip\noindent
D'après le paragraphe précédent, on peut supposer $f : X \to S$
surjectif et entier, $X$ affine et $S$ séparé.
Puisque $f$ est surjectif et $X$ quasi-compact, $S$ est aussi
quasi-compact.

\medskip\noindent
D'après le Lemme \ref{lemma-integral-limit-finite-and-finite-presentation},
on peut écrire $X = \lim_i X_i$, où $X_i \to S$ est fini. D'après le
Lemme \ref{lemma-limit-affine},
pour $i$ assez grand, le schéma $X_i$ est affine.
De plus, puisque $X \to S$ se factorise par chaque $X_i$,
$X_i \to S$ est surjectif. On conclut donc que $S$ est affine par le
Lemme \ref{lemma-affine}.
\end{proof}

\begin{lemma}
\label{lemma-affines-glued-in-closed-affine}
Soit $X$ un schéma qui est, ensemblistement, réunion d'un
nombre fini de sous-schémas fermés affines. Alors $X$ est affine.
\end{lemma}

\begin{proof}
Soient $Z_i \subset X$, $i = 1, \ldots, n$, des sous-schémas fermés affines tels que
$X = \bigcup Z_i$ ensemblistement. Alors $\coprod Z_i \to X$ est surjectif
et entier, de source affine. Ainsi $X$ est affine par la
Proposition \ref{proposition-affine}.
\end{proof}

\begin{lemma}
\label{lemma-ample-on-reduction}
Soit $i : Z \to X$ une immersion fermée de schémas
induisant un homéomorphisme des espaces topologiques sous-jacents.
Soit $\mathcal{L}$ un faisceau inversible sur $X$.
Alors $i^*\mathcal{L}$ est ample sur $Z$ si et seulement si $\mathcal{L}$
est ample sur $X$.
\end{lemma}

\begin{proof}
Si $\mathcal{L}$ est ample, alors $i^*\mathcal{L}$ est ample, par
exemple d'après Morphismes, Lemme
\ref{morphisms-lemma-pullback-ample-tensor-relatively-ample}.
Supposons $i^*\mathcal{L}$ ample. Alors $Z$ est quasi-compact
(Propriétés, Définition \ref{properties-definition-ample})
et séparé
(Propriétés, Lemme \ref{properties-lemma-ample-separated}).
Puisque $i$ est surjectif, $X$ est quasi-compact.
Puisque $i$ est universellement fermé et surjectif,
$X$ est séparé (Morphismes, Lemme
\ref{morphisms-lemma-image-universally-closed-separated}).

\medskip\noindent
D'après la Proposition \ref{proposition-approximate}, on peut écrire
$X = \lim X_i$ comme une limite projective de schémas de type fini sur $\mathbf{Z}$
à morphismes de transition affines. On peut trouver un $i$ et un faisceau
inversible $\mathcal{L}_i$ sur $X_i$ dont l'image inverse sur $X$ est isomorphe à
$\mathcal{L}$ ; voir Lemme \ref{lemma-descend-modules-finite-presentation}.

\medskip\noindent
Pour chaque $i$, soit $Z_i \subset X_i$ l'image schématique
du morphisme $Z \to X_i$. Si $\Spec(A_i) \subset X_i$ est un sous-schéma ouvert
affine dont l'image réciproque est $\Spec(A)$ dans $X$ et si
$Z \cap \Spec(A)$ est défini par l'idéal $I \subset A$,
alors $Z_i \cap \Spec(A_i)$ est défini par l'idéal $I_i \subset A_i$,
image réciproque de $I$ dans $A_i$ par l'homomorphisme
$A_i \to A$ ; voir
Morphismes, Exemple \ref{morphisms-example-scheme-theoretic-image}.
Puisque $\colim A_i/I_i = A/I$, on en déduit $\lim Z_i = Z$.
D'après le Lemme \ref{lemma-limit-ample}, $\mathcal{L}_i|_{Z_i}$
est ample pour un certain $i$. Puisque $Z$, et donc $X$, a son image dans $Z_i$
ensemblistement, $X_{i'} \to X_i$ a son image dans $Z_i$
ensemblistement pour un certain $i' \geq i$ ; voir
Lemme \ref{lemma-limit-contained-in-constructible}.
(Comme $X_i$ est noethérien, toute partie fermée
de $X_i$ est constructible.) Soit $T \subset X_{i'}$
l'image réciproque schématique de $Z_i$ dans $X_{i'}$.
Remarquons que $\mathcal{L}_{i'}|_T$ est l'image inverse
de $\mathcal{L}_i|_{Z_i}$, donc est ample d'après
Morphismes, Lemme \ref{morphisms-lemma-pullback-ample-tensor-relatively-ample},
car $T \to Z_i$ est un morphisme affine.
Ainsi $\mathcal{L}_{i'}$ est ample sur $X_{i'}$
d'après Cohomologie des schémas, Lemme \ref{coherent-lemma-ample-on-reduction}.
En prenant l'image inverse sur $X$ (à l'aide du même lemme que ci-dessus),
on trouve que $\mathcal{L}$ est ample.
\end{proof}

\begin{lemma}
\label{lemma-thickening-quasi-affine}
Soit $i : Z \to X$ une immersion fermée de schémas
induisant un homéomorphisme des espaces topologiques sous-jacents.
Alors $X$ est quasi-affine si et seulement si $Z$ est quasi-affine.
\end{lemma}

\begin{proof}
Rappelons qu'un schéma est quasi-affine
si et seulement si son faisceau structural est ample ; voir
Propriétés, Lemme \ref{properties-lemma-quasi-affine-O-ample}.
Ainsi, si $Z$ est quasi-affine, $\mathcal{O}_Z$ est ample,
donc $\mathcal{O}_X$ est ample par le
Lemme \ref{lemma-ample-on-reduction}, et
$X$ est donc quasi-affine. La réciproque, qui
se démontre aussi de manière élémentaire, s'obtient
en lisant le raisonnement précédent en sens inverse.
\end{proof}

\noindent
Le lemme suivant n'a pas vraiment sa place dans cette section.

\begin{lemma}
\label{lemma-ample-profinite-set-in-principal-affine}
Soit $X$ un schéma. Soit $\mathcal{L}$ un faisceau inversible ample sur $X$.
Supposons donnés des morphismes de schémas
$$
\Spec(k) \leftarrow \Spec(A) \to W \subset X
$$
où $k$ est un corps, $A$ une $k$-algèbre entière et $W$ un ouvert de $X$.
Il existe alors un $n > 0$ et une section
$s \in \Gamma(X, \mathcal{L}^{\otimes n})$ tels que
$X_s$ soit affine, $X_s \subset W$, et que $\Spec(A) \to W$ se factorise par $X_s$.
\end{lemma}

\begin{proof}
Puisque $\Spec(A)$ est quasi-compact, on peut remplacer $W$ par un ouvert
quasi-compact contenant encore l'image de $\Spec(A) \to X$.
Rappelons que $X$ est quasi-séparé et quasi-compact du fait
qu'il possède un faisceau inversible ample ; voir
Propriétés, Définition \ref{properties-definition-ample} et
Lemme \ref{properties-lemma-affine-s-opens-cover-quasi-separated}.
D'après la Proposition \ref{proposition-approximate}, on peut
écrire $X = \lim X_i$ comme la limite d'un système projectif filtrant
de schémas de type fini sur $\mathbf{Z}$
à morphismes de transition affines.
Pour un certain $i$, le faisceau inversible ample $\mathcal{L}$ sur $X$
descend en un faisceau inversible ample $\mathcal{L}_i$ sur $X_i$,
et l'ouvert $W$ est l'image réciproque d'un ouvert
quasi-compact $W_i \subset X_i$ ; voir
Lemmes \ref{lemma-limit-ample}, \ref{lemma-descend-invertible-modules} et
\ref{lemma-descend-opens}.
On peut remplacer $X, W, \mathcal{L}$ par $X_i, W_i, \mathcal{L}_i$
et supposer $X$ de présentation finie sur $\mathbf{Z}$.
Écrivons $A = \colim A_j$ comme la limite inductive de ses $k$-sous-algèbres finies.
Pour un certain $j$, le morphisme $\Spec(A) \to X$ se factorise alors par
un morphisme $\Spec(A_j) \to X$ ; voir
Proposition \ref{proposition-characterize-locally-finite-presentation}.
Puisque $\Spec(A_j)$ est fini, on est ramené à
Propriétés, Lemme \ref{properties-lemma-ample-finite-set-in-principal-affine}.
\end{proof}



















\section{Variantes du lemme de Chow}
\label{section-chows-lemma}

\noindent
Dans cette section, nous démontrons plusieurs variantes du lemme de Chow.
La version la plus intéressante est probablement le cas noethérien,
énoncé et démontré dans
Cohomologie des schémas, section \ref{coherent-section-chows-lemma}.

\begin{lemma}
\label{lemma-chow-finite-type}
Soit $S$ un schéma quasi-compact et quasi-séparé.
Soit $f : X \to S$ un morphisme séparé de type fini.
Il existe alors un $n \geq 0$ et un diagramme
$$
\xymatrix{
X \ar[rd] & X' \ar[d] \ar[l]^\pi \ar[r] & \mathbf{P}^n_S \ar[dl] \\
& S &
}
$$
où $X' \to \mathbf{P}^n_S$ est une immersion et
$\pi : X' \to X$ est propre et surjectif.
\end{lemma}

\begin{proof}
D'après la Proposition \ref{proposition-separated-closed-in-finite-presentation},
on peut trouver une immersion fermée $X \to Y$, où $Y$ est séparé
et de présentation finie sur $S$. Si l'on démontre l'assertion
pour $Y$, le résultat pour $X$ s'en déduit immédiatement. On peut donc supposer
$X$ de présentation finie sur $S$.

\medskip\noindent
Écrivons $S = \lim_i S_i$ comme une limite projective de schémas noethériens ; voir
Proposition \ref{proposition-approximate}. D'après le
Lemme \ref{lemma-descend-finite-presentation}, on peut
trouver un indice $i \in I$ et un schéma $X_i \to S_i$ de présentation finie
tels que $X = S \times_{S_i} X_i$.
D'après le Lemme \ref{lemma-descend-separated-finite-presentation},
on peut supposer $X_i \to S_i$ séparé.
Si l'on démontre l'assertion pour
$X_i$ sur $S_i$, elle vaut évidemment pour $X$. Le cas
$X_i \to S_i$ est traité dans
Cohomologie des schémas, Lemme \ref{coherent-lemma-chow-Noetherian}.
\end{proof}

\begin{remark}
\label{remark-chow-finite-type}
Dans la situation du Lemme de Chow \ref{lemma-chow-finite-type} :
\begin{enumerate}
\item Le morphisme $\pi$ est en fait H-projectif (donc projectif ; voir
Morphismes, Lemme \ref{morphisms-lemma-H-projective}),
puisque le morphisme $X' \to \mathbf{P}^n_S \times_S X = \mathbf{P}^n_X$
est une immersion fermée (utiliser le fait que $\pi$ est propre ; voir
Morphismes, Lemme \ref{morphisms-lemma-image-proper-scheme-closed}).
\item On peut supposer $X'$ réduit, car on peut remplacer $X'$ par
son réduit sans changer les autres assertions du lemme.
\item On peut supposer $X' \to X$ de présentation finie
sans changer les autres assertions du lemme. Cela se déduit
de la démonstration du Lemme \ref{lemma-chow-finite-type}, mais on peut aussi le démontrer
directement comme suit. D'après (1), on a une immersion fermée
$X' \to \mathbf{P}^n_X$. D'après le
Lemme \ref{lemma-closed-is-limit-closed-and-finite-presentation}, on peut écrire
$X' = \lim X'_i$, où $X'_i \to \mathbf{P}^n_X$ est une immersion
fermée de présentation finie. En particulier, $X'_i \to X$ est
de présentation finie, propre et surjectif. Pour
$i$ assez grand, le morphisme $X'_i \to \mathbf{P}^n_S$ est une immersion par le
Lemme \ref{lemma-finite-type-eventually-closed}. En remplaçant
$X'$ par $X'_i$, on obtient le résultat voulu.
\end{enumerate}
Bien entendu, on ne peut pas, en général, obtenir simultanément (2) et (3).
\end{remark}

\noindent
Voici une variante du lemme de Chow dans laquelle on suppose que le schéma
source possède un nombre fini de composantes irréductibles.

\begin{lemma}
\label{lemma-chow-EGA}
Soit $S$ un schéma quasi-compact et quasi-séparé.
Soit $f : X \to S$ un morphisme séparé de type fini.
Supposons que $X$ possède un nombre fini de composantes irréductibles.
Il existe alors un $n \geq 0$ et un diagramme
$$
\xymatrix{
X \ar[rd] & X' \ar[d] \ar[l]^\pi \ar[r] & \mathbf{P}^n_S \ar[dl] \\
& S &
}
$$
où $X' \to \mathbf{P}^n_S$ est une immersion et
$\pi : X' \to X$ est propre et surjectif. De plus, il existe
un sous-schéma ouvert dense $U \subset X$ tel que $\pi^{-1}(U) \to U$
soit un isomorphisme de schémas.
\end{lemma}

\begin{proof}
Soit $X = Z_1 \cup \ldots \cup Z_n$ la décomposition de $X$
en composantes irréductibles. Soit $\eta_j \in Z_j$ le point générique.

\medskip\noindent
Il y a au moins deux façons de poursuivre la démonstration.
La première consiste à reprendre la démonstration de
Cohomologie des schémas, Lemme \ref{coherent-lemma-chow-Noetherian},
en utilisant le résultat général
Propriétés, Lemme \ref{properties-lemma-point-and-maximal-points-affine},
pour trouver des ouverts affines convenables de $X$. (C'est la démonstration « standard ».)
La seconde consiste à utiliser l'approximation noethérienne absolue, comme dans
la démonstration du Lemme \ref{lemma-chow-finite-type} ci-dessus.
C'est ce que nous allons faire ici.

\medskip\noindent
D'après la Proposition \ref{proposition-separated-closed-in-finite-presentation},
on peut trouver une immersion fermée $X \to Y$, où $Y$ est séparé
et de présentation finie sur $S$.
Écrivons $S = \lim_i S_i$ comme une limite projective de schémas noethériens ; voir
Proposition \ref{proposition-approximate}. D'après le
Lemme \ref{lemma-descend-finite-presentation}, on peut
trouver un indice $i \in I$ et un schéma $Y_i \to S_i$ de présentation finie
tels que $Y = S \times_{S_i} Y_i$.
D'après le Lemme \ref{lemma-descend-separated-finite-presentation},
on peut supposer $Y_i \to S_i$ séparé.
On a le diagramme suivant :
$$
\xymatrix{
\eta_j \in Z_j \ar[r] & X \ar[r] \ar[rd] & Y \ar[r] \ar[d] & Y_i \ar[d] \\
& & S \ar[r] & S_i
}
$$
Notons $h : X \to Y_i$ la composée.

\medskip\noindent
Pour $i' \geq i$, posons $Y_{i'} = S_{i'} \times_{S_i} Y_i$.
Alors $Y = \lim_{i' \geq i} Y_{i'}$ ; voir
Lemme \ref{lemma-scheme-over-limit}.
Choisissons $j, j' \in \{1, \ldots, n\}$, $j \not = j'$.
Remarquons que $\eta_j$ n'est pas une spécialisation de $\eta_{j'}$.
D'après le Lemme \ref{lemma-topology-limit},
on peut remplacer $i$ par un indice plus grand et supposer
que $h(\eta_j)$ n'est pas une spécialisation de $h(\eta_{j'})$
pour tous les couples $(j, j')$ ci-dessus.
Pour un tel indice, soit
$Y' \subset Y_i$ l'image schématique de
$h : X \to Y_i$ ; voir
Morphismes, Définition \ref{morphisms-definition-scheme-theoretic-image}.
Le morphisme $h$ est quasi-compact, étant la composée des morphismes quasi-compacts
$X \to Y$ et $Y \to Y_i$ (ce dernier est affine).
D'après
Morphismes, Lemme \ref{morphisms-lemma-quasi-compact-scheme-theoretic-image},
le morphisme $X \to Y'$ est donc dominant. Ainsi, les points génériques
de $Y'$ appartiennent tous à l'ensemble
$\{h(\eta_1), \ldots, h(\eta_n)\}$ ; voir
Morphismes, Lemme \ref{morphisms-lemma-quasi-compact-dominant}.
Puisqu'aucun des $h(\eta_j)$ n'est une spécialisation d'un autre,
les points $h(\eta_1), \ldots, h(\eta_n)$ sont deux à deux
distincts et chacun est un point générique de $Y'$.

\medskip\noindent
Appliquons Cohomologie des schémas, Lemme
\ref{coherent-lemma-chow-Noetherian}, au morphisme
$Y' \to S_i$. On obtient un diagramme
$$
\xymatrix{
Y' \ar[rd] & Y^* \ar[d] \ar[l]^\pi \ar[r] & \mathbf{P}^n_{S_i} \ar[dl] \\
& S_i &
}
$$
tel que $\pi$ soit propre et surjectif, et un isomorphisme au-dessus
d'un sous-schéma ouvert dense $V \subset Y'$. Par le choix de $i$ ci-dessus,
on sait que $h(\eta_1), \ldots, h(\eta_n) \in V$. Considérons
le diagramme commutatif
$$
\xymatrix{
X' \ar@{=}[r] &
X \times_{Y'} Y^* \ar[r] \ar[d] &
Y^* \ar[r] \ar[d] &
\mathbf{P}^n_{S_i} \ar[ddl] \\
& X \ar[r] \ar[d] & Y' \ar[d] & \\
& S \ar[r] & S_i &
}
$$
Remarquons que $X' \to X$ est un isomorphisme au-dessus du sous-schéma ouvert
$U = h^{-1}(V)$, qui contient chacun des $\eta_j$, donc est
dense dans $X$. On conclut que $X \leftarrow X' \rightarrow \mathbf{P}^n_S$
résout le problème posé dans le lemme.
\end{proof}













\section{Applications du lemme de Chow}
\label{section-apply-chow}

\noindent
Voici une première application du lemme de Chow.

\begin{lemma}
\label{lemma-eventually-proper}
\begin{slogan}
Si le changement de base d'un schéma à une limite est propre, alors
le changement de base est déjà propre à un étage fini.
\end{slogan}
Reprenons les hypothèses et les notations de la Situation \ref{situation-descent-property}.
Si
\begin{enumerate}
\item $f$ est propre,
\item $f_0$ est localement de type fini,
\end{enumerate}
alors il existe un $i$ tel que $f_i$ soit propre.
\end{lemma}

\begin{proof}
D'après le Lemme \ref{lemma-descend-separated-finite-presentation},
$f_i$ est séparé pour un certain $i \geq 0$. En remplaçant
$0$ par $i$, on peut supposer $f_0$ séparé.
Remarquons que $f_0$ est quasi-compact ; voir
Schémas, Lemme \ref{schemes-lemma-quasi-compact-permanence}.
D'après le Lemme \ref{lemma-chow-finite-type}, on peut choisir un diagramme
$$
\xymatrix{
X_0 \ar[rd] & X_0' \ar[d] \ar[l]^\pi \ar[r] & \mathbf{P}^n_{Y_0} \ar[dl] \\
& Y_0 &
}
$$
où $X_0' \to \mathbf{P}^n_{Y_0}$ est une immersion et
$\pi : X_0' \to X_0$ est propre et surjectif. Posons
$X' = X_0' \times_{Y_0} Y$ et $X_i' = X_0' \times_{Y_0} Y_i$.
D'après Morphismes, Lemmes \ref{morphisms-lemma-composition-proper} et
\ref{morphisms-lemma-base-change-proper},
$X' \to Y$ est propre. Ainsi $X' \to \mathbf{P}^n_Y$ est
une immersion fermée (Morphismes, Lemme
\ref{morphisms-lemma-image-proper-scheme-closed}). D'après
Morphismes, Lemme \ref{morphisms-lemma-image-proper-is-proper},
il suffit de montrer que $X'_i \to Y_i$ est propre pour un certain $i$.
D'après le Lemme \ref{lemma-descend-closed-immersion-finite-presentation},
$X'_i \to \mathbf{P}^n_{Y_i}$ est
une immersion fermée pour $i$ assez grand. Alors $X'_i \to Y_i$
est propre, ce qui conclut.
\end{proof}

\begin{lemma}
\label{lemma-proper-limit-of-proper-finite-presentation}
Soit $f : X \to S$ un morphisme propre, où $S$ est quasi-compact et
quasi-séparé. Alors $X = \lim X_i$ est une limite projective de schémas
$X_i$ propres et de présentation finie sur $S$, telle que
tous les morphismes de transition et les morphismes $X \to X_i$ soient des immersions
fermées.
\end{lemma}

\begin{proof}
D'après la Proposition \ref{proposition-separated-closed-in-finite-presentation},
on peut trouver une immersion fermée $X \to Y$, où $Y$ est séparé et de
présentation finie sur $S$. D'après le Lemme \ref{lemma-chow-finite-type},
on peut trouver un diagramme
$$
\xymatrix{
Y \ar[rd] & Y' \ar[d] \ar[l]^\pi \ar[r] & \mathbf{P}^n_S \ar[dl] \\
& S &
}
$$
où $Y' \to \mathbf{P}^n_S$ est une immersion et
$\pi : Y' \to Y$ est propre et surjectif. D'après le
Lemme \ref{lemma-closed-is-limit-closed-and-finite-presentation},
on peut écrire $X = \lim X_i$, où $X_i \to Y$ est une immersion fermée de
présentation finie. Notons $X'_i \subset Y'$, resp.\ $X' \subset Y'$,
l'image réciproque schématique de $X_i \subset Y$, resp.\ $X \subset Y$.
Alors $\lim X'_i = X'$. Puisque $X' \to S$ est propre
(Morphismes, Lemmes \ref{morphisms-lemma-composition-proper}),
$X' \to \mathbf{P}^n_S$ est une immersion fermée (Morphismes, Lemme
\ref{morphisms-lemma-image-proper-scheme-closed}). Pour $i$ assez grand,
$X'_i \to \mathbf{P}^n_S$ est donc une immersion fermée d'après le
Lemme \ref{lemma-eventually-closed-immersion}.
Ainsi $X'_i$ est propre sur $S$.
Pour un tel $i$, le morphisme $X_i \to S$ est propre d'après
Morphismes, Lemme \ref{morphisms-lemma-image-proper-is-proper}.
\end{proof}

\begin{lemma}
\label{lemma-proper-limit-of-proper-finite-presentation-noetherian}
Soit $f : X \to S$ un morphisme propre, où $S$ est quasi-compact et
quasi-séparé. Il existe alors un ensemble ordonné filtrant $I$ et un
système projectif $(f_i : X_i \to S_i)$ de morphismes de schémas indexé par $I$,
tels que les morphismes de transition $X_i \to X_{i'}$ et $S_i \to S_{i'}$
soient affines, que $f_i$ soit propre, que $S_i$ soit de type
fini sur $\mathbf{Z}$ et que
$(X \to S) = \lim (X_i \to S_i)$.
\end{lemma}

\begin{proof}
D'après le Lemme \ref{lemma-proper-limit-of-proper-finite-presentation},
on peut écrire $X = \lim_{k \in K} X_k$, où $X_k \to S$ est propre et
de présentation finie. Ensuite, par approximation noethérienne absolue
(Proposition \ref{proposition-approximate}), on peut
écrire $S = \lim_{j \in J} S_j$, où $S_j$ est de type fini sur $\mathbf{Z}$.
Pour chaque $k$, il existe un $j$ et un morphisme $X_{k, j} \to S_j$
de présentation finie tels que $X_k \cong S \times_{S_j} X_{k, j}$
comme schémas sur $S$ ; voir
Lemme \ref{lemma-descend-finite-presentation}.
En augmentant $j$, on peut supposer $X_{k, j} \to S_j$ propre ; voir
Lemme \ref{lemma-eventually-proper}. L'ensemble $I$ sera constitué
de ces couples $(k, j)$, et le morphisme correspondant sera $X_{k, j} \to S_j$.
Pour tout $k' \geq k$, on peut trouver un $j' \geq j$ et un morphisme
$X_{j', k'} \to X_{j, k}$ au-dessus de $S_{j'} \to S_j$ dont le changement de base à $S$
donne le morphisme $X_{k'} \to X_k$ (cela découle encore du
Lemme \ref{lemma-descend-finite-presentation}).
Ces morphismes constituent les morphismes de transition du système. Certains détails
sont omis.
\end{proof}

\begin{lemma}
\label{lemma-finite-type-eventually-proper}
Soit $S$ un schéma. Soit $X = \lim X_i$ une limite projective de
schémas sur $S$ à morphismes de transition affines. Soit $Y \to X$
un morphisme de schémas sur $S$.
Si $Y \to X$ est propre, si les $X_i$ sont quasi-compacts et quasi-séparés et si
$Y$ est localement de type fini sur $S$, alors $Y \to X_i$ est propre
pour $i$ assez grand.
\end{lemma}

\begin{proof}
Choisissons une immersion fermée $Y \to Y'$, où $Y'$ est propre et de présentation
finie sur $X$ ; voir
Lemme \ref{lemma-proper-limit-of-proper-finite-presentation}.
Choisissons ensuite un $i$ et un morphisme propre $Y'_i \to X_i$
tels que $Y' = X \times_{X_i} Y'_i$. C'est possible d'après les
Lemmes \ref{lemma-descend-finite-presentation} et
\ref{lemma-eventually-proper}. En remplaçant $i$
par un indice plus grand, $Y \to Y'_i$ est alors une immersion
fermée ; voir Lemme \ref{lemma-finite-type-eventually-closed}.
\end{proof}

\noindent
Rappelons la notion de support schématique
d'un module quasi-cohérent de type fini ; voir
Morphismes, Définition \ref{morphisms-definition-scheme-theoretic-support}.

\begin{lemma}
\label{lemma-eventually-proper-support}
Reprenons les hypothèses et les notations de la Situation \ref{situation-descent-property}.
Soit $\mathcal{F}_0$ un $\mathcal{O}_{X_0}$-module quasi-cohérent.
Notons $\mathcal{F}$ et $\mathcal{F}_i$ les images inverses de
$\mathcal{F}_0$ sur $X$ et $X_i$. Faisons les hypothèses suivantes :
\begin{enumerate}
\item $f_0$ est localement de type fini,
\item $\mathcal{F}_0$ est de type fini,
\item le support schématique de $\mathcal{F}$ est propre sur $Y$.
\end{enumerate}
Alors le support schématique de $\mathcal{F}_i$ est propre sur $Y_i$
pour un certain $i$.
\end{lemma}

\begin{proof}
On peut remplacer $X_0$ par le support schématique de $\mathcal{F}_0$.
D'après Morphismes, Lemme \ref{morphisms-lemma-support-finite-type}, cela
garantit que $X_i$ est le support de $\mathcal{F}_i$ et $X$ le
support de $\mathcal{F}$. Si $Z \subset X$ désigne le support
schématique de $\mathcal{F}$, alors $Z \to X$ est un homéomorphisme
universel. On conclut que $X \to Y$ est propre, puisque
$Z \to Y$ l'est par hypothèse ; voir
Morphismes, Lemme \ref{morphisms-lemma-image-proper-is-proper}.
D'après le Lemme \ref{lemma-eventually-proper}, $X_i \to Y$ est propre
pour un certain $i$. Le support schématique $Z_i$ de
$\mathcal{F}_i$ est donc propre sur $Y$ d'après
Morphismes, Lemmes \ref{morphisms-lemma-closed-immersion-proper} et
\ref{morphisms-lemma-composition-proper}.
\end{proof}











\section{Morphismes universellement fermés}
\label{section-universally-closed}

\noindent
Dans cette section, nous étudions quand un morphisme quasi-compact (mais
pas nécessairement séparé) est universellement fermé. Nous démontrons d'abord
un lemme permettant de tester cette propriété après un changement de base
localement de présentation finie.

\begin{lemma}
\label{lemma-separate}
Soit $f : X \to S$ un morphisme quasi-compact de schémas.
Soit $g : T \to S$ un morphisme de schémas.
Soit $t \in T$ un point et $Z \subset X_T$ un sous-schéma
fermé tel que $Z \cap X_t = \emptyset$.
Il existe alors un voisinage ouvert
$V \subset T$ de $t$, un diagramme commutatif
$$
\xymatrix{
V \ar[d] \ar[r]_a & T' \ar[d]^b \\
T \ar[r]^g & S,
}
$$
et un sous-schéma fermé $Z' \subset X_{T'}$ tels que
\begin{enumerate}
\item le morphisme $b : T' \to S$ soit localement de présentation finie,
\item en posant $t' = a(t)$, on ait $Z' \cap X_{t'} = \emptyset$,
\item $Z \cap X_V$ ait son image dans $Z'$ par le morphisme $X_V \to X_{T'}$.
\end{enumerate}
De plus, on peut supposer $V$ et $T'$ affines.
\end{lemma}

\begin{proof}
Posons $s = g(t)$. Au cours de la démonstration, on peut toujours remplacer $T$ par un
voisinage ouvert de $t$. On peut donc aussi remplacer $S$ par un voisinage ouvert
de $s$. On peut ainsi supposer, et l'on suppose, $T$ et $S$ affines.
Écrivons $S = \Spec(A)$, $T = \Spec(B)$ ; soit $g$ donné par
l'homomorphisme $A \to B$, et soit $t$ associé à l'idéal premier
$\mathfrak q \subset B$.

\medskip\noindent
Puisque $X \to S$ est quasi-compact et $S$ affine, on peut écrire
$X = \bigcup_{i = 1, \ldots, n} U_i$ comme réunion finie d'ouverts affines.
Écrivons $U_i = \Spec(C_i)$. On a en particulier
$X_T = \bigcup_{i = 1, \ldots, n} U_{i, T} =
\bigcup_{i = 1, \ldots n} \Spec(C_i \otimes_A B)$.
Soit $I_i \subset C_i \otimes_A B$ l'idéal correspondant au
sous-schéma fermé $Z \cap U_{i, T}$. La condition
$Z \cap X_t = \emptyset$ signifie que $I_i$ engendre
l'idéal unité de l'anneau
$$
C_i \otimes_A \kappa(\mathfrak q) =
(B \setminus \mathfrak q)^{-1}\left(
C_i \otimes_A B/\mathfrak q C_i \otimes_A B \right)
$$
Puisque $I_i (B \setminus \mathfrak q)^{-1}(C_i \otimes_A B) =
(B \setminus \mathfrak q)^{-1} I_i$, cela signifie que $1 = x_i/g_i$
pour certains $x_i \in I_i$ et $g_i \in B$, $g_i \not \in \mathfrak q$.
En chassant les dénominateurs, on trouve donc une relation de la forme
$$
x_i + \sum\nolimits_j f_{i, j}c_{i, j} = g_i
$$
avec $x_i \in I_i$, $f_{i, j} \in \mathfrak q$, $c_{i, j} \in C_i \otimes_A B$
et $g_i \in B$, $g_i \not \in \mathfrak q$. En remplaçant $B$ par
$B_{g_1 \ldots g_n}$, c'est-à-dire $T$ par un voisinage affine plus petit
de $t$, on peut supposer les équations de la forme
$$
x_i + \sum\nolimits_j f_{i, j}c_{i, j} = 1
$$
avec $x_i \in I_i$, $f_{i, j} \in \mathfrak q$, $c_{i, j} \in C_i \otimes_A B$.

\medskip\noindent
Pour terminer, écrivons $B$ comme limite inductive de
$A$-algèbres de présentation finie $B_\lambda$ indexée par un ensemble ordonné filtrant $\Lambda$.
Pour chaque $\lambda$, posons
$\mathfrak q_\lambda = (B_\lambda \to B)^{-1}(\mathfrak q)$.
Pour $\lambda \in \Lambda$ assez grand, on peut trouver
\begin{enumerate}
\item un élément
$x_{i, \lambda} \in C_i \otimes_A B_\lambda$ d'image $x_i$,
\item des éléments $f_{i, j, \lambda} \in \mathfrak q_{i, \lambda}$
d'images $f_{i, j}$,
\item des éléments $c_{i, j, \lambda} \in C_i \otimes_A B_\lambda$
d'images $c_{i, j}$.
\end{enumerate}
En augmentant encore $\lambda$, l'équation
$$
x_{i, \lambda} + \sum\nolimits_j f_{i, j, \lambda}c_{i, j, \lambda} = 1
$$
sera satisfaite. Fixons un tel $\lambda$ et posons $T' = \Spec(B_\lambda)$.
Alors $t' \in T'$ est le point associé à l'idéal premier $\mathfrak q_\lambda$.
Enfin, soit $Z' \subset X_{T'}$ l'image schématique de
$Z \to X_T \to X_{T'}$. Puisque $X_T \to X_{T'}$ est affine, on peut calculer $Z'$
sur les ouverts affines $U_{i, T'}$ comme le sous-schéma fermé associé
à $\Ker(C_i \otimes_A B_\lambda \to C_i \otimes_A B/I_i)$ ; voir
Morphismes, Exemple \ref{morphisms-example-scheme-theoretic-image}.
Ainsi $x_{i, \lambda}$ appartient à l'idéal définissant $Z'$. La dernière
équation affichée montre donc que $Z' \cap X_{t'}$ est vide.
\end{proof}

\begin{lemma}
\label{lemma-test-universally-closed}
Soit $f : X \to S$ un morphisme quasi-compact de schémas.
Les conditions suivantes sont équivalentes :
\begin{enumerate}
\item $f$ est universellement fermé,
\item pour tout morphisme $S' \to S$ localement de présentation finie,
le changement de base $X_{S'} \to S'$ est fermé,
\item pour tout $n$, le morphisme
$\mathbf{A}^n \times X \to \mathbf{A}^n \times S$
est fermé.
\end{enumerate}
\end{lemma}

\begin{proof}
Il est clair que (1) entraîne (2). Montrons que (2) entraîne (1).
Supposons que le changement de base $X_T \to T$ ne soit pas fermé pour un certain
schéma $T$ sur $S$. D'après
Schémas, Lemme \ref{schemes-lemma-quasi-compact-closed},
il existe alors une spécialisation $t_1 \leadsto t$ dans
$T$ et un point $\xi \in X_T$ d'image $t_1$ tel que $\xi$ ne
se spécialise en aucun point de la fibre au-dessus de $t$. Posons
$Z = \overline{\{\xi\}} \subset X_T$. Alors $Z \cap X_t = \emptyset$. Appliquons le
Lemme \ref{lemma-separate}.
On trouve un voisinage ouvert $V \subset T$ de $t$, un diagramme commutatif
$$
\xymatrix{
V \ar[d] \ar[r]_a & T' \ar[d]^b \\
T \ar[r]^g & S,
}
$$
et un sous-schéma fermé $Z' \subset X_{T'}$ tels que
\begin{enumerate}
\item le morphisme $b : T' \to S$ soit localement de présentation finie,
\item en posant $t' = a(t)$, on ait $Z' \cap X_{t'} = \emptyset$,
\item $Z \cap X_V$ ait son image dans $Z'$ par le morphisme $X_V \to X_{T'}$.
\end{enumerate}
Cela signifie que $X_{T'} \to T'$ envoie le fermé $Z'$
sur une partie de $T'$ contenant $a(t_1)$ mais non $t' = a(t)$.
Puisque $a(t_1) \leadsto a(t) = t'$, on conclut que $X_{T'} \to T'$
n'est pas fermé. Ainsi, si $X \to S$ n'est pas universellement fermé,
$X_{T'} \to T'$ n'est pas fermé pour un certain $T' \to S$
localement de présentation finie. Autrement dit, (2)
entraîne (1).

\medskip\noindent
Supposons que $\mathbf{A}^n \times X \to \mathbf{A}^n \times S$ soit
fermé pour tout entier $n$. Nous voulons montrer que $X_T \to T$ est
fermé pour tout schéma $T$ localement de présentation finie
sur $S$. On peut bien sûr supposer $T$ affine et son image contenue dans
un ouvert affine $V$ de $S$ (car le fait que $X_T \to T$ soit fermé est local sur $T$).
Dans ce cas, il existe une immersion fermée $T \to \mathbf{A}^n \times V$,
car $\mathcal{O}_T(T)$ est une
$\mathcal{O}_S(V)$-algèbre de présentation finie ; voir
Morphismes,
Lemme \ref{morphisms-lemma-locally-finite-presentation-characterize}.
Alors $T \to \mathbf{A}^n \times S$ est une immersion localement fermée.
On obtient donc un diagramme cartésien
$$
\xymatrix{
X_T \ar[d]_{f_T} \ar[r] & \mathbf{A}^n \times X \ar[d]^{f_n} \\
T \ar[r] & \mathbf{A}^n \times S
}
$$
de schémas dont les flèches horizontales sont des immersions localement fermées.
Toute partie fermée $Z \subset X_T$ peut donc s'écrire
$X_T \cap Z'$ pour une partie fermée $Z' \subset \mathbf{A}^n \times X$.
Alors $f_T(Z) = T \cap f_n(Z')$, et si $f_n$ est fermé,
$f_T$ l'est aussi.
\end{proof}

\begin{lemma}
\label{lemma-limited-base-change}
Soit $S$ un schéma.
Soit $f : X \to S$ un morphisme séparé de type fini.
Les conditions suivantes sont équivalentes :
\begin{enumerate}
\item Le morphisme $f$ est propre.
\item Pour tout morphisme $S' \to S$ localement de type fini,
le changement de base $X_{S'} \to S'$ est fermé.
\item Pour tout $n \geq 0$, le morphisme
$\mathbf{A}^n \times X \to \mathbf{A}^n \times S$ est fermé.
\end{enumerate}
\end{lemma}

\begin{proof}[Première démonstration]
Puisqu'un morphisme propre est exactement un morphisme
séparé, de type fini et universellement fermé, ce
lemme est un cas particulier du Lemme \ref{lemma-test-universally-closed}.
\end{proof}

\begin{proof}[Seconde démonstration]
Il est clair que (1) entraîne (2) et que (2) entraîne (3) ; il suffit donc de montrer que (3)
entraîne (1).
Ramenons-nous d'abord au cas où $S$ est affine. Supposons que (3) entraîne (1)
lorsque la base est affine. Soit maintenant $f: X \to S$ un morphisme séparé de
type fini. La propreté est locale sur la base
(voir Morphismes, Lemme \ref{morphisms-lemma-proper-local-on-the-base}) ; si
$S = \bigcup_\alpha S_\alpha$ est un recouvrement ouvert affine et si
l'on note $X_\alpha := f^{-1}(S_\alpha)$, il
suffit donc de montrer que $f|_{X_\alpha}: X_\alpha \to S_\alpha$ est propre
pour tout $\alpha$. Puisque $S_\alpha$ est affine, si
l'application $f|_{X_\alpha}$ satisfait (3), elle satisfera (1)
par hypothèse, donc sera propre. Pour achever la réduction au
cas où $S$ est affine, il faut montrer que, si $f: X \to S$ est séparé de
type fini et satisfait (3), alors $f|_{X_\alpha} : X_\alpha \to S_\alpha$
est séparé de type fini et satisfait (3). Le caractère séparé et le type
fini sont clairs. Pour (3), remarquons que
$\mathbf{A}^n \times X_\alpha$ est l'image réciproque ouverte de
$\mathbf{A}^n \times S_\alpha$ par l'application $1 \times f$. Fixons un
fermé $Z \subset \mathbf A^n \times X_\alpha$. Notons $\bar Z$
l'adhérence de $Z$ dans $\mathbf{A}^n \times X$. Pour des raisons
topologiques, on a
$$
1 \times f(\bar Z) \cap \mathbf{A}^n \times S_\alpha  = 1 \times f(Z).
$$
Ainsi $1 \times f(Z)$ est fermé, et la démonstration de
(3) $\Rightarrow$ (1) est ramenée au cas affine.

\medskip\noindent
Supposons $S$ affine et $f : X \to S$ séparé de type fini.
On peut appliquer le Lemme de Chow \ref{lemma-chow-finite-type}
pour obtenir $\pi : X' \to X$ propre surjectif et $X' \to \mathbf{P}^n_S$
une immersion. Si $X$ est propre sur $S$, alors $X' \to S$ est propre
(Morphismes, Lemme \ref{morphisms-lemma-composition-proper}). Puisque
$\mathbf{P}^n_S \to S$ est séparé, on conclut que $X' \to
\mathbf{P}^n_S$ est propre
(Morphismes, Lemme \ref{morphisms-lemma-image-proper-scheme-closed}),
donc est une immersion fermée
(Schémas, Lemme \ref{schemes-lemma-immersion-when-closed}).
Réciproquement, supposons que $X' \to \mathbf{P}^n_S$ soit une immersion fermée.
Considérons le diagramme :
\begin{equation}
\label{equation-check-proper}
\xymatrix{
X' \ar[r] \ar@{->>}[d]_{\pi} &
\mathbf{P}^n_S \ar[d] \\
X \ar[r]^f & S
}
\end{equation}
Tous les morphismes sont a priori propres, sauf $X \to S$.
On en conclut que $X \to S$ est propre par
Morphismes, Lemme \ref{morphisms-lemma-image-proper-is-proper}.
Nous avons donc montré que $X \to S$ est propre si et seulement si
$X' \to \mathbf{P}^n_S$ est une immersion fermée.

\medskip\noindent
Supposons $S$ affine et (3) satisfaite, et soient $n, X', \pi$ comme ci-dessus.
Puisque le caractère fermé d'un morphisme est local sur la base, l'application
$X \times \mathbf{P}^n \to S \times \mathbf{P}^n$ est fermée : d'après (3),
$X \times \mathbf{A}^n \to S \times \mathbf{A}^n$ est fermée, et
l'espace projectif est recouvert par des copies de l'espace affine de dimension $n$ ; voir
Constructions,
Lemme \ref{constructions-lemma-standard-covering-projective-space}.
D'après Morphismes, Lemme \ref{morphisms-lemma-base-change-proper},
le morphisme
$$
X' \times_S \mathbf{P}^n_S
\to
X \times_S \mathbf{P}^n_S =
X \times \mathbf{P}^n
$$
est propre. Puisque $\mathbf{P}^n$ est séparé, la projection
$$
X' \times_S \mathbf{P}^n_S = \mathbf{P}^n_{X'} \to X'
$$
est séparée, étant simplement un changement de base d'un morphisme
séparé. Par conséquent, l'application $X' \to X' \times_S \mathbf{P}^n_S$ est propre,
car elle est une section d'un morphisme séparé (voir
Schémas, Lemme \ref{schemes-lemma-section-immersion}).
En composant ces morphismes
$$
X' \to X' \times_S \mathbf{P}^n_S \to X \times_S \mathbf{P}^n_S
= X \times \mathbf{P}^n \to S \times \mathbf{P}^n = \mathbf{P}^n_S
$$
on trouve que l'immersion $X' \to \mathbf{P}^n_S$ est fermée, donc
est une immersion fermée.
\end{proof}





\section{Critère valuatif noethérien}
\label{section-Noetherian-valuative-criterion}

\noindent
Si la base est noethérienne, on peut établir le critère valuatif
en ne considérant que des anneaux de valuation discrète.

\medskip\noindent
Beaucoup des résultats de cette section peuvent (et devraient peut-être)
se démontrer à l'aide du lemme suivant, bien que nous ne l'ayons
pas toujours fait.

\begin{lemma}
\label{lemma-reach-point-closure-Noetherian}
Soit $f : X \to Y$ un morphisme de schémas.
Supposons $f$ de type fini et $Y$ localement noethérien.
Soit $y \in Y$ un point de l'adhérence de l'image de $f$.
Il existe alors un diagramme commutatif
$$
\xymatrix{
\Spec(K) \ar[r] \ar[d] & X \ar[d]^f \\
\Spec(A) \ar[r] & Y
}
$$
où $A$ est un anneau de valuation discrète et $K$ son corps des fractions,
envoyant le point fermé de $\Spec(A)$ sur $y$. De plus, on peut supposer
que le point image de $\Spec(K) \to X$ est un point générique $\eta$
d'une composante irréductible de $X$ et que $K = \kappa(\eta)$.
\end{lemma}

\begin{proof}
D'après la version non noethérienne de ce lemme
(Morphismes, Lemme \ref{morphisms-lemma-reach-points-scheme-theoretic-image}),
il existe un point $x \in X$ tel que $f(x)$ se spécialise en $y$.
On peut remplacer $x$ par tout point se spécialisant en $x$ ; on peut donc
supposer que $x$ est un point générique d'une composante irréductible de $X$.
On obtient ainsi un homomorphisme $\mathcal{O}_{Y, y} \to \kappa(x)$
(voir Schémas, section \ref{schemes-section-points}).
Soit $R \subset \kappa(x)$ son image. Alors $R$ est noethérien comme quotient
de l'anneau local noethérien $\mathcal{O}_{Y, y}$.
D'autre part, $\kappa(x)$ est
une extension de type fini du corps des fractions de $R$,
puisque $f$ est de type fini.
Il existe donc un anneau de valuation discrète $A \subset \kappa(x)$
de corps des fractions $\kappa(x)$ dominant $R$, d'après
Algèbre, Lemme \ref{algebra-lemma-exists-dvr}. Alors
$$
\xymatrix{
\Spec(\kappa(x)) \ar[d] \ar[rrr] & & & X \ar[d] \\
\Spec(A) \ar[r] & \Spec(R) \ar[r] & \Spec(\mathcal{O}_{Y, y}) \ar[r] & Y
}
$$
est le diagramme cherché.
\end{proof}

\noindent
Énonçons d'abord le résultat concernant
le caractère séparé. Nous utiliserons souvent des diagrammes commutatifs de morphismes
de schémas, dont les flèches pleines sont données, de la forme suivante :
\begin{equation}
\label{equation-valuative}
\vcenter{
\xymatrix{
\Spec(K) \ar[r] \ar[d] & X \ar[d] \\
\Spec(A) \ar[r] \ar@{-->}[ru] & S
}
}
\end{equation}
où $A$ est un anneau de valuation et $K$ son corps des fractions.

\begin{lemma}
\label{lemma-Noetherian-dvr-valuative-separation}
Soit $S$ un schéma localement noethérien.
Soit $f : X \to S$ un morphisme de schémas.
Supposons $f$ localement de type fini.
Les conditions suivantes sont équivalentes :
\begin{enumerate}
\item Le morphisme $f$ est séparé.
\item Pour tout diagramme (\ref{equation-valuative}), il existe au plus
une flèche pointillée.
\item Pour tout diagramme (\ref{equation-valuative}) où $A$ est un anneau
de valuation discrète, il existe au plus une flèche pointillée.
\item Pour toute composante irréductible $X_0$ de $X$ de
point générique $\eta \in X_0$, tout anneau de valuation discrète
$A \subset K = \kappa(\eta)$ de corps des fractions $K$ et tout
diagramme (\ref{equation-valuative}) tel que
le morphisme $\Spec(K) \to X$ soit le morphisme canonique
(voir Schémas, section \ref{schemes-section-points}),
il existe au plus une flèche pointillée.
\end{enumerate}
\end{lemma}

\begin{proof}
Il est clair que (1) entraîne (2), que (2) entraîne (3) et que (3) entraîne (4).
Il reste à montrer que (4) entraîne (1). Supposons (4).
Ramenons-nous d'abord au cas où $S$ est affine. Le caractère séparé est local
sur la base (voir
Schémas, Lemme \ref{schemes-lemma-characterize-separated}).
Il suffit donc de montrer que, lorsque
$X \to S$ satisfait (4), la restriction $X_\alpha \to S_\alpha$ satisfait (4),
où $S_\alpha \subset S$ est un ouvert (affine) et $X_\alpha :=
f^{-1}(S_\alpha)$. Les
points génériques des composantes irréductibles de $X_\alpha$ sont des
points génériques de composantes irréductibles de $X$, puisque $X_\alpha$ est
ouvert dans $X$. Deux flèches pointillées distinctes dans le diagramme
\begin{equation}
\label{equation-valuative-alpha}
\xymatrix{
\Spec(K) \ar[r] \ar[d] & X_\alpha \ar[d] \\
\Spec(A) \ar[r] \ar@{-->}[ru] & S_\alpha
}
\end{equation}
donneraient donc deux flèches distinctes dans le diagramme
(\ref{equation-valuative}), par les morphismes $X_\alpha \to X$ et
$S_\alpha \to S$, ce qui est contradictoire. On est ainsi ramené
au cas où $S$ est affine. Remarquons qu'au cours de cette
réduction, nous montrons que, si $X \to S$ satisfait (4), la restriction $U
\to V$ satisfait (4) pour des ouverts $U \subset X$ et $V \subset S$ tels que
$f(U) \subset V$.

\medskip\noindent
Ramenons-nous ensuite au cas où $X \to S$ est de type fini. Supposons
que (4) entraîne (1) lorsque $X$ est de type fini. Puisque
$S$ est noethérien et $X$ localement de type fini sur $S$,
$X$ est lui aussi localement noethérien (voir Morphismes,
Lemme \ref{morphisms-lemma-finite-type-noetherian}).
Ainsi $X \to S$ est quasi-séparé (voir
Propriétés, Lemme \ref{properties-lemma-locally-Noetherian-quasi-separated}),
et l'on peut donc appliquer le critère valuatif pour déterminer si $X$
est séparé (voir
Schémas, Lemme \ref{schemes-lemma-valuative-criterion-separatedness}).
Soit $X = \bigcup_\alpha X_\alpha$ un recouvrement ouvert
affine de $X$. Étant données deux flèches pointillées dans un diagramme
(\ref{equation-valuative}), les images du point fermé de
$\Spec A$
appartiennent à deux ouverts $X_\alpha$ et $X_\beta$. Puisque $X_\alpha \cup
X_\beta$ est ouvert, il contient, pour des raisons topologiques, l'image de
$\Spec(A)$ par chacune des deux applications. Les deux flèches pointillées se factorisent donc
par $X_\alpha \cup X_\beta \to X$, dont la source est un schéma de type fini sur
$S$. Puisque $X_\alpha \cup X_\beta$ est un ouvert de $X$, notre
remarque précédente montre que $X_\alpha \cup X_\beta$ satisfait (4), donc
est séparé par hypothèse. Les deux flèches pointillées
sont donc égales. On est ainsi ramené au cas où $X \to S$ est de type fini.

\medskip\noindent
Supposons $X \to S$ de type fini et (4) satisfaite.
Puisque $X \to S$ est de type fini et $S$ un schéma affine
noethérien, $X$ est lui aussi noethérien (voir
Morphismes, Lemme \ref{morphisms-lemma-finite-type-noetherian}).
Ainsi $X \to X \times_S X$ est
une immersion quasi-compacte de schémas noethériens. Raisonnons par
l'absurde. Supposons que $X \to X \times_S X$ ne soit pas fermé.
Il existe alors un $y \in X \times_S X$ dans l'adhérence de l'image, mais
hors de l'image. Comme $X$ est noethérien, il possède un nombre fini de composantes
irréductibles. Le point $y$ appartient donc à l'adhérence de l'image d'une
composante irréductible $X_0 \subset X$. Munissons $X_0$ de la structure
réduite induite. La composée $X_0 \to X \to X \times_S X$
se factorise par le sous-schéma fermé $X_0 \times_S X_0 \subset X \times_S X$.
Notons l'adhérence de $\Delta(X_0)$ dans $X_0 \times_S X_0$
par $\bar X_0$ (encore munie de la structure de sous-schéma fermé réduit). Ainsi $y \in \bar X_0$.
Puisque $X_0 \to X_0 \times_S X_0$ est une immersion, l'image de $X_0$
est ouverte dans $\bar X_0$. Ainsi $X_0$ et $\bar X_0$ sont
birationnels. Puisque $\bar{X}_0$ est un sous-schéma fermé d'un
schéma noethérien, il est noethérien. L'anneau local
$\mathcal O_{{\bar X_0, y}}$ est donc un anneau local noethérien intègre de corps des
fractions $K$ égal au corps des fonctions de $X_0$. D'après le théorème de Krull--Akizuki
(voir Algèbre, Lemme \ref{algebra-lemma-exists-dvr}), il existe un
anneau de valuation discrète $A$ dominant $\mathcal O_{{\bar X_0, y}}$
et de corps des fractions $K$. Cela permet de construire un diagramme :
\begin{equation}
\label{equation-valuative-generic}
\xymatrix{
\Spec(K) \ar[r] \ar[d] & X_0 \ar[d]^{\Delta} \\
\Spec(A) \ar[r] \ar@{-->}[ur]& X_0 \times_S X_0 \\
}
\end{equation}
qui envoie $\Spec K$ sur le point générique de $\Delta(X_0)$ et
le point fermé de $A$ sur $y \in X_0 \times_S X_0$ (utiliser
Schémas, section \ref{schemes-section-points}, pour construire les flèches).
Il ne peut même pas exister de
flèche pointillée ensembliste, puisque $y$ n'est pas dans l'image de
$\Delta$, par le choix de $y$. Par un argument catégorique, l'existence
de la flèche pointillée dans le diagramme ci-dessus équivaut à l'unicité
de la flèche pointillée dans le diagramme suivant :
\begin{equation}
\label{equation-valuative-nonexistent}
\xymatrix{
\Spec(K) \ar[r] \ar[d] & X_0 \ar[d]\\
\Spec(A) \ar[r] \ar@{-->}[ur] & S \\
}
\end{equation}
La non-existence dans le premier diagramme donne donc la non-unicité dans
ce dernier. Ainsi $X_0$ ne satisfait pas
l'unicité pour les anneaux de valuation discrète ; puisque $X_0$ est une
composante irréductible de $X$, le morphisme $X \to S$ ne satisfait pas
(4). Nous avons donc montré que (4) entraîne (1).
\end{proof}

\begin{lemma}
\label{lemma-Noetherian-dvr-valuative-proper}
Soit $S$ un schéma localement noethérien.
Soit $f : X \to S$ un morphisme de type fini.
Les conditions suivantes sont équivalentes :
\begin{enumerate}
\item Le morphisme $f$ est propre.
\item Pour tout diagramme (\ref{equation-valuative}), il existe exactement
une flèche pointillée.
\item Pour tout diagramme (\ref{equation-valuative}) où $A$ est un anneau
de valuation discrète, il existe exactement une flèche pointillée.
\item Pour toute composante irréductible $X_0$ de $X$ de
point générique $\eta \in X_0$, tout anneau de valuation discrète
$A \subset K = \kappa(\eta)$ de corps des fractions $K$ et tout
diagramme (\ref{equation-valuative}) tel que
le morphisme $\Spec(K) \to X$ soit le morphisme canonique
(voir
Schémas, section \ref{schemes-section-points}),
il existe exactement une flèche pointillée.
\end{enumerate}
\end{lemma}

\begin{proof}
(1) entraîne (2), qui entraîne (3), qui entraîne (4). Montrons maintenant que (4) entraîne
(1). Comme dans la démonstration du Lemme \ref{lemma-Noetherian-dvr-valuative-separation},
on se ramène au
cas où $S$ est affine : la propreté est locale sur la base, et si $X
\to S$ satisfait (4), alors $X_\alpha \to S_\alpha$ la satisfait aussi pour
un ouvert $S_\alpha \subset S$ et $X_\alpha = f^{-1}(S_\alpha)$.

\medskip\noindent
Le schéma $S$ est maintenant noethérien ; $X$ l'est donc aussi, puisque $X \to
S$ est de type fini. On peut alors utiliser le lemme de Chow
(Cohomologie des schémas, Lemme \ref{coherent-lemma-chow-Noetherian})
pour obtenir un morphisme surjectif, propre et birationnel
$X' \to X$ et une immersion $X' \to \mathbf{P}^n_S$. Nous voulons
montrer que $X \to S$ est universellement fermé. Comme dans la démonstration du Lemme
\ref{lemma-limited-base-change}, il suffit de vérifier que
$X' \to \mathbf{P}^n_S$ est une immersion fermée.
Raisonnons par l'absurde et supposons que $X' \to
\mathbf{P}^n_S$ ne soit pas une immersion fermée. Il existe alors un $y
\in \mathbf{P}^n_S$ dans l'adhérence de l'image de $X'$,
mais hors de l'image. Ainsi $y$ est dans l'adhérence de l'image d'une
composante irréductible $X_0'$ de $X'$, mais hors de l'image.
Soit $\bar X_0' \subset \mathbf{P}^n_S$ l'adhérence de
l'image de $X_0'$. Puisque $X' \to \mathbf{P}^n_S$ est une immersion
de schémas noethériens, le morphisme $X'_0 \to \bar X_0'$ est
ouvert et dense. D'après
Algèbre, Lemme \ref{algebra-lemma-exists-dvr},
ou
Propriétés, Lemme \ref{properties-lemma-locally-Noetherian-specialization-dvr},
on peut trouver un anneau de valuation discrète $A$ dominant
$\mathcal{O}_{\bar X_0', y}$ et de même corps
des fractions $K$. Il est clair que
$K$ est le corps résiduel du point générique de $X_0'$.
On obtient donc le diagramme commutatif de flèches pleines
\begin{equation}
\label{equation-solid}
\xymatrix{
\Spec K \ar[r] \ar[d] & X' \ar [r] \ar[d] &
\mathbf{P}^n_S \ar[d] \\
\Spec A \ar@{-->}[r] \ar@{-->}[ru] \ar[urr] & X \ar[r] & S\\
}
\end{equation}
Remarquons que le point fermé de $A$ a pour image $y \in \mathbf{P}^n_S$. Par
construction, il n'existe pas de relèvement ensembliste à $X'$.
Puisque $X' \to X$ est birationnel, l'image de $X'_0$ dans $X$ est une
composante irréductible $X_0$ de $X$, et $K$ s'identifie aussi
au corps des fonctions de $X_0$. Puisque $X \to S$ satisfait (4) par hypothèse,
la flèche pointillée $\Spec(A) \to X$ existe.
Puisque $X' \to X$ est propre, cette flèche pointillée
se relève en une flèche pointillée $\Spec(A) \to X'$ (utiliser Schémas,
Proposition \ref{schemes-proposition-characterize-universally-closed}).
En la composant avec l'immersion $X' \to \mathbf{P}^n_S$, on obtient
un autre morphisme (non représenté dans le diagramme)
$\Spec(A) \to \mathbf{P}^n_S$. Puisque $\mathbf{P}^n_S$
est propre sur $S$, il satisfait (2), et ces deux morphismes
coïncident. C'est une contradiction, car nous avons construit
le relèvement interdit de l'application initiale $\Spec(A) \to \mathbf{P}^n_S$
à $X'$.
\end{proof}

\begin{lemma}
\label{lemma-check-universally-closed-Noetherian}
Soit $f : X \to S$ un morphisme de type fini de schémas.
Supposons $S$ localement noethérien. Les conditions suivantes sont alors équivalentes :
\begin{enumerate}
\item $f$ est universellement fermé,
\item pour tout $n$, le morphisme
$\mathbf{A}^n \times X \to \mathbf{A}^n \times S$ est fermé,
\item pour tout diagramme (\ref{equation-valuative}), il existe une
flèche pointillée,
\item pour tout diagramme (\ref{equation-valuative}) où $A$ est un anneau
de valuation discrète, il existe une flèche pointillée.
\end{enumerate}
\end{lemma}

\begin{proof}
L'équivalence de (1) et (2) est un cas particulier du
Lemme \ref{lemma-test-universally-closed}.
L'équivalence de (1) et (3) est un cas particulier de
Schémas, Proposition \ref{schemes-proposition-characterize-universally-closed}.
Il est évident que (3) entraîne (4).
Il suffit donc de montrer que (4) entraîne (2).
Nous montrerons que $\mathbf{A}^n \times X \to \mathbf{A}^n \times S$
est fermé à l'aide du critère de
Schémas, Lemme \ref{schemes-lemma-quasi-compact-closed}.
Choisissons $n$, une spécialisation non triviale $z \leadsto z'$ de points
de $\mathbf{A}^n \times S$ et un point $y \in \mathbf{A}^n \times X$
au-dessus de $z$. Remarquons que $\kappa(y)$ est une extension de corps de type fini
de $\kappa(z)$, puisque $\mathbf{A}^n \times X \to \mathbf{A}^n \times S$
est de type fini. D'après
Propriétés, Lemme \ref{properties-lemma-locally-Noetherian-specialization-dvr},
ou
Algèbre, Lemme \ref{algebra-lemma-exists-dvr},
il existe donc un anneau de valuation discrète $A \subset \kappa(y)$
de corps des fractions $\kappa(z)$ dominant l'image de
$\mathcal{O}_{\mathbf{A}^n \times S, z'}$ dans $\kappa(z)$.
Cela donne un diagramme commutatif
$$
\xymatrix{
\Spec(\kappa(y)) \ar[r] \ar[d] &
\mathbf{A}^n \times X \ar[d] \ar[r] & X \ar[d] \\
\Spec(A) \ar[r] & \mathbf{A}^n \times S \ar[r] & S
}
$$
La propriété (4) entraîne alors l'existence d'un morphisme
$\Spec(A) \to X$ s'insérant dans ce diagramme.
Puisque la flèche horizontale inférieure gauche fournit déjà
le morphisme $\Spec(A) \to \mathbf{A}^n$,
on obtient aussi un morphisme
$\Spec(A) \to \mathbf{A}^n \times X$ s'insérant dans le
carré de gauche. L'image $y' \in \mathbf{A}^n \times X$
du point fermé est donc une spécialisation de $y$ au-dessus de $z'$.
Cela montre que les spécialisations se relèvent le long de
$\mathbf{A}^n \times X \to \mathbf{A}^n \times S$,
ce qui conclut.
\end{proof}









\section{Critères valuatifs noethériens affinés}
\label{section-refined-valuative-criteria}

\noindent
Il n'est généralement pas nécessaire de considérer tous les diagrammes possibles
avec des anneaux de valuation pour vérifier les critères valuatifs. Un exemple
est donné par Morphismes, Lemme
\ref{morphisms-lemma-refined-valuative-criterion-universally-closed}.
Dans le cadre noethérien, nous l'avons également vu dans les
Lemmes \ref{lemma-Noetherian-dvr-valuative-separation} et
\ref{lemma-Noetherian-dvr-valuative-proper}. Voici une autre variante.

\begin{lemma}
\label{lemma-refined-valuative-criterion-proper}
Soient $f : X \to S$ et $h : U \to X$ des morphismes de schémas.
Supposons $S$ localement noethérien, $f$ et $h$ de type fini,
$f$ séparé et $h(U)$ dense dans $X$.
Si, pour tout diagramme commutatif de flèches pleines
$$
\xymatrix{
\Spec(K) \ar[r] \ar[d] & U \ar[r]^h & X \ar[d]^f \\
\Spec(A) \ar[rr] \ar@{-->}[rru] & & S
}
$$
où $A$ est un anneau de valuation discrète de corps des fractions $K$,
il existe une flèche pointillée rendant le diagramme commutatif, alors $f$ est propre.
\end{lemma}

\begin{proof}
On se ramène immédiatement au cas où $S$ est affine.
Alors $U$ est quasi-compact.
Soit $U = U_1 \cup \ldots \cup U_n$ un recouvrement ouvert affine.
On peut remplacer $U$ par $U_1 \amalg \ldots \amalg U_n$ sans
changer les hypothèses ; on peut donc supposer $U$ affine.
On peut alors trouver une immersion ouverte $U \to Y$ sur $X$,
où $Y$ est propre sur $X$. (Plonger d'abord $U$ dans $\mathbf{A}^n_X$
à l'aide de Morphismes, Lemme \ref{morphisms-lemma-quasi-affine-finite-type-over-S},
puis prendre l'adhérence dans $\mathbf{P}^n_X$, ou utiliser directement
Morphismes, Lemme \ref{morphisms-lemma-quasi-projective-open-projective}.)
On peut supposer $U$ dense dans $Y$ (remplacer $Y$ par l'adhérence schématique
de $U$ si nécessaire ; voir Morphismes, section
\ref{morphisms-section-scheme-theoretic-closure}).
Remarquons que $g : Y \to X$ est surjectif, car son image est fermée
et contient la partie dense $h(U)$.
Nous allons montrer que $Y \to S$ est propre. Cela entraînera que
$X \to S$ est propre par
Morphismes, Lemme \ref{morphisms-lemma-image-proper-is-proper},
et achèvera la démonstration.
Pour montrer que $Y \to S$ est propre, nous utiliserons
le point (4) du Lemme \ref{lemma-Noetherian-dvr-valuative-proper}.
Considérons à cet effet un diagramme
$$
\xymatrix{
\Spec(K) \ar[r]_y \ar[d] & Y \ar[d]^{f \circ g} \\
\Spec(A) \ar[r] \ar@{..>}[ru] & S
}
$$
où $A$ est un anneau de valuation discrète de corps des fractions $K$
et où $y : \Spec(K) \to Y$ est l'inclusion d'un point générique.
Il faut montrer qu'il existe une unique flèche pointillée.
L'unicité résulte de la réciproque du critère valuatif
de séparation
(Schémas, Lemme \ref{schemes-lemma-separated-implies-valuative}),
puisque $Y \to S$ est séparé comme
composée des morphismes séparés $Y \to X$ et $X \to S$
(Schémas, Lemme \ref{schemes-lemma-separated-permanence}).
L'existence se voit comme suit.
Puisque $y$ est un point générique de $Y$, il appartient à $U$.
L'hypothèse du lemme
fournit un morphisme $a : \Spec(A) \to X$ tel que
$$
\xymatrix{
\Spec(K) \ar[r]_y \ar[d] & U \ar[r] & X \ar[d]^f \\
\Spec(A) \ar[rr] \ar[rru]^a & & S
}
$$
soit commutatif. Puisque $Y \to X$ est propre, on peut ensuite
appliquer le critère valuatif de propreté
(Morphismes, Lemme \ref{morphisms-lemma-characterize-proper})
pour trouver un morphisme $b : \Spec(A) \to Y$ tel que
$$
\xymatrix{
\Spec(K) \ar[r]_y \ar[d] & Y \ar[d]^g \\
\Spec(A) \ar[r]^a \ar[ru]^b & X
}
$$
soit commutatif. Cela achève la démonstration, puisque
$b$ peut servir de flèche pointillée ci-dessus.
\end{proof}

\begin{lemma}
\label{lemma-refined-valuative-criterion-separated}
Soient $f : X \to S$ et $h : U \to X$ des morphismes de schémas.
Supposons $S$ localement noethérien, $f$ localement de type fini,
$h$ de type fini et $h(U)$ dense dans $X$.
Si, pour tout diagramme commutatif de flèches pleines
$$
\xymatrix{
\Spec(K) \ar[r] \ar[d] & U \ar[r]^h & X \ar[d]^f \\
\Spec(A) \ar[rr] \ar@{-->}[rru] & & S
}
$$
où $A$ est un anneau de valuation discrète de corps des fractions $K$,
il existe au plus une flèche pointillée rendant le diagramme commutatif, alors $f$ est
séparé.
\end{lemma}

\begin{proof}
Nous allons appliquer le Lemme \ref{lemma-refined-valuative-criterion-proper}
aux morphismes $U \to X$ et $\Delta : X \to X \times_S X$.
Vérifions les conditions. Remarquons que $\Delta$ est quasi-compact d'après
Propriétés, Lemme \ref{properties-lemma-locally-Noetherian-quasi-separated}
(et Schémas, Lemme \ref{schemes-lemma-compose-after-separated}).
Bien entendu, $\Delta$ est localement de type fini et séparé (ce qui est vrai
de tout morphisme diagonal).
Enfin, donnons-nous un diagramme commutatif de flèches pleines
$$
\xymatrix{
\Spec(K) \ar[r] \ar[d] & U \ar[r]^h & X \ar[d]^\Delta \\
\Spec(A) \ar[rr]^{(a, b)} \ar@{-->}[rru] & & X \times_S X
}
$$
où $A$ est un anneau de valuation discrète de corps des fractions $K$.
Alors $a$ et $b$ donnent deux flèches pointillées dans le diagramme du lemme
et doivent être égales. On peut donc prendre $a = b$ comme flèche pointillée,
ce qui donne l'existence et achève la démonstration.
\end{proof}

\begin{lemma}
\label{lemma-refined-valuative-criterion-universally-closed}
Soient $f : X \to S$ et $h : U \to X$ des morphismes de schémas.
Supposons $S$ localement noethérien, $f$ et $h$ de type fini et
$h(U)$ dense dans $X$. Si, pour tout diagramme commutatif de flèches pleines
$$
\xymatrix{
\Spec(K) \ar[r] \ar[d] & U \ar[r]^h & X \ar[d]^f \\
\Spec(A) \ar[rr] \ar@{-->}[rru] & & S
}
$$
où $A$ est un anneau de valuation discrète de corps des fractions $K$,
il existe une unique flèche pointillée rendant le diagramme commutatif, alors $f$ est propre.
\end{lemma}

\begin{proof}
Combiner les Lemmes \ref{lemma-refined-valuative-criterion-separated} et
\ref{lemma-refined-valuative-criterion-proper}.
\end{proof}









\section{Critères valuatifs sur une base de Nagata}
\label{section-nagata-valuative}

\noindent
Pour les schémas localement de type fini sur une
base de Nagata, on peut se limiter aux anneaux de valuation discrète
essentiellement de type fini sur la base.
Les résultats suivants ne sont que quelques exemples de ce que l'on peut obtenir.

\begin{lemma}
\label{lemma-essentially-finite-type-criterion-universally-closed}
Soit $S$ un schéma de Nagata (donc, en particulier, localement noethérien).
Soit $f : X \to Y$ un morphisme quasi-compact de schémas localement
de type fini sur $S$. Les conditions suivantes sont équivalentes :
\begin{enumerate}
\item $f$ est universellement fermé,
\item pour tout $n$, le morphisme
$\mathbf{A}^n \times X \to \mathbf{A}^n \times Y$ est fermé,
\item pour tout diagramme commutatif
$$
\xymatrix{
U \ar[r] \ar[d] & X \ar[d]^f \\
C \ar[r] \ar@{..>}[ru] & Y
}
$$
de schémas sur $S$ tel que
\begin{enumerate}
\item $C$ soit un schéma intègre normal de type fini sur $S$,
\item $U = C \setminus \{c\}$ pour un point fermé $c \in C$,
\item $A = \mathcal{O}_{C, c}$ soit de dimension $1$\footnote{Il s'ensuit
que $A$ est un anneau de valuation discrète ; voir
Algèbre, Lemme \ref{algebra-lemma-characterize-dvr}.
De plus, $c$ a pour image un point de type fini $s \in S$, et
$A$ est essentiellement de type fini sur $\mathcal{O}_{S, s}$.}
\end{enumerate}
alors, dans le diagramme commutatif
$$
\xymatrix{
\Spec(K) \ar[r] \ar[d] & X \ar[d]^f \\
\Spec(A) \ar[r] \ar@{-->}[ru] & Y
}
$$
où $K = \text{Frac}(A)$, il existe une flèche pointillée\footnote{D'après le
Lemme \ref{lemma-morphism-glueing-near-closed-point}, cela
équivaut à demander l'existence d'une flèche pointillée
rendant commutatif le premier diagramme.}
rendant le diagramme commutatif.
\end{enumerate}
\end{lemma}

\begin{proof}
Nous avons vu l'équivalence de (1) et (2), ainsi que le
fait que ces conditions entraînent (3), dans le
Lemme \ref{lemma-check-universally-closed-Noetherian}.
Il suffit donc de montrer que (3) entraîne (2).
Remarquons que, si la condition (3) est satisfaite par $f : X \to Y$,
elle l'est aussi par
$1 \times f : \mathbf{A}^n \times X \to \mathbf{A}^n \times Y$
(voir l'argument de la démonstration du
Lemme \ref{lemma-check-universally-closed-Noetherian}).
Il suffit donc de montrer que (3) entraîne que $f$ est fermé.

\medskip\noindent
Réduction au cas où $Y$ et $S$ sont affines ; on peut
sauter ce paragraphe.
Soit $S' \subset S$ un ouvert affine et $Y' \subset Y$ un
ouvert affine dont l'image est contenue dans $S'$. Posons $X' = f^{-1}(Y')$. Nous
affirmons que la restriction $f' : X' \to Y'$ de $f$, considérée comme morphisme
de schémas sur $S'$, possède aussi la propriété (3). Nous omettons les détails.
Si l'on montre que $f'$ est fermé pour tous les choix de
$S'$ et $Y'$, il s'ensuit que $f$ est fermé. Cela nous ramène
au cas traité au paragraphe suivant.

\medskip\noindent
Supposons $S$ et $Y$ affines. Soit $Z \subset X$ une partie fermée.
Munissons $Z$ de la structure de sous-schéma fermé réduit de $X$.
Il faut montrer que $E = f(Z)$ est fermé. Choisissons un point fermé $y \in Y$
dans l'adhérence de $f(Z)$. Il suffit de montrer $y \in E$.
Supposons $y \not \in E$ pour obtenir une contradiction.
L'image $s \in S$ de $y$ est un point de type fini de $S$ ; voir
Morphismes, Lemme \ref{morphisms-lemma-finite-type-points-morphism}.
Rappelons que $E$ est constructible
(Morphismes, Lemme \ref{morphisms-lemma-chevalley}).
Considérons l'intersection $\Spec(\mathcal{O}_{Y, y}) \cap E$.
C'est une partie constructible du spectre
(Morphismes, Lemme \ref{morphisms-lemma-inverse-image-constructible})
qui ne contient pas le point fermé.
Puisque le spectre épointé
$\Spec(\mathcal{O}_{Y, y}) \setminus \{y\}$ est de Jacobson
(Morphismes, Lemme \ref{morphisms-lemma-ubiquity-Jacobson-schemes}), on trouve un
point fermé $t \in \Spec(\mathcal{O}_{Y, y}) \setminus \{y\}$
tel que $t \in E$ (voir Topologie, Lemme \ref{topology-lemma-jacobson-inherited}).
Autrement dit, $t \in E$ est un point de $Y$
admettant une spécialisation immédiate $t \leadsto y$.
Puisque $t \in E$, la fibre schématique $Z_t$ est non vide.
Choisissons un point fermé $x \in Z_t$. On a en particulier
$[\kappa(x) : \kappa(t)] < \infty$ par le Nullstellensatz de Hilbert
(Morphismes, Lemme
\ref{morphisms-lemma-closed-point-fibre-locally-finite-type}).

\medskip\noindent
Notons $T = \overline{\{t\}} \subset Y$ le sous-schéma fermé intègre
dont l'espace topologique sous-jacent est celui indiqué
(Schémas, Définition \ref{schemes-definition-reduced-induced-scheme}).
Alors $t \in T$ est le point générique.
Notons $C \to T$ la normalisation de $T$ dans $\kappa(x)$ ; voir
Morphismes, section \ref{morphisms-section-normalization-X-in-Y}
(plus précisément, $C \to T$ est la normalisation de $T$ dans $x$,
où l'on considère $x = \Spec(\kappa(x)) \to T$ comme schéma sur $T$).
Puisque $S$ est un schéma de Nagata, $T$ l'est aussi
(Morphismes, Lemme \ref{morphisms-lemma-finite-type-nagata}).
Ainsi $C \to T$ est fini
(Morphismes, Lemme \ref{morphisms-lemma-nagata-normalization-finite-general}).
Puisque $t$ appartient à l'image, $C \to T$ est surjectif (car
l'image est fermée et $T$ est l'adhérence de $t$ dans $Y$).
Choisissons un point $c \in C$ d'image $y \in T$.
Puisque $y$ est un point fermé de $T$, $c$ est
un point fermé de $C$. Puisque $\dim(\mathcal{O}_{T, y}) = 1$,
on a $\dim(\mathcal{O}_{C, c}) = 1$ (la dimension est
au moins $1$, car $c$ n'est pas le point générique de $C$, et
au plus $1$, car $C \to T$ est fini).
Puisque le corps des fonctions de $C$ est $\kappa(x)$ et que $x$
est un point de $X$, on a une application rationnelle sur $Y$ de $C$ vers $X$
(voir par exemple Morphismes, Lemme
\ref{morphisms-lemma-rational-map-finite-presentation}).
Soit $C \supset U \to X$ un représentant (en particulier, $U$ est non vide).
On peut supposer $c \not \in U$ (remplacer $U$ par $U \setminus \{c\}$).
Puisque $c$ est un point fermé de codimension $1$ du schéma intègre $C$,
on a $C = U \amalg \{c\} \amalg \Sigma$
pour une partie fermée propre $\Sigma \subset C$.
En remplaçant $C$ par $C \setminus \Sigma$,
on obtient un diagramme commutatif comme dans (3).
D'après la seconde note de bas de page de l'énoncé,
l'existence de la flèche pointillée prolonge
l'application rationnelle à tout $C$ ; on obtient une contradiction,
car l'image de $c$ est alors un point de $Z$ d'image $y$.
\end{proof}

\begin{lemma}
\label{lemma-essentially-finite-type-criterion-separation}
Soit $S$ un schéma de Nagata (donc, en particulier, localement noethérien).
Soit $f : X \to Y$ un morphisme de schémas localement de type fini sur $S$.
Les conditions suivantes sont équivalentes :
\begin{enumerate}
\item $f$ est séparé,
\item pour tout diagramme commutatif
$$
\xymatrix{
U \ar[r] \ar[d] & X \ar[d]^f \\
C \ar[r] \ar@{..>}[ru] & Y
}
$$
de schémas sur $S$ tel que
\begin{enumerate}
\item $C$ soit un schéma intègre normal de type fini sur $S$,
\item $U = C \setminus \{c\}$ pour un point fermé $c \in C$,
\item $A = \mathcal{O}_{C, c}$ soit de dimension $1$\footnote{Il s'ensuit
que $A$ est un anneau de valuation discrète ; voir
Algèbre, Lemme \ref{algebra-lemma-characterize-dvr}.
De plus, $c$ a pour image un point de type fini $s \in S$, et
$A$ est essentiellement de type fini sur $\mathcal{O}_{S, s}$.}
\end{enumerate}
alors, dans le diagramme commutatif
$$
\xymatrix{
\Spec(K) \ar[r] \ar[d] & X \ar[d]^f \\
\Spec(A) \ar[r] \ar@{-->}[ru] & Y
}
$$
où $K = \text{Frac}(A)$, il existe au plus une flèche pointillée\footnote{D'après le
Lemme \ref{lemma-morphism-glueing-near-closed-point}, cela
équivaut à demander qu'il existe au plus une flèche pointillée
rendant commutatif le premier diagramme.}
rendant le diagramme commutatif.
\end{enumerate}
\end{lemma}

\begin{proof}
D'après le Lemme \ref{lemma-Noetherian-dvr-valuative-separation},
(1) entraîne (2). Supposons (2). Pour montrer que $f$ est séparé,
il faut montrer que $\Delta : X \to X \times_Y X$ est fermé. D'après
Morphismes, Lemme \ref{morphisms-lemma-finite-type-Noetherian-quasi-separated},
le morphisme $\Delta$ est quasi-compact.
D'après le Lemme \ref{lemma-essentially-finite-type-criterion-universally-closed},
il suffit de montrer que, pour tout diagramme commutatif
$$
\xymatrix{
U \ar[rr] \ar[d] & & X \ar[d]^\Delta \\
C \ar[rr]^{(a_1, a_2)} \ar@{..>}[rru] & & X \times_Y X
}
$$
de schémas sur $S$ tel que
\begin{enumerate}
\item $C$ soit un schéma intègre normal de type fini sur $S$,
\item $U = C \setminus \{c\}$ pour un point fermé $c \in C$,
\item $A = \mathcal{O}_{C, c}$ soit de dimension $1$,
\end{enumerate}
alors, dans le diagramme commutatif
$$
\xymatrix{
\Spec(K) \ar[r] \ar[d] & X \ar[d]^\Delta \\
\Spec(A) \ar[r] \ar@{-->}[ru] & X \times_Y X
}
$$
où $K = \text{Frac}(A)$, il existe une flèche pointillée
rendant le diagramme commutatif. D'après le
Lemme \ref{lemma-morphism-glueing-near-closed-point},
l'existence de la flèche pointillée dans le second diagramme
équivaut à celle de la flèche pointillée
dans le premier. De plus, cette existence
équivaut à demander $a_1 = a_2$. Or $a_1|_U = a_2|_U$ ;
l'hypothèse d'unicité (2) donne donc cette égalité,
ce qui achève la démonstration.
\end{proof}

\begin{lemma}
\label{lemma-essentially-finite-type-criterion-proper}
Soit $S$ un schéma de Nagata (donc, en particulier, localement noethérien).
Soit $f : X \to Y$ un morphisme quasi-compact de schémas localement
de type fini sur $S$. Les conditions suivantes sont équivalentes :
\begin{enumerate}
\item $f$ est propre,
\item pour tout diagramme commutatif
$$
\xymatrix{
U \ar[r] \ar[d] & X \ar[d]^f \\
C \ar[r] \ar@{..>}[ru] & Y
}
$$
de schémas sur $S$ tel que
\begin{enumerate}
\item $C$ soit un schéma intègre normal de type fini sur $S$,
\item $U = C \setminus \{c\}$ pour un point fermé $c \in C$,
\item $A = \mathcal{O}_{C, c}$ soit de dimension $1$\footnote{Il s'ensuit
que $A$ est un anneau de valuation discrète ; voir
Algèbre, Lemme \ref{algebra-lemma-characterize-dvr}.
De plus, $c$ a pour image un point de type fini $s \in S$, et
$A$ est essentiellement de type fini sur $\mathcal{O}_{S, s}$.}
\end{enumerate}
alors, dans le diagramme commutatif
$$
\xymatrix{
\Spec(K) \ar[r] \ar[d] & X \ar[d]^f \\
\Spec(A) \ar[r] \ar@{-->}[ru] & Y
}
$$
où $K = \text{Frac}(A)$, il existe exactement une flèche pointillée\footnote{D'après le
Lemme \ref{lemma-morphism-glueing-near-closed-point}, cela
équivaut à demander l'existence et l'unicité de la flèche pointillée
rendant commutatif le premier diagramme.}
rendant le diagramme commutatif.
\end{enumerate}
\end{lemma}

\begin{proof}
Cela découle formellement des Lemmes
\ref{lemma-essentially-finite-type-criterion-universally-closed} et
\ref{lemma-essentially-finite-type-criterion-separation}
et de la définition d'un morphisme propre : un morphisme de type fini, séparé
et universellement fermé.
\end{proof}






\section{Limites et dimensions des fibres}
\label{section-limits-dimension}

\noindent
Le lemme suivant est le plus souvent utilisé dans la situation du
Lemme \ref{lemma-descend-finite-presentation}
pour garantir que, si les fibres de la limite sont de dimension $\leq d$,
les fibres à un certain étage fini sont aussi de dimension $\leq d$.

\begin{lemma}
\label{lemma-limit-dimension}
Soit $I$ un ensemble ordonné filtrant.
Soit $(f_i : X_i \to S_i)$ un système projectif de morphismes de schémas
indexé par $I$. Faisons les hypothèses suivantes :
\begin{enumerate}
\item tous les morphismes $S_{i'} \to S_i$ sont affines,
\item tous les schémas $S_i$ sont quasi-compacts et quasi-séparés,
\item les morphismes $f_i$ sont de type fini,
\item les morphismes $X_{i'} \to X_i \times_{S_i} S_{i'}$ sont des immersions
fermées.
\end{enumerate}
Soit $f : X = \lim_i X_i \to S = \lim_i S_i$ la limite.
Soit $d \geq 0$.
Si toute fibre de $f$ est de dimension $\leq d$, alors, pour un certain $i$,
toute fibre de $f_i$ est de dimension $\leq d$.
\end{lemma}

\begin{proof}
Pour chaque $i$, soit $U_i = \{x \in X_i \mid \dim_x((X_i)_{f_i(x)}) \leq d\}$.
C'est un ouvert de $X_i$ ; voir
Morphismes, Lemme \ref{morphisms-lemma-openness-bounded-dimension-fibres}.
Posons $Z_i = X_i \setminus U_i$ (muni de la structure réduite induite).
Il faut montrer que $Z_i = \emptyset$ pour un certain $i$.
Sinon, $Z = \lim Z_i \not = \emptyset$ ; voir
Lemme \ref{lemma-limit-nonempty}.
Soit $z \in Z$ un point. Remarquons que $Z \subset X$ est un sous-schéma fermé.
Posons $s = f(z)$. Pour chaque $i$, soit $s_i \in S_i$ l'image
de $s$. Remarquons que $Z_s$ est la limite des schémas $(Z_i)_{s_i}$,
et que $Z_s$ est aussi la limite des schémas $(Z_i)_{s_i}$ après changement
de base à $\kappa(s)$. De plus, tous les morphismes
$$
Z_s
\longrightarrow
(Z_{i'})_{s_{i'}} \times_{\Spec(\kappa(s_{i'}))} \Spec(\kappa(s))
\longrightarrow
(Z_i)_{s_i} \times_{\Spec(\kappa(s_i))} \Spec(\kappa(s))
\longrightarrow
X_s
$$
sont des immersions fermées par l'hypothèse (4). Ainsi $Z_s$ est l'intersection
schématique des sous-schémas fermés
$(Z_i)_{s_i} \times_{\Spec(\kappa(s_i))} \Spec(\kappa(s))$
dans $X_s$. Puisque toutes les composantes irréductibles des schémas
$(Z_i)_{s_i} \times_{\Spec(\kappa(s_i))} \Spec(\kappa(s))$
sont de dimension $> d$ et contiennent $z$, on conclut que
$Z_s$ contient une composante irréductible de dimension $> d$ passant
par $z$, ce qui contredit $Z_s \subset X_s$ et
$\dim(X_s) \leq d$.
\end{proof}

\begin{lemma}
\label{lemma-descend-quasi-finite}
Reprenons les notations et les hypothèses de la Situation \ref{situation-descent-property}.
Si
\begin{enumerate}
\item $f$ est un morphisme quasi-fini,
\item $f_0$ est localement de type fini,
\end{enumerate}
alors il existe un $i \geq 0$ tel que $f_i$ soit quasi-fini.
\end{lemma}

\begin{proof}
Cela découle immédiatement du Lemme \ref{lemma-limit-dimension}.
\end{proof}

\begin{lemma}
\label{lemma-eventually-relative-dimension}
Reprenons les hypothèses et les notations de la Situation \ref{situation-descent-property}.
Soit $d \geq 0$. Si
\begin{enumerate}
\item $f$ est de dimension relative $\leq d$
(Morphismes, Définition \ref{morphisms-definition-relative-dimension-d}),
\item $f_0$ est localement de type fini,
\end{enumerate}
alors il existe un $i$ tel que $f_i$ soit de dimension relative $\leq d$.
\end{lemma}

\begin{proof}
Cela découle immédiatement du Lemme \ref{lemma-limit-dimension}.
\end{proof}

\begin{lemma}
\label{lemma-descend-dimension-d}
Reprenons les notations et les hypothèses de la Situation \ref{situation-descent-property}.
Si
\begin{enumerate}
\item $f$ est de dimension relative $d$,
\item $f_0$ est localement de présentation finie,
\end{enumerate}
alors il existe un $i \geq 0$ tel que $f_i$
soit de dimension relative $d$.
\end{lemma}

\begin{proof}
D'après le Lemme \ref{lemma-limit-dimension}, on peut supposer que toutes les fibres
de $f_0$ sont de dimension $\leq d$. D'après Morphismes, Lemme
\ref{morphisms-lemma-openness-bounded-dimension-fibres-finite-presentation},
l'ensemble $U_0 \subset X_0$ des points $x \in X_0$ tels que
la dimension de la fibre de $X_0 \to Y_0$ en $x$ soit $\leq d - 1$
est ouvert et rétrocompact dans $X_0$. Son complémentaire
$E = X_0 \setminus U_0$ est donc constructible.
De plus, l'image de $X \to X_0$ est contenue dans $E$
d'après Morphismes, Lemme \ref{morphisms-lemma-dimension-fibre-after-base-change}.
Ainsi, pour $i \gg 0$, l'image
de $X_i \to X_0$ est contenue dans $E$
(Lemme \ref{lemma-limit-contained-in-constructible}). Toutes les fibres
de $X_i \to Y_i$ sont alors de dimension $d$ d'après le résultat déjà cité
Morphismes, Lemme \ref{morphisms-lemma-dimension-fibre-after-base-change}.
\end{proof}

\begin{lemma}
\label{lemma-approximate-given-relative-dimension}
Soit $S$ un schéma quasi-compact et quasi-séparé.
Soit $f : X \to S$ un morphisme de présentation finie.
Soit $d \geq 0$ un entier.
Si $Z \subset X$ est un sous-schéma fermé tel que
$\dim(Z_s) \leq d$ pour tout $s \in S$, il existe un
sous-schéma fermé $Z' \subset X$ tel que
\begin{enumerate}
\item $Z \subset Z'$,
\item $Z' \to X$ soit de présentation finie,
\item $\dim(Z'_s) \leq d$ pour tout $s \in S$.
\end{enumerate}
\end{lemma}

\begin{proof}
D'après la
Proposition \ref{proposition-approximate},
on peut écrire $S = \lim S_i$ comme la limite d'un système projectif
filtrant de schémas noethériens à morphismes de transition affines. D'après le
Lemme \ref{lemma-descend-finite-presentation},
on peut supposer qu'il existe un système de morphismes
$f_i : X_i \to S_i$ de présentation finie tels que
$X_{i'} = X_i \times_{S_i} S_{i'}$
pour tout $i' \geq i$ et que $X = X_i \times_{S_i} S$.
Soit $Z_i \subset X_i$ l'image schématique de
$Z \to X \to X_i$. Pour $i' \geq i$, le morphisme $X_{i'} \to X_i$
envoie alors $Z_{i'}$ dans $Z_i$, et le morphisme induit
$Z_{i'} \to Z_i \times_{S_i} S_{i'}$ est une immersion fermée. D'après le
Lemme \ref{lemma-limit-dimension},
les fibres de $Z_i \to S_i$
sont toutes de dimension $\leq d$ pour un indice convenable $i \in I$.
Fixons un tel $i$ et posons $Z' = Z_i \times_{S_i} S \subset X$.
Puisque $S_i$ est noethérien, $X_i$ l'est aussi, et
le morphisme $Z_i \to X_i$ est donc de présentation finie.
Son changement de base $Z' \to X$ est donc lui aussi de présentation finie.
De plus, les fibres de $Z' \to S$ sont des changements de base des fibres
de $Z_i \to S_i$, donc sont de dimension $\leq d$.
\end{proof}






\section{Changement de base en degré maximal}
\label{section-top-degree}

\noindent
Pour un morphisme propre et un module quasi-cohérent de type fini,
l'application de changement de base est un isomorphisme en degré maximal.

\begin{lemma}
\label{lemma-top-cohomology-functor}
Soit $f : X \to Y$ un morphisme de schémas. Soit $d \geq 0$. Faisons les hypothèses suivantes :
\begin{enumerate}
\item $X$ et $Y$ sont quasi-compacts et quasi-séparés,
\item $R^if_*\mathcal{F} = 0$ pour $i > d$ et
tout $\mathcal{O}_X$-module quasi-cohérent $\mathcal{F}$.
\end{enumerate}
On a alors les propriétés suivantes :
\begin{enumerate}
\item[(a)] pour tout diagramme de changement de base
$$
\xymatrix{
X' \ar[d]_{f'} \ar[r]_{g'} & X \ar[d]^f \\
Y' \ar[r]^g & Y
}
$$
on a $R^if'_*\mathcal{F}' = 0$ pour $i > d$ et tout
$\mathcal{O}_{X'}$-module quasi-cohérent $\mathcal{F}'$,
\item[(b)]
$R^df'_*(\mathcal{F}' \otimes_{\mathcal{O}_{X'}} (f')^*\mathcal{G}') =
R^df'_*\mathcal{F}' \otimes_{\mathcal{O}_{Y'}} \mathcal{G}'$
pour tout $\mathcal{O}_{Y'}$-module quasi-cohérent $\mathcal{G}'$,
\item[(c)] la formation de $R^df'_*\mathcal{F}'$ commute à tout
changement de base ultérieur (voir l'explication dans la démonstration).
\end{enumerate}
\end{lemma}

\begin{proof}
Avant de démontrer ces assertions, expliquons le sens de (c). Supposons
qu'un carré cartésien supplémentaire
$$
\xymatrix{
X'' \ar[d]_{f''} \ar[r]_{h'} &
X' \ar[d]_{f'} \ar[r]_{g'} &
X \ar[d]^f \\
Y'' \ar[r]^h &
Y' \ar[r]^g &
Y
}
$$
soit accolé au diagramme donné. Si (a) est satisfaite, il existe une application
canonique $\gamma : h^*R^df'_*\mathcal{F}' \to R^df''_*(h')^*\mathcal{F}'$.
En effet, $\gamma$ est l'application sur les faisceaux de cohomologie de degré $d$
induite par la composée
$$
Lh^*Rf'_*\mathcal{F}' \longrightarrow
Rf''_*L(h')^*\mathcal{F}' \longrightarrow Rf''_*(h')^*\mathcal{F}'
$$
Ici, la première flèche est l'application de changement de base
(Cohomologie, Remarque \ref{cohomology-remark-base-change}), et
la seconde provient de l'application canonique
$L(g')^*\mathcal{F} \to (g')^*\mathcal{F}$.
De même, puisque $Rf'_*\mathcal{F}$ n'a pas de faisceau de cohomologie non nul en
degré $> d$ d'après (a), on a
$H^d(Lh^*Rf_*\mathcal{F}') = h^*R^df_*\mathcal{F}$.
L'assertion (c) signifie que $\gamma$ est un isomorphisme.

\medskip\noindent
Ceci étant, on peut vérifier (a), (b) et (c) localement sur $Y'$ et $Y''$.
Soit $V \subset Y$ un sous-schéma ouvert quasi-compact. Nous
affirmons que (1) et (2) sont satisfaites par $f|_{f^{-1}(V)} : f^{-1}(V) \to V$.
En effet, (1) est immédiate, et (2) résulte du fait que tout module quasi-cohérent
sur $f^{-1}(V)$ est la restriction d'un module quasi-cohérent
sur $X$ (Propriétés, Lemme \ref{properties-lemma-extend-trivial}) et que la formation
des images directes supérieures commute à la restriction aux ouverts.
On peut donc aussi travailler localement sur $Y$. Autrement dit, on peut
supposer $Y''$, $Y'$ et $Y$ affines.

\medskip\noindent
Démonstration de (a) lorsque $Y'$ et $Y$ sont affines. Dans ce cas, les morphismes
$g$ et $g'$ sont affines. Ainsi $g_* = Rg_*$ et $g'_* = Rg'_*$
(Cohomologie des schémas, Lemme \ref{coherent-lemma-relative-affine-vanishing}),
et $g_*$ s'identifie au foncteur de restriction des scalaires sur les modules
(Schémas, Lemme \ref{schemes-lemma-widetilde-pullback}).
Alors
$$
g_*(R^if'_*\mathcal{F}') = H^i(Rg_*Rf'_*\mathcal{F}') =
H^i(Rf_*Rg'_*\mathcal{F}') =
H^i(Rf_*g'_*\mathcal{F}') =
Rf^i_*g'_*\mathcal{F}'
$$
est nul par l'hypothèse (2). D'où (a), compte tenu de notre description de $g_*$.

\medskip\noindent
Démonstration de (b) lorsque $Y'$ est affine, disons $Y' = \Spec(R')$.
D'après (a), on a $H^{d + 1}(X', \mathcal{F}') = 0$
pour tout $\mathcal{O}_{X'}$-module quasi-cohérent $\mathcal{F}'$ ; voir
Cohomologie des schémas, Lemme
\ref{coherent-lemma-quasi-coherence-higher-direct-images-application}.
Considérons le foncteur $F$ sur les $R'$-modules défini par
$$
F(M) = H^d(X', \mathcal{F}' \otimes_{\mathcal{O}_{X'}} (f')^*\widetilde{M})
$$
D'après Cohomologie, Lemme \ref{cohomology-lemma-quasi-separated-cohomology-colimit},
ce foncteur commute aux sommes directes (c'est ici qu'intervient le fait que
$X$, et donc $X'$, est quasi-compact et quasi-séparé).
D'autre part, si $M_1 \to M_2 \to M_3 \to 0$ est une suite exacte,
alors
$$
\mathcal{F}' \otimes_{\mathcal{O}_{X'}} (f')^*\widetilde{M}_1 \to
\mathcal{F}' \otimes_{\mathcal{O}_{X'}} (f')^*\widetilde{M}_2 \to
\mathcal{F}' \otimes_{\mathcal{O}_{X'}} (f')^*\widetilde{M}_3 \to 0
$$
est une suite exacte de modules quasi-cohérents sur $X'$,
et l'annulation de la cohomologie supérieure établie ci-dessus donne une suite exacte
$$
F(M_1) \to F(M_2) \to F(M_3) \to 0
$$
Autrement dit, $F$ est exact à droite. Tout foncteur $R'$-linéaire exact à droite
$F : \text{Mod}_{R'} \to \text{Mod}_{R'}$
qui commute aux sommes directes est donné par le produit tensoriel avec un $R'$-module
(démonstration omise et laissée en exercice).
On obtient ainsi $F(M) = H^d(X', \mathcal{F}') \otimes_{R'} M$.
Puisque $R^d(f')_*\mathcal{F}'$ et
$R^d(f')_*(\mathcal{F}' \otimes_{\mathcal{O}_{X'}} (f')^*\widetilde{M})$
sont quasi-cohérents (Cohomologie des schémas, Lemme
\ref{coherent-lemma-quasi-coherence-higher-direct-images}),
l'égalité $F(M) = H^d(X', \mathcal{F}') \otimes_{R'} M$
se traduit par l'énoncé de (b).

\medskip\noindent
Démonstration de (c) lorsque $Y'' \to Y' \to Y$ sont des morphismes de schémas affines.
Écrivons $Y'' = \Spec(R'')$ et $Y' = \Spec(R')$.
Alors $R^df''_*(h')^*\mathcal{F}'$
est le module quasi-cohérent sur $Y'$ associé au $R''$-module
$H^d(X'', (h')^*\mathcal{F}')$. Or $h' : X'' \to X'$ est affine ;
ainsi $H^d(X'', (h')^*\mathcal{F}') = H^d(X, h'_*(h')^*\mathcal{F}')$
d'après le résultat déjà utilisé
Cohomologie des schémas, Lemme \ref{coherent-lemma-relative-affine-cohomology}.
On a
$$
h'_*(h')^*\mathcal{F}' =
\mathcal{F}' \otimes_{\mathcal{O}_{X'}} (f')^*\widetilde{R''}
$$
comme on le vérifie sur un recouvrement ouvert affine. Ainsi
$H^d(X'', (h')^*\mathcal{F}') = H^d(X', \mathcal{F}') \otimes_{R'} R''$
d'après (b), appliqué à $f'$, ce qui achève la démonstration.
\end{proof}

\begin{lemma}
\label{lemma-higher-direct-images-zero-above-dimension-fibre}
Soit $f : X \to Y$ un morphisme de schémas. Soit $y \in Y$.
Supposons $f$ propre et $\dim(X_y) = d$.
Alors
\begin{enumerate}
\item pour $\mathcal{F} \in \QCoh(\mathcal{O}_X)$,
on a $(R^if_*\mathcal{F})_y = 0$ pour tout $i > d$,
\item il existe un voisinage ouvert affine $V \subset Y$
de $y$ tel que $f^{-1}(V) \to V$ et $d$ satisfassent les
hypothèses et les conclusions du Lemme \ref{lemma-top-cohomology-functor}.
\end{enumerate}
\end{lemma}

\begin{proof}
D'après Morphismes, Lemme
\ref{morphisms-lemma-openness-bounded-dimension-fibres},
et puisque $f$ est fermé, on peut trouver un voisinage ouvert affine
$V$ de $y$ tel que les fibres au-dessus des points de $V$ soient toutes de dimension
$\leq d$. On peut donc supposer que $X \to Y$ est un morphisme propre dont
toutes les fibres sont de dimension $\leq d$, avec $Y$ affine.
Nous montrerons (2), ce qui entraînera immédiatement (1)
pour tout $y \in Y$.

\medskip\noindent
D'après le Lemme \ref{lemma-proper-limit-of-proper-finite-presentation},
on peut écrire $X = \lim X_i$ comme une limite projective, où $X_i \to Y$ est propre
et de présentation finie, et où les morphismes $X \to X_i$
ainsi que les morphismes de transition sont des immersions fermées.
Pour un certain $i$, les fibres de $X_i \to Y$ sont de dimension $\leq d$ ;
voir Lemme \ref{lemma-limit-dimension}.
Pour un $\mathcal{O}_X$-module quasi-cohérent $\mathcal{F}$, on a
$R^pf_*\mathcal{F} = R^pf_{i, *}(X \to X_i)_*\mathcal{F}$ d'après
Cohomologie des schémas, Lemme \ref{coherent-lemma-relative-affine-vanishing},
et Leray (Cohomologie, Lemme \ref{cohomology-lemma-relative-Leray}).
On peut donc remplacer $X$ par $X_i$ et
se ramener au cas traité au paragraphe suivant.

\medskip\noindent
Supposons $Y$ affine et $f : X \to Y$ propre et de présentation finie,
toutes les fibres étant de dimension $\leq d$. Il suffit de montrer que
$H^p(X, \mathcal{F}) = 0$ pour $p > d$. En effet, d'après
Cohomologie des schémas, Lemme
\ref{coherent-lemma-quasi-coherence-higher-direct-images-application},
on a $H^p(X, \mathcal{F}) = H^0(Y, R^pf_*\mathcal{F})$.
D'autre part, $R^pf_*\mathcal{F}$ est quasi-cohérent sur $Y$
d'après Cohomologie des schémas, Lemme
\ref{coherent-lemma-quasi-coherence-higher-direct-images} ;
l'annulation de ses sections globales entraîne donc son annulation.
Écrivons $Y = \lim_{i \in I} Y_i$ comme une limite projective de schémas affines,
où $Y_i$ est le spectre d'un anneau noethérien
(par exemple d'une $\mathbf{Z}$-algèbre de type fini).
On peut choisir un élément $0 \in I$ et un morphisme de type fini
$X_0 \to Y_0$ tels que $X \cong Y \times_{Y_0} X_0$ ; voir
Lemme \ref{lemma-descend-finite-presentation}.
En augmentant $0$, on peut supposer $X_0 \to Y_0$ propre
(Lemme \ref{lemma-eventually-proper})
et les fibres de $X_0 \to Y_0$ de dimension $\leq d$
(Lemme \ref{lemma-limit-dimension}).
Puisque $X \to X_0$ est affine, on a
$H^p(X, \mathcal{F}) = H^p(X_0, (X \to X_0)_*\mathcal{F})$ d'après
Cohomologie des schémas, Lemme \ref{coherent-lemma-relative-affine-cohomology}.
Cela nous ramène au cas traité au paragraphe suivant.

\medskip\noindent
Supposons $Y$ affine noethérien et $f : X \to Y$ propre,
toutes les fibres étant de dimension $\leq d$.
Dans ce cas, on peut écrire $\mathcal{F} = \colim \mathcal{F}_i$
comme une limite inductive filtrante de $\mathcal{O}_X$-modules cohérents ; voir
Propriétés, Lemme
\ref{properties-lemma-directed-colimit-finite-presentation}.
Alors $H^p(X, \mathcal{F}) = \colim H^p(X, \mathcal{F}_i)$ d'après
Cohomologie, Lemme \ref{cohomology-lemma-quasi-separated-cohomology-colimit}.
On peut donc supposer $\mathcal{F}$ cohérent.
Dans ce cas, on a $(R^pf_*\mathcal{F})_y = 0$ pour
tout $y \in Y$ d'après Cohomologie des schémas, Lemme
\ref{coherent-lemma-higher-direct-images-zero-above-dimension-fibre}.
Ainsi $R^pf_*\mathcal{F} = 0$, donc
$H^p(X, \mathcal{F}) = 0$ (voir ci-dessus), ce qui conclut.
\end{proof}

\begin{lemma}
\label{lemma-proper-top-cohomology-finite-type}
Soit $f : X \to Y$ un morphisme de schémas. Soit $d \geq 0$. Soit $\mathcal{F}$
un $\mathcal{O}_X$-module. Faisons les hypothèses suivantes :
\begin{enumerate}
\item $f$ est un morphisme propre dont toutes les fibres sont de dimension $\leq d$,
\item $\mathcal{F}$ est un $\mathcal{O}_X$-module quasi-cohérent de type fini.
\end{enumerate}
Alors $R^df_*\mathcal{F}$ est un $\mathcal{O}_X$-module quasi-cohérent
de type fini.
\end{lemma}

\begin{proof}
Le module $R^df_*\mathcal{F}$ est quasi-cohérent d'après
Cohomologie des schémas, Lemme
\ref{coherent-lemma-quasi-coherence-higher-direct-images}.
La question est locale sur $Y$ ; on peut donc supposer $Y$ affine.
Écrivons $Y = \Spec(R)$. Il suffit alors de montrer que $H^d(X, \mathcal{F})$
est un $R$-module de type fini.

\medskip\noindent
D'après le Lemme \ref{lemma-proper-limit-of-proper-finite-presentation},
on peut écrire $X = \lim X_i$ comme une limite projective, où $X_i \to Y$ est propre
et de présentation finie, et où les morphismes $X \to X_i$
ainsi que les morphismes de transition sont des immersions fermées.
Pour un certain $i$, les fibres de $X_i \to Y$ sont de dimension $\leq d$ ;
voir Lemme \ref{lemma-limit-dimension}. On a
$R^pf_*\mathcal{F} = R^pf_{i, *}(X \to X_i)_*\mathcal{F}$ d'après
Cohomologie des schémas, Lemme \ref{coherent-lemma-relative-affine-vanishing},
et Leray (Cohomologie, Lemme \ref{cohomology-lemma-relative-Leray}).
On peut donc remplacer $X$ par $X_i$ et
se ramener au cas traité au paragraphe suivant.

\medskip\noindent
Supposons $Y$ affine et $f : X \to Y$ propre et de présentation finie,
toutes les fibres étant de dimension $\leq d$.
On peut écrire $\mathcal{F}$ comme quotient d'un
$\mathcal{O}_X$-module de présentation finie $\mathcal{F}'$ ; voir
Propriétés, Lemme
\ref{properties-lemma-finite-directed-colimit-surjective-maps}.
L'application $H^d(X, \mathcal{F}') \to H^d(X, \mathcal{F})$ est
surjective, car $H^{d + 1}(X, \Ker(\mathcal{F}' \to \mathcal{F})) = 0$
d'après l'annulation de la cohomologie supérieure établie dans le
Lemme \ref{lemma-higher-direct-images-zero-above-dimension-fibre}
(ou sa démonstration). On est ainsi ramené au cas traité au paragraphe suivant.

\medskip\noindent
Supposons $Y = \Spec(R)$ affine et $f : X \to Y$
propre et de présentation finie,
toutes les fibres étant de dimension $\leq d$, et $\mathcal{F}$ un
$\mathcal{O}_X$-module de présentation finie.
Écrivons $Y = \lim_{i \in I} Y_i$ comme une limite projective de schémas affines,
où $Y_i = \Spec(R_i)$ est le spectre d'un anneau noethérien
(par exemple d'une $\mathbf{Z}$-algèbre de type fini).
On peut choisir un élément $0 \in I$ et un morphisme de type fini
$X_0 \to Y_0$ tels que $X \cong Y \times_{Y_0} X_0$ ; voir
Lemme \ref{lemma-descend-finite-presentation}.
En augmentant $0$, on peut supposer $X_0 \to Y_0$ propre
(Lemme \ref{lemma-eventually-proper})
et les fibres de $X_0 \to Y_0$ de dimension $\leq d$
(Lemme \ref{lemma-limit-dimension}).
En augmentant $0$, on peut supposer qu'il existe un
$\mathcal{O}_{X_0}$-module cohérent $\mathcal{F}_0$ dont l'image
inverse est $\mathcal{F}$ ; voir
Lemme \ref{lemma-descend-modules-finite-presentation}.
D'après le Lemme \ref{lemma-top-cohomology-functor},
on a
$$
H^d(X, \mathcal{F}) = H^d(X_0, \mathcal{F}_0) \otimes_{R_0} R
$$
Cela achève la démonstration, puisque le module de cohomologie
$H^d(X_0, \mathcal{F}_0)$ est de type fini d'après
Cohomologie des schémas, Lemme
\ref{coherent-lemma-proper-over-affine-cohomology-finite}.
\end{proof}

\begin{lemma}
\label{lemma-proper-top-cohomology-finite-presentation}
Soit $f : X \to Y$ un morphisme de schémas. Soit $d \geq 0$.
Soit $\mathcal{F}$ un $\mathcal{O}_X$-module. Faisons les hypothèses suivantes :
\begin{enumerate}
\item $f$ est un morphisme propre de présentation finie
dont toutes les fibres sont de dimension $\leq d$,
\item $\mathcal{F}$ est un $\mathcal{O}_X$-module de présentation finie.
\end{enumerate}
Alors $R^df_*\mathcal{F}$ est un $\mathcal{O}_X$-module
de présentation finie.
\end{lemma}

\begin{proof}
La démonstration est exactement la même que celle du
Lemme \ref{lemma-proper-top-cohomology-finite-type},
à ceci près que l'on peut omettre le troisième paragraphe.
Nous omettons les détails.
\end{proof}











\section{Recollement au voisinage des fibres fermées}
\label{section-change-over-closed-points}

\noindent
En appliquant la théorie précédente au spectre d'un anneau local, on obtient
le résultat de recollement suivant pour les schémas relatifs.

\begin{lemma}
\label{lemma-glueing-near-closed-point}
Soit $S$ un schéma. Soit $s \in S$ un point fermé tel que
$U = S \setminus \{s\} \to S$ soit quasi-compact. En posant
$V = \Spec(\mathcal{O}_{S, s}) \setminus \{s\}$, on a
une équivalence de catégories
$$
\left\{
\begin{matrix}
X \to S\text{ de présentation finie}
\end{matrix}
\right\}
\longrightarrow
\left\{
\vcenter{
\xymatrix{
X' \ar[d] & Y' \ar[d] \ar[l] \ar[r] & Y \ar[d] \\
U & V \ar[l] \ar[r] & \Spec(\mathcal{O}_{S, s})
}
}
\right\}
$$
où, à droite, on considère les diagrammes commutatifs
dont les carrés sont cartésiens et les flèches verticales
de présentation finie.
\end{lemma}

\begin{proof}
Soit $W \subset S$ un voisinage ouvert de $s$. Par
recollement des schémas relatifs, voir
Constructions, section \ref{constructions-section-relative-glueing},
le foncteur
$$
\left\{
\begin{matrix}
X \to S\text{ de présentation finie}
\end{matrix}
\right\}
\longrightarrow
\left\{
\vcenter{
\xymatrix{
X' \ar[d] & Y' \ar[d] \ar[l] \ar[r] & Y \ar[d] \\
U & W \setminus \{s\} \ar[l] \ar[r] & W
}
}
\right\}
$$
est une équivalence de catégories. On a
$\mathcal{O}_{S, s} = \colim \mathcal{O}_W(W)$, où
$W$ parcourt les voisinages ouverts affines de $s$.
Ainsi $\Spec(\mathcal{O}_{S, s}) = \lim W$, où $W$
parcourt les voisinages ouverts affines de $s$.
La catégorie des schémas de présentation finie
sur $\Spec(\mathcal{O}_{S, s})$ est donc la limite des
catégories de schémas de présentation finie sur
$W$, où $W$ parcourt les voisinages ouverts affines
de $s$ ; voir
Lemme \ref{lemma-descend-finite-presentation}.
Pour tout ouvert affine $s \in W$, l'intersection $U \cap W$
est quasi-compacte, puisque $U \to S$ est quasi-compact.
Ainsi $V = \lim W \cap U = \lim W \setminus \{s\}$ est une limite de
schémas quasi-compacts et quasi-séparés (voir
Lemme \ref{lemma-directed-inverse-system-has-limit}).
La catégorie des schémas de présentation finie
sur $V$ est donc elle aussi la limite des
catégories de schémas de présentation finie sur
$W \cap U$, où $W$ parcourt les voisinages ouverts affines
de $s$. Le lemme découle formellement de la combinaison
de ces résultats.
\end{proof}

\begin{lemma}
\label{lemma-glueing-near-closed-point-modules}
Soit $S$ un schéma. Soit $s \in S$ un point fermé tel que
$U = S \setminus \{s\} \to S$ soit quasi-compact. En posant
$V = \Spec(\mathcal{O}_{S, s}) \setminus \{s\}$, on a
une équivalence de catégories
$$
\left\{
\mathcal{O}_S\text{-modules }\mathcal{F}\text{ de présentation finie}
\right\}
\longrightarrow
\left\{
(\mathcal{G}, \mathcal{H}, \alpha)
\right\}
$$
où, à droite, on considère les triplets
constitués d'un $\mathcal{O}_U$-module $\mathcal{G}$ de
présentation finie, d'un $\mathcal{O}_{\Spec(\mathcal{O}_{S, s})}$-module
$\mathcal{H}$ de présentation finie et d'un isomorphisme
$\alpha : \mathcal{G}|_V \to \mathcal{H}|_V$ de
$\mathcal{O}_V$-modules.
\end{lemma}

\begin{proof}
On peut démontrer ce résultat
en reprenant la démonstration du Lemme \ref{lemma-glueing-near-closed-point}
à l'aide du Lemme \ref{lemma-descend-modules-finite-presentation},
ou le déduire du Lemme \ref{lemma-glueing-near-closed-point}
à l'aide de l'équivalence entre les modules quasi-cohérents et
les « fibrés vectoriels » de
Constructions, section \ref{constructions-section-vector-bundle}.
Nous omettons les détails.
\end{proof}

\begin{lemma}
\label{lemma-glueing-near-point}
Soit $S$ un schéma. Soit $U \subset S$ un ouvert rétrocompact.
Soit $s \in S$ un point du complémentaire de $U$. En posant
$V = \Spec(\mathcal{O}_{S, s}) \cap U$, on a
une équivalence de catégories
$$
\colim_{s \in U' \supset U\text{ ouvert}}
\left\{
\vcenter{
\xymatrix{
X \ar[d] \\
U'
}
}
\right\}
\longrightarrow
\left\{
\vcenter{
\xymatrix{
X' \ar[d] & Y' \ar[d] \ar[l] \ar[r] & Y \ar[d] \\
U & V \ar[l] \ar[r] & \Spec(\mathcal{O}_{S, s})
}
}
\right\}
$$
où, à gauche, la flèche verticale est de présentation
finie, et où, à droite, on considère les diagrammes commutatifs
dont les carrés sont cartésiens et les flèches verticales
de présentation finie.
\end{lemma}

\begin{proof}
Soit $W \subset S$ un voisinage ouvert de $s$. Par
recollement des schémas relatifs, voir
Constructions, section \ref{constructions-section-relative-glueing},
le foncteur
$$
\left\{
\begin{matrix}
X \to U' = U \cup W \text{ de présentation finie}
\end{matrix}
\right\}
\longrightarrow
\left\{
\vcenter{
\xymatrix{
X' \ar[d] & Y' \ar[d] \ar[l] \ar[r] & Y \ar[d] \\
U & W \cap U \ar[l] \ar[r] & W
}
}
\right\}
$$
est une équivalence de catégories. On a
$\mathcal{O}_{S, s} = \colim \mathcal{O}_W(W)$, où
$W$ parcourt les voisinages ouverts affines de $s$.
Ainsi $\Spec(\mathcal{O}_{S, s}) = \lim W$, où $W$
parcourt les voisinages ouverts affines de $s$.
La catégorie des schémas de présentation finie
sur $\Spec(\mathcal{O}_{S, s})$ est donc la limite des
catégories de schémas de présentation finie sur
$W$, où $W$ parcourt les voisinages ouverts affines
de $s$ ; voir
Lemme \ref{lemma-descend-finite-presentation}.
Pour tout ouvert affine $s \in W$, l'intersection $U \cap W$
est quasi-compacte, puisque $U \to S$ est quasi-compact.
Ainsi $V = \lim W \cap U$ est une limite de
schémas quasi-compacts et quasi-séparés (voir
Lemme \ref{lemma-directed-inverse-system-has-limit}).
La catégorie des schémas de présentation finie
sur $V$ est donc elle aussi la limite des
catégories de schémas de présentation finie sur
$W \cap U$, où $W$ parcourt les voisinages ouverts affines
de $s$. Le lemme découle formellement de la combinaison
de ces résultats.
\end{proof}

\begin{lemma}
\label{lemma-glueing-near-point-properties}
Reprenons les notations et les hypothèses du Lemme \ref{lemma-glueing-near-point}.
Soit $U \subset U' \subset X$ un ouvert contenant $s$.
\begin{enumerate}
\item Soit $f' : X \to U'$ correspondant à $f : X' \to U$ et
$g : Y \to \Spec(\mathcal{O}_{S, s})$
par l'équivalence. Si $f$ et $g$ sont séparés, propres, finis ou étales,
alors, quitte à rétrécir $U'$, le morphisme $f'$ possède la même propriété.
\item Soit $a : X_1 \to X_2$
un morphisme de schémas de présentation finie sur $U'$,
dont les changements de base sont $a' : X'_1 \to X'_2$ sur $U$ et
$b : Y_1 \to Y_2$ sur $\Spec(\mathcal{O}_{S, s})$.
Si $a'$ et $b$ sont séparés, propres, finis ou étales,
alors, quitte à rétrécir $U'$, le morphisme $a$ possède la même propriété.
\end{enumerate}
\end{lemma}

\begin{proof}
Démonstration de (1). Rappelons que $\Spec(\mathcal{O}_{S, s})$ est la limite des
voisinages ouverts affines de $s$ dans $S$. Puisque $g$ possède la propriété
considérée, la restriction de $f'$ à l'un de ces
voisinages ouverts affines la possède aussi ; voir
Lemmes \ref{lemma-descend-separated-finite-presentation},
\ref{lemma-eventually-proper},
\ref{lemma-descend-finite-finite-presentation} et
\ref{lemma-descend-etale}.
Puisque $f'$ possède la propriété donnée sur $U$, car $f$ la possède,
on conclut en vérifiant cette propriété localement sur la base.

\medskip\noindent
Démonstration de (2). Si l'on écrit $\Spec(\mathcal{O}_{S, s}) = \lim W$,
où $W$ parcourt les voisinages ouverts affines de $s$ dans $S$,
on a $Y_i = \lim W \times_S X_i$. On peut donc utiliser
exactement les mêmes arguments que dans la démonstration de (1).
\end{proof}

\begin{lemma}
\label{lemma-glueing-near-multiple-closed-points}
Soit $S$ un schéma. Soient $s_1, \ldots, s_n \in S$ des points fermés
deux à deux distincts tels que
$U = S \setminus \{s_1, \ldots, s_n\} \to S$ soit quasi-compact. En posant
$S_i = \Spec(\mathcal{O}_{S, s_i})$ et $U_i = S_i \setminus \{s_i\}$,
on a une équivalence de catégories
$$
FP_S \longrightarrow
FP_U \times_{(FP_{U_1} \times \ldots \times FP_{U_n})}
(FP_{S_1} \times \ldots \times FP_{S_n})
$$
où $FP_T$ est la catégorie des schémas de présentation finie sur
le schéma $T$.
\end{lemma}

\begin{proof}
Pour $n = 1$, c'est le Lemme \ref{lemma-glueing-near-closed-point}.
Pour $n > 1$, le lemme peut se démontrer exactement de la même manière ou
s'en déduire. Supposons par exemple que les $f_i : X_i \to S_i$
soient des objets de $FP_{S_i}$, que $f : X \to U$ soit un objet
de $FP_U$ et que l'on se donne des isomorphismes $X_i \times_{S_i} U_i = X \times_U U_i$.
D'après le Lemme \ref{lemma-glueing-near-closed-point}, on peut trouver
un morphisme $f' : X' \to U' = S \setminus \{s_1, \ldots, s_{n - 1}\}$
de présentation finie, isomorphe à
$X_i$ sur $S_i$ et à $X$ sur $U$, ces
isomorphismes étant compatibles avec l'isomorphisme donné
$X_i \times_{S_n} U_n = X \times_U U_n$.
On peut alors appliquer l'hypothèse de récurrence à
$f_i : X_i \to S_i$, $i \leq n - 1$,
$f' : X' \to U'$ et aux isomorphismes
induits $X_i \times_{S_i} U_i = X' \times_{U'} U_i$, $i \leq n - 1$.
Cela montre la surjectivité essentielle. Nous omettons la démonstration de
la pleine fidélité.
\end{proof}






\section{Application aux modifications}
\label{section-modifications-at-a-point}

\noindent
À l'aide des résultats de la section \ref{section-change-over-closed-points},
on peut décrire la catégorie des modifications d'un
schéma au-dessus d'un point fermé en termes d'anneau local.

\begin{lemma}
\label{lemma-modifications}
Soit $S$ un schéma. Soit $s \in S$ un point fermé tel que
$U = S \setminus \{s\} \to S$ soit quasi-compact. En posant
$V = \Spec(\mathcal{O}_{S, s}) \setminus \{s\}$, le foncteur de changement de base
$$
\left\{
\begin{matrix}
f : X \to S\text{ de présentation finie} \\
f^{-1}(U) \to U\text{ est un isomorphisme}
\end{matrix}
\right\}
\longrightarrow
\left\{
\begin{matrix}
g : Y \to \Spec(\mathcal{O}_{S, s})\text{ de présentation finie} \\
g^{-1}(V) \to V\text{ est un isomorphisme}
\end{matrix}
\right\}
$$
est une équivalence de catégories.
\end{lemma}

\begin{proof}
C'est un cas particulier du Lemme \ref{lemma-glueing-near-closed-point}.
\end{proof}

\begin{lemma}
\label{lemma-modifications-properties}
Reprenons les notations et les hypothèses du Lemme \ref{lemma-modifications}.
Soit $f : X \to S$ correspondant à $g : Y \to \Spec(\mathcal{O}_{S, s})$
par l'équivalence. Alors $f$ est séparé, propre, fini, étale,
et d'autres propriétés à ajouter ici, si et seulement si $g$ l'est.
\end{lemma}

\begin{proof}
Le caractère séparé, propre, entier, fini, etc.,
est stable par changement de base. Voir
Schémas, Lemme \ref{schemes-lemma-separated-permanence},
et
Morphismes, Lemmes \ref{morphisms-lemma-base-change-proper} et
\ref{morphisms-lemma-base-change-finite}.
Ainsi, si $f$ possède la propriété, $g$ la possède aussi.
La réciproque découle du Lemme \ref{lemma-glueing-near-point-properties},
mais nous en donnons aussi une démonstration directe.
En effet, si $g$ possède la propriété, $f$ la possède dans un voisinage de $s$ d'après les
Lemmes \ref{lemma-descend-separated-finite-presentation},
\ref{lemma-eventually-proper},
\ref{lemma-descend-finite-finite-presentation} et
\ref{lemma-descend-etale}.
Puisque $f$ possède évidemment la propriété considérée sur $S \setminus \{s\}$,
on conclut en vérifiant cette propriété localement sur la base.
\end{proof}

\begin{remark}
\label{remark-more-general-modification}
Le lemme précédent se généralise comme suit. Soit $S$ un schéma et
soit $T \subset S$ une partie fermée. Supposons qu'il existe un système cofinal
de voisinages ouverts $T \subset W_i$ tel que
(1) $W_i \setminus T$ soit quasi-compact et
(2) $W_i \subset W_j$ soit un morphisme affine.
Alors $W = \lim W_i$ est un schéma contenant $T$
comme sous-schéma fermé. Posons $U = X \setminus T$ et $V = W \setminus T$.
Le foncteur de changement de base
$$
\left\{
\begin{matrix}
f : X \to S\text{ de présentation finie} \\
f^{-1}(U) \to U\text{ est un isomorphisme}
\end{matrix}
\right\}
\longrightarrow
\left\{
\begin{matrix}
g : Y \to W\text{ de présentation finie} \\
g^{-1}(V) \to V\text{ est un isomorphisme}
\end{matrix}
\right\}
$$
est alors une équivalence de catégories. Si nous en avons besoin un jour,
nous transformerons cette remarque en lemme et en donnerons une démonstration détaillée.
\end{remark}



\section{Descente des schémas de type fini}
\label{section-finite-type-quasi-separated}

\noindent
Cette section poursuit le thème de la
section \ref{section-finite-type-closed-in-finite-presentation}
dans l'esprit des résultats de la
section \ref{section-descending-relative}.

\begin{situation}
\label{situation-limit-noetherian}
Soit $S = \lim_{i \in I} S_i$ la limite d'un système projectif filtrant de
schémas noethériens à morphismes de transition affines $S_{i'} \to S_i$
pour $i' \geq i$.
\end{situation}

\begin{lemma}
\label{lemma-good-diagram}
Plaçons-nous dans la Situation \ref{situation-limit-noetherian}.
Soit $X \to S$ quasi-séparé et de type fini.
Il existe alors un $i \in I$ et un diagramme
\begin{equation}
\label{equation-good-diagram}
\vcenter{
\xymatrix{
X \ar[r] \ar[d] & W \ar[d] \\
S \ar[r] & S_i
}
}
\end{equation}
tels que $W \to S_i$ soit de type fini et que
le morphisme induit $X \to S \times_{S_i} W$ soit une immersion
fermée.
\end{lemma}

\begin{proof}
D'après le Lemme \ref{lemma-finite-type-closed-in-finite-presentation},
on peut trouver une immersion fermée $X \to X'$
sur $S$, où $X'$ est un schéma de présentation finie sur $S$.
D'après le Lemme \ref{lemma-descend-finite-presentation},
on peut trouver un $i$ et un morphisme de présentation finie
$X'_i \to S_i$ dont l'image inverse est $X'$. Posons $W = X'_i$.
\end{proof}

\begin{lemma}
\label{lemma-limit-from-good-diagram}
Plaçons-nous dans la Situation \ref{situation-limit-noetherian}.
Soit $X \to S$ quasi-séparé et de type fini.
Étant donnés $i \in I$ et un diagramme
$$
\vcenter{
\xymatrix{
X \ar[r] \ar[d] & W \ar[d] \\
S \ar[r] & S_i
}
}
$$
comme dans (\ref{equation-good-diagram}), pour $i' \geq i$, soit
$X_{i'}$ l'image schématique de $X \to S_{i'} \times_{S_i} W$.
Alors $X = \lim_{i' \geq i} X_{i'}$.
\end{lemma}

\begin{proof}
Puisque $X$ est quasi-compact et quasi-séparé, la formation de
l'image schématique de $X \to S_{i'} \times_{S_i} W$
commute à la restriction aux sous-schémas ouverts
(Morphismes, Lemme \ref{morphisms-lemma-quasi-compact-scheme-theoretic-image}).
On peut donc supposer, et l'on suppose, $W$ affine et son image contenue dans un ouvert affine
$U_i$ de $S_i$. Soient $U \subset S$, $U_{i'} \subset S_{i'}$
les images réciproques de $U_i$. Alors $U$, $U_{i'}$,
$S_{i'} \times_{S_i} W = U_{i'} \times_{U_i} W$ et
$S \times_{S_i} W = U \times_{U_i} W$ sont tous affines.
Il s'ensuit que $X$ est affine, puisque $X \to S \times_{S_i} W$ est
une immersion fermée. Cela montre aussi que l'homomorphisme
$$
\mathcal{O}(U) \otimes_{\mathcal{O}(U_i)} \mathcal{O}(W) \to
\mathcal{O}(X)
$$
est surjectif. Soit $I$ son noyau. Alors $X_{i'}$
est le spectre de l'anneau
$$
\mathcal{O}(X_{i'}) =
\mathcal{O}(U_{i'}) \otimes_{\mathcal{O}(U_i)} \mathcal{O}(W)/I_{i'}
$$
où $I_{i'}$ est l'image réciproque de l'idéal $I$ (voir
Morphismes, Exemple \ref{morphisms-example-scheme-theoretic-image}).
Puisque $\mathcal{O}(U) = \colim \mathcal{O}(U_{i'})$,
on a $I = \colim I_{i'}$, et l'on conclut
que $\colim \mathcal{O}(X_{i'}) = \mathcal{O}(X)$.
\end{proof}

\begin{lemma}
\label{lemma-morphism-good-diagram}
Plaçons-nous dans la Situation \ref{situation-limit-noetherian}.
Soit $f : X \to Y$ un morphisme de schémas quasi-séparés
et de type fini sur $S$. Soient
$$
\vcenter{
\xymatrix{
X \ar[r] \ar[d] & W \ar[d] \\
S \ar[r] & S_{i_1}
}
}
\quad\text{et}\quad
\vcenter{
\xymatrix{
Y \ar[r] \ar[d] & V \ar[d] \\
S \ar[r] & S_{i_2}
}
}
$$
des diagrammes comme dans (\ref{equation-good-diagram}). Soient
$X = \lim_{i \geq i_1} X_i$ et
$Y = \lim_{i \geq i_2} Y_i$ les descriptions
correspondantes comme limites, données par le Lemme \ref{lemma-limit-from-good-diagram}.
Il existe alors un $i_0 \geq \max(i_1, i_2)$ et un morphisme
$$
(f_i)_{i \geq i_0} : (X_i)_{i \geq i_0} \to (Y_i)_{i \geq i_0}
$$
de systèmes projectifs sur $(S_i)_{i \geq i_0}$ tel que
$f = \lim_{i \geq i_0} f_i$.
Si $(g_i)_{i \geq i_0} : (X_i)_{i \geq i_0} \to (Y_i)_{i \geq i_0}$
est un second morphisme de systèmes projectifs sur $(S_i)_{i \geq i_0}$ tel que
$f = \lim_{i \geq i_0} g_i$,
alors $f_i = g_i$ pour tout $i \gg i_0$.
\end{lemma}

\begin{proof}
Puisque $V \to S_{i_2}$ est de présentation finie et
$X = \lim_{i \geq i_1} X_i$, on peut appliquer la Proposition
\ref{proposition-characterize-locally-finite-presentation}
pour trouver un $i_0 \geq \max(i_1, i_2)$ et un morphisme $h : X_{i_0} \to V$
sur $S_{i_2}$ tels que $X \to X_{i_0} \to V$ soit égal à $X \to Y \to V$.
Pour $i \geq i_0$, on obtient un diagramme commutatif de flèches pleines
$$
\xymatrix{
X \ar[d] \ar[r] &
X_i \ar[r] \ar@{..>}[d] \ar@/_2pc/[dd] |!{[d];[ld]}\hole &
X_{i_0} \ar[d]^h \\
Y \ar[r] \ar[d] & Y_i \ar[r] \ar[d] & V \ar[d] \\
S \ar[r] & S_i \ar[r] & S_{i_0}
}
$$
Puisque $X \to X_i$ a une image schématiquement dense
et que $Y_i$ est l'image schématique de
$Y \to S_i \times_{S_{i_2}} V$,
le morphisme $X_i \to S_i \times_{S_{i_2}} V$
induit par le diagramme
se factorise par $Y_i$ (Morphismes, Lemme \ref{morphisms-lemma-factor-factor}).
Cela démontre l'existence.

\medskip\noindent
Unicité. Soit $E_i \subset X_i$ l'égalisateur de $f_i$ et $g_i$
pour $i \geq i_0$. D'après
Schémas, Lemme \ref{schemes-lemma-where-are-they-equal},
$E_i$ est un sous-schéma localement fermé de $X_i$.
Puisque $X_i$ est un sous-schéma fermé de $S_i \times_{S_{i_0}} X_{i_0}$,
et de même pour $Y_i$, on a
$$
E_i = X_i \times_{(S_i \times_{S_{i_0}} X_{i_0})} (S_i \times_{S_{i_0}} E_{i_0})
$$
Pour achever la démonstration, il suffit donc de montrer que $X_i \to X_{i_0}$
se factorise par $E_{i_0}$ pour un certain $i \geq i_0$.
Pour cela, on utilise le fait que $X \to X_{i_0}$ se factorise par $E_{i_0}$,
puisque $f_{i_0}$ et $g_{i_0}$ sont tous deux compatibles avec $f$.
Puisque $X_i$ est noethérien, l'espace topologique
sous-jacent $|E_{i_0}|$ est une partie constructible de $|X_{i_0}|$
(Topologie, Lemme \ref{topology-lemma-constructible-Noetherian-space}).
Ainsi $X_i \to X_{i_0}$ se factorise ensemblistement par $E_{i_0}$
pour $i$ assez grand, d'après le Lemme \ref{lemma-limit-contained-in-constructible}.
Pour un tel $i$, l'image réciproque schématique
$(X_i \to X_{i_0})^{-1}(E_{i_0})$ est un sous-schéma fermé de $X_i$
par lequel $X$ se factorise ; elle est donc égale à $X_i$, puisque
$X \to X_i$ a une image schématiquement dense par construction.
Cela conclut la démonstration.
\end{proof}

\begin{remark}
\label{remark-finite-type-gives-well-defined-system}
Dans la Situation \ref{situation-limit-noetherian},
les Lemmes \ref{lemma-good-diagram}, \ref{lemma-limit-from-good-diagram} et
\ref{lemma-morphism-good-diagram}
montrent que la catégorie des schémas quasi-séparés et
de type fini sur $S$ est équivalente à certains types de
systèmes projectifs de schémas sur $(S_i)_{i \in I}$, à savoir
ceux obtenus en appliquant le Lemme \ref{lemma-limit-from-good-diagram}
à un diagramme de la forme (\ref{equation-good-diagram}).
Par exemple, étant donné $X \to S$ de type fini et quasi-séparé,
si l'on choisit deux diagrammes différents $X \to V_1 \to S_{i_1}$
et $X \to V_2 \to S_{i_2}$ comme dans (\ref{equation-good-diagram}),
le Lemme \ref{lemma-morphism-good-diagram} appliqué à $\text{id}_X$
(dans les deux sens)
montre que les descriptions correspondantes de
$X$ comme limite sont canoniquement isomorphes (quitte à restreindre
l'ensemble ordonné filtrant $I$). Et ainsi de suite.
\end{remark}

\begin{lemma}
\label{lemma-morphism-good-diagram-flat}
Reprenons les notations et les hypothèses du Lemme \ref{lemma-morphism-good-diagram}.
Si $f$ est plat et de présentation finie,
il existe un $i_3 \geq i_0$ tel que, pour $i \geq i_3$,
$f_i$ soit plat, $X_i = Y_i \times_{Y_{i_3}} X_{i_3}$ et
$X = Y \times_{Y_{i_3}} X_{i_3}$.
\end{lemma}

\begin{proof}
D'après le Lemme \ref{lemma-descend-finite-presentation},
on peut choisir un $i \geq i_2$ et un morphisme
$U \to Y_i$ de présentation finie tels que $X = Y \times_{Y_i} U$
(c'est ici qu'intervient la présentation finie de $f$).
En augmentant $i$, on peut supposer $U \to Y_i$ plat ; voir
Lemme \ref{lemma-descend-flat-finite-presentation}.
Comme expliqué dans la Remarque \ref{remark-finite-type-gives-well-defined-system},
on peut remplacer, et l'on remplace, le diagramme initial définissant le système
$(X_i)_{i \geq i_1}$ par le système correspondant à
$X \to U \to S_i$. Ainsi $X_{i'}$, pour $i' \geq i$, est défini comme
l'image schématique de $X \to S_{i'} \times_{S_i} U$.

\medskip\noindent
Puisque $U \to Y_i$ est plat (c'est ici qu'intervient la platitude de $f$),
que $X = Y \times_{Y_i} U$ et
que l'image schématique de $Y \to Y_i$ est $Y_i$,
l'image schématique de $X \to U$ est $U$
(Morphismes, Lemme
\ref{morphisms-lemma-flat-base-change-scheme-theoretic-image}).
Remarquons que $Y_{i'} \to S_{i'} \times_{S_i} Y_i$ est une immersion
fermée pour $i' \geq i$, par construction du système des $Y_j$.
Le même argument que ci-dessus montre alors que l'image schématique
de $X \to S_{i'} \times_{S_i} U$
est égale au sous-schéma fermé $Y_{i'} \times_{Y_i} U$.
Ainsi $X_{i'} = Y_{i'} \times_{Y_i} U$ pour tout $i' \geq i$,
et le lemme est satisfait avec $i_3 = i$.
\end{proof}

\begin{lemma}
\label{lemma-morphism-good-diagram-smooth}
Reprenons les notations et les hypothèses du Lemme \ref{lemma-morphism-good-diagram}.
Si $f$ est lisse, il existe un $i_3 \geq i_0$ tel que, pour
$i \geq i_3$, le morphisme $f_i$ soit lisse.
\end{lemma}

\begin{proof}
Combiner les Lemmes \ref{lemma-morphism-good-diagram-flat} et
\ref{lemma-descend-smooth}.
\end{proof}

\begin{lemma}
\label{lemma-morphism-good-diagram-proper}
Reprenons les notations et les hypothèses du Lemme \ref{lemma-morphism-good-diagram}.
Si $f$ est propre, il existe un $i_3 \geq i_0$ tel que, pour
$i \geq i_3$, le morphisme $f_i$ soit propre.
\end{lemma}

\begin{proof}
D'après la discussion de la
Remarque \ref{remark-finite-type-gives-well-defined-system},
le choix de $i_1$ et de $W$ s'insérant dans un diagramme comme dans
(\ref{equation-good-diagram}) n'influe pas sur la validité
du lemme. Choisissons donc $W$ comme suit.
Choisissons d'abord une immersion fermée $X \to X'$
avec $X' \to S$ propre et de présentation finie ; voir
Lemme \ref{lemma-proper-limit-of-proper-finite-presentation}.
Choisissons ensuite un $i_3 \geq i_2$ et un morphisme propre $W \to Y_{i_3}$
tels que $X' = Y \times_{Y_{i_3}} W$. C'est possible puisque
$Y = \lim_{i \geq i_2} Y_i$, d'après les
Lemmes \ref{lemma-descend-finite-presentation} et
\ref{lemma-eventually-proper}.
Avec ce choix de $W$, la construction montre immédiatement que,
pour $i \geq i_3$, le schéma $X_i$ est un sous-schéma fermé de
$Y_i \times_{Y_{i_3}} W \subset S_i \times_{S_{i_3}} W$,
donc est propre sur $Y_i$.
\end{proof}

\begin{lemma}
\label{lemma-good-diagram-fibre-product}
Dans la Situation \ref{situation-limit-noetherian}, supposons donné un
diagramme cartésien
$$
\xymatrix{
X^1 \ar[r]_p \ar[d]_q & X^3 \ar[d]^a \\
X^2 \ar[r]^b & X^4
}
$$
de schémas quasi-séparés et de type fini sur $S$.
Pour chaque $j = 1, 2, 3, 4$, choisissons $i_j \in I$ et un diagramme
$$
\xymatrix{
X^j \ar[r] \ar[d] & W^j \ar[d] \\
S \ar[r] & S_{i_j}
}
$$
comme dans (\ref{equation-good-diagram}). Soit
$X^j = \lim_{i \geq i_j} X^j_i$ la description correspondante comme limite,
donnée par le Lemme \ref{lemma-morphism-good-diagram}.
Soient $(a_i)_{i \geq i_5}$, $(b_i)_{i \geq i_6}$, $(p_i)_{i \geq i_7}$ et
$(q_i)_{i \geq i_8}$ les morphismes de systèmes correspondants, construits
dans le Lemme \ref{lemma-morphism-good-diagram}. Il existe alors un
$i_9 \geq \max(i_5, i_6, i_7, i_8)$ tel que, pour $i \geq i_9$,
$a_i \circ p_i = b_i \circ q_i$ et que
$$
(q_i, p_i) : X^1_i \longrightarrow X^2_i \times_{b_i, X^4_i, a_i} X^3_i
$$
soit une immersion fermée.
Si $a$ et $b$ sont plats et de présentation finie, il existe un
$i_{10} \geq \max(i_5, i_6, i_7, i_8, i_9)$ tel que, pour $i \geq i_{10}$,
le dernier morphisme affiché soit un isomorphisme.
\end{lemma}

\begin{proof}
D'après la discussion de la
Remarque \ref{remark-finite-type-gives-well-defined-system},
le choix de $W^1$ s'insérant dans un diagramme comme dans
(\ref{equation-good-diagram}) n'influe pas sur la validité
du lemme. On peut donc choisir $W^1 = W^2 \times_{W^4} W^3$.
La construction de $X^1_i$ montre alors immédiatement que
$a_i \circ p_i = b_i \circ q_i$ et que
$$
(q_i, p_i) : X^1_i \longrightarrow X^2_i \times_{b_i, X^4_i, a_i} X^3_i
$$
est une immersion fermée.

\medskip\noindent
Si $a$ et $b$ sont plats et de présentation finie, il en va de même pour
$p$ et $q$, qui sont des changements de base de $a$ et $b$. On peut donc appliquer le
Lemme \ref{lemma-morphism-good-diagram-flat}
à chacun des morphismes $a$, $b$, $p$, $q$ et $a \circ p = b \circ q$.
Il s'ensuit qu'il existe un $i_9 \in I$ tel que
$$
(q_i, p_i) : X^1_i \to X^2_i \times_{X^4_i} X^3_i
$$
soit le changement de base de $(q_{i_9}, p_{i_9})$ par le morphisme
$X^4_i \to X^4_{i_9}$ pour tout $i \geq i_9$.
On conclut que $(q_i, p_i)$ est un isomorphisme pour tout
$i$ assez grand, d'après le Lemme \ref{lemma-descend-isomorphism}.
\end{proof}






\begin{multicols}{2}[\section{Autres chapitres}]
\noindent
Préliminaires
\begin{enumerate}
\item \hyperref[introduction-section-phantom]{Introduction}
\item \hyperref[conventions-section-phantom]{Conventions}
\item \hyperref[sets-section-phantom]{Théorie des ensembles}
\item \hyperref[categories-section-phantom]{Catégories}
\item \hyperref[topology-section-phantom]{Topologie}
\item \hyperref[sheaves-section-phantom]{Faisceaux sur les espaces}
\item \hyperref[sites-section-phantom]{Sites et faisceaux}
\item \hyperref[stacks-section-phantom]{Champs}
\item \hyperref[fields-section-phantom]{Corps}
\item \hyperref[algebra-section-phantom]{Algèbre commutative}
\item \hyperref[brauer-section-phantom]{Groupes de Brauer}
\item \hyperref[homology-section-phantom]{Algèbre homologique}
\item \hyperref[derived-section-phantom]{Catégories dérivées}
\item \hyperref[simplicial-section-phantom]{Méthodes simpliciales}
\item \hyperref[more-algebra-section-phantom]{Compléments d'algèbre}
\item \hyperref[smoothing-section-phantom]{Lissage des morphismes d'anneaux}
\item \hyperref[modules-section-phantom]{Faisceaux de modules}
\item \hyperref[sites-modules-section-phantom]{Modules sur les sites}
\item \hyperref[injectives-section-phantom]{Objets injectifs}
\item \hyperref[cohomology-section-phantom]{Cohomologie des faisceaux}
\item \hyperref[sites-cohomology-section-phantom]{Cohomologie sur les sites}
\item \hyperref[dga-section-phantom]{Algèbre différentielle graduée}
\item \hyperref[dpa-section-phantom]{Algèbre à puissances divisées}
\item \hyperref[sdga-section-phantom]{Faisceaux différentiels gradués}
\item \hyperref[hypercovering-section-phantom]{Hyperrecouvrements}
\end{enumerate}
Schémas
\begin{enumerate}
\setcounter{enumi}{25}
\item \hyperref[schemes-section-phantom]{Schémas}
\item \hyperref[constructions-section-phantom]{Constructions de schémas}
\item \hyperref[properties-section-phantom]{Propriétés des schémas}
\item \hyperref[morphisms-section-phantom]{Morphismes de schémas}
\item \hyperref[coherent-section-phantom]{Cohomologie des schémas}
\item \hyperref[divisors-section-phantom]{Diviseurs}
\item \hyperref[limits-section-phantom]{Limites de schémas}
\item \hyperref[varieties-section-phantom]{Variétés}
\item \hyperref[topologies-section-phantom]{Topologies sur les schémas}
\item \hyperref[descent-section-phantom]{Descente}
\item \hyperref[perfect-section-phantom]{Catégories dérivées de schémas}
\item \hyperref[more-morphisms-section-phantom]{Compléments sur les morphismes}
\item \hyperref[flat-section-phantom]{Compléments sur la platitude}
\item \hyperref[groupoids-section-phantom]{Schémas en groupoïdes}
\item \hyperref[more-groupoids-section-phantom]{Compléments sur les schémas en groupoïdes}
\item \hyperref[etale-section-phantom]{Morphismes étales de schémas}
\end{enumerate}
Thèmes de théorie des schémas
\begin{enumerate}
\setcounter{enumi}{41}
\item \hyperref[chow-section-phantom]{Homologie de Chow}
\item \hyperref[intersection-section-phantom]{Théorie de l'intersection}
\item \hyperref[pic-section-phantom]{Schémas de Picard des courbes}
\item \hyperref[weil-section-phantom]{Théories cohomologiques de Weil}
\item \hyperref[adequate-section-phantom]{Modules adéquats}
\item \hyperref[dualizing-section-phantom]{Complexes dualisants}
\item \hyperref[duality-section-phantom]{Dualité pour les schémas}
\item \hyperref[discriminant-section-phantom]{Discriminants et différentes}
\item \hyperref[derham-section-phantom]{Cohomologie de de Rham}
\item \hyperref[local-cohomology-section-phantom]{Cohomologie locale}
\item \hyperref[algebraization-section-phantom]{Géométrie algébrique et formelle}
\item \hyperref[curves-section-phantom]{Courbes algébriques}
\item \hyperref[resolve-section-phantom]{Résolution des surfaces}
\item \hyperref[models-section-phantom]{Réduction semi-stable}
\item \hyperref[functors-section-phantom]{Foncteurs et morphismes}
\item \hyperref[equiv-section-phantom]{Catégories dérivées de variétés}
\item \hyperref[pione-section-phantom]{Groupes fondamentaux de schémas}
\item \hyperref[etale-cohomology-section-phantom]{Cohomologie étale}
\item \hyperref[crystalline-section-phantom]{Cohomologie cristalline}
\item \hyperref[proetale-section-phantom]{Cohomologie pro-étale}
\item \hyperref[relative-cycles-section-phantom]{Cycles relatifs}
\item \hyperref[more-etale-section-phantom]{Compléments de cohomologie étale}
\item \hyperref[trace-section-phantom]{La formule des traces}
\end{enumerate}
Espaces algébriques
\begin{enumerate}
\setcounter{enumi}{64}
\item \hyperref[spaces-section-phantom]{Espaces algébriques}
\item \hyperref[spaces-properties-section-phantom]{Propriétés des espaces algébriques}
\item \hyperref[spaces-morphisms-section-phantom]{Morphismes d'espaces algébriques}
\item \hyperref[decent-spaces-section-phantom]{Espaces algébriques décents}
\item \hyperref[spaces-cohomology-section-phantom]{Cohomologie des espaces algébriques}
\item \hyperref[spaces-limits-section-phantom]{Limites d'espaces algébriques}
\item \hyperref[spaces-divisors-section-phantom]{Diviseurs sur les espaces algébriques}
\item \hyperref[spaces-over-fields-section-phantom]{Espaces algébriques sur des corps}
\item \hyperref[spaces-topologies-section-phantom]{Topologies sur les espaces algébriques}
\item \hyperref[spaces-descent-section-phantom]{Descente et espaces algébriques}
\item \hyperref[spaces-perfect-section-phantom]{Catégories dérivées d'espaces}
\item \hyperref[spaces-more-morphisms-section-phantom]{Compléments sur les morphismes d'espaces}
\item \hyperref[spaces-flat-section-phantom]{Platitude sur les espaces algébriques}
\item \hyperref[spaces-groupoids-section-phantom]{Groupoïdes dans les espaces algébriques}
\item \hyperref[spaces-more-groupoids-section-phantom]{Compléments sur les groupoïdes dans les espaces}
\item \hyperref[bootstrap-section-phantom]{Amorçage}
\item \hyperref[spaces-pushouts-section-phantom]{Sommes amalgamées d'espaces algébriques}
\end{enumerate}
Thèmes de géométrie
\begin{enumerate}
\setcounter{enumi}{81}
\item \hyperref[spaces-chow-section-phantom]{Groupes de Chow des espaces}
\item \hyperref[groupoids-quotients-section-phantom]{Quotients de groupoïdes}
\item \hyperref[spaces-more-cohomology-section-phantom]{Compléments sur la cohomologie des espaces}
\item \hyperref[spaces-simplicial-section-phantom]{Espaces simpliciaux}
\item \hyperref[spaces-duality-section-phantom]{Dualité pour les espaces}
\item \hyperref[formal-spaces-section-phantom]{Espaces algébriques formels}
\item \hyperref[restricted-section-phantom]{Algébrisation des espaces formels}
\item \hyperref[spaces-resolve-section-phantom]{Résolution des surfaces revisitée}
\end{enumerate}
Théorie des déformations
\begin{enumerate}
\setcounter{enumi}{89}
\item \hyperref[formal-defos-section-phantom]{Théorie formelle des déformations}
\item \hyperref[defos-section-phantom]{Théorie des déformations}
\item \hyperref[cotangent-section-phantom]{Le complexe cotangent}
\item \hyperref[examples-defos-section-phantom]{Problèmes de déformation}
\end{enumerate}
Champs algébriques
\begin{enumerate}
\setcounter{enumi}{93}
\item \hyperref[algebraic-section-phantom]{Champs algébriques}
\item \hyperref[examples-stacks-section-phantom]{Exemples de champs}
\item \hyperref[stacks-sheaves-section-phantom]{Faisceaux sur les champs algébriques}
\item \hyperref[criteria-section-phantom]{Critères de représentabilité}
\item \hyperref[artin-section-phantom]{Axiomes d'Artin}
\item \hyperref[quot-section-phantom]{Espaces de Quot et de Hilbert}
\item \hyperref[stacks-properties-section-phantom]{Propriétés des champs algébriques}
\item \hyperref[stacks-morphisms-section-phantom]{Morphismes de champs algébriques}
\item \hyperref[stacks-limits-section-phantom]{Limites de champs algébriques}
\item \hyperref[stacks-cohomology-section-phantom]{Cohomologie des champs algébriques}
\item \hyperref[stacks-perfect-section-phantom]{Catégories dérivées de champs}
\item \hyperref[stacks-introduction-section-phantom]{Introduction aux champs algébriques}
\item \hyperref[stacks-more-morphisms-section-phantom]{Compléments sur les morphismes de champs}
\item \hyperref[stacks-geometry-section-phantom]{La géométrie des champs}
\end{enumerate}
Thèmes de théorie des espaces de modules
\begin{enumerate}
\setcounter{enumi}{107}
\item \hyperref[moduli-section-phantom]{Champs de modules}
\item \hyperref[moduli-curves-section-phantom]{Espaces de modules de courbes}
\end{enumerate}
Divers
\begin{enumerate}
\setcounter{enumi}{109}
\item \hyperref[examples-section-phantom]{Exemples}
\item \hyperref[exercises-section-phantom]{Exercices}
\item \hyperref[guide-section-phantom]{Guide bibliographique}
\item \hyperref[desirables-section-phantom]{Desiderata}
\item \hyperref[coding-section-phantom]{Style de programmation}
\item \hyperref[obsolete-section-phantom]{Obsolète}
\item \hyperref[fdl-section-phantom]{Licence de documentation libre GNU}
\item \hyperref[index-section-phantom]{Index généré automatiquement}
\end{enumerate}
\end{multicols}


\bibliography{my}
\bibliographystyle{amsalpha}

\end{document}

